% !TEX root = ../main.tex
\clearpage
\chapter{부록}
\label{ch:appendix}

% Demote headings one level under Chapter 8 (sections become 8.1.x).
\renewcommand{\dissertationSubheading}[2][]{%
  \phantomsection
  \subsection{#2}%
  \if\relax\detokenize{#1}\relax\else\label{#1}\fi
}
\renewcommand{\dissertationSubsubheading}[2][]{%
  \phantomsection
  \subsubsection{#2}%
  \if\relax\detokenize{#1}\relax\else\label{#1}\fi
}

\section{평가 처리 과정과 재현성}
\label{ch:appendix-eval}

이 부록은 \ref{ch:results}장의 바탕이 되는 평가 처리 과정(파이프라인)을 기록해요. 처리 과정의 단계, 추출 설계, 규모에 따른 변화 표, 계정 종류 나누기, 시간을 잰 모든 실행의 기록, 기본 매개변수(미리 정해 둔 설정값), 대리 지표(대신 재는 값), 그리고 용어 풀이를 담고 있어요. 처리 과정은 뽑아 낸 같은 로그(계약이 남기는 기록)에서 두 가지 분석을 계산해요. 하나는 분석 계획과 등록 문서를 따른 등록 분석(미리 등록한 그대로 한 분석)이에요(부록~\ref{sec:map-registered}). 다른 하나는 \ref{sec:amount-revision}절의 수정 분석(등록 뒤에 금액을 달러로 바꿔 다시 한 분석)이에요. \ref{ch:results}장이 보고하는 것은 이 수정 분석이에요. 경로는 공개 저장소 \url{https://github.com/abitocodes/phd_works}의 \texttt{1-Dissertation-Works/ERC20-Allowance-PageRank-Wallet-Reputation/experiments} 폴더를 기준으로 한 상대 경로로 적어요.

\dissertationSubheading[sec:pipeline-stages]{처리 과정의 단계}

평가는 일곱 단계로 돌아가요. 1단계부터 6단계까지가 등록 분석을 만들고, 7단계가 수정 분석을 만들어요. 2단계부터 4단계까지는 관측 기간(점수를 매길 자료를 모은 기간)인 2023년 10월 1일부터 2026년 6월 30일까지의 로그를 읽어요(UTC, 양 끝 날 포함). 5단계는 2026년 7월부터 9월까지를 더해요. 6단계와 7단계는 프로젝트 표에서만 이벤트를 읽고, 공개 로그는 하나도 훑지 않아요. 6단계는 3단계에서 만든 관측 기간의 이벤트를 읽고, 7단계는 2단계부터 5단계까지에서 만든 두 기간의 이벤트를 읽어요. 1단계는 로그를 하나도 읽지 않아요.

\begin{enumerate}
\item 계획과 설정: 분석 계획(\texttt{docs/analysis\_plan.md}), 설정 파일(\texttt{config/contract\_reputation.yaml}), 로그를 미리 표로 만들어 두는 질의(\texttt{sql/01\_materialize\_logs.sql}), 그리고 빅쿼리(BigQuery) 실행 프로그램(\texttt{scripts/run\_bq.py})을 \texttt{43d8045}에 함께 커밋했어요. 이 연구의 로그를 하나도 뽑기 전이었어요. 계획은 코호트(묶음), 점수, 대리 지표와 라벨(나중에 채점에 쓰는 정답 값), 통계, 그리고 결합 연산자(두 그래프를 함께 걷는 계산)에 대한 규칙~A를 정해 두어요. 7단계는 몇 군데에서 계획과 달라요. 토큰과 금액이 다르고, 그에 따라 코호트도 달라요. 그리고 강건성 점검(조건을 바꿔도 결과가 버티는지 보기)에서는 상위 토큰 부분 그래프 대신 금액 기울기 점검을 해요.
\item 스캔~A와 코호트: \texttt{scripts/run\_bq.py} 스크립트는 모든 질의를 먼저 미리 돌려 봐요. 그리고 어림값을 설정 파일의 예산과 견주어 본 뒤 질의를 돌리고, 요금이 매겨진 바이트 수를 \texttt{data/bq\_ledger.json}에 적어요. 스캔~A는 관측 기간의 로그를 미리 표로 만들어 두어요(부록~\ref{sec:bq-extraction}). \texttt{sql/02\_spender\_candidates.sql} 질의는 기간 안에 서로 다른 주인(허락을 해 주는 쪽) 적어도 세 명에게서 0이 아닌 허락(토큰을 얼마까지 대신 써도 된다는 허락)을 받은 스펜더(spender, 허락을 받아 대신 쓰는 쪽)를 나열해요. \texttt{scripts/check\_account\_code.py} 스크립트는 그 스펜더 하나하나의 계정 종류를 읽어요(부록~\ref{sec:account-type-classification}). 그리고 \texttt{scripts/build\_cohort.py} 스크립트는 그 가운데 스마트 계약(블록체인 위에서 저절로 실행되는 프로그램)을 계획의 코호트로 남기고, 그 주소들을 빅쿼리에 올려요. 7단계의 수정 코호트는 이 코호트의 일부예요.
\item 그래프 표, 대리 지표, 라벨: \texttt{sql/03\_ego\_events.sql} 질의는 코호트의 에고 네트워크(코호트 계약에 닿는 화살표만 모은 그래프)를 뽑아요. \texttt{sql/04\_anchor\_pairs.sql} 질의는 그것을 모아서 $T_{\text{obs}}$와 $t_1$에서 두 그래프의 쌍 표를 만들어요. \texttt{sql/05\_proxies\_same\_window.sql} 질의는 송금(transfer) 대리 지표와 시빌(가짜 주소를 많이 만들어 속이는 공격) 안정성 대리 지표를 계산하고, \texttt{sql/06\_holdout\_labels.sql} 질의는 W0의 라벨을 계산해요. \texttt{sql/07\_gmx\_accounts.sql} 질의와 \texttt{sql/09\_gmx\_contract\_closes.sql} 질의는 계약 계정이 GMX~V2에서 포지션을 닫은 기록을 골라요. 그러면 \texttt{scripts/decode\_gmx\_events.py} 스크립트가 \texttt{data/0-gmx-abi-for-decoding-logs/}에 둔 이벤트 ABI(이벤트를 읽는 규격)로 그 기록을 풀어 읽어요. \texttt{sql/08\_export\_ids.sql} 질의는 주소를 정수로 된 점(노드) 번호로 바꿔요. \texttt{scripts/download\_graph\_tables.sh} 스크립트는 표를 parquet 조각으로 내려받아요. 그리고 \texttt{sql/10\_describe.sql} 질의는 계획의 에고 표에 있는 이벤트 수와 허락 금액을 세요.
\item 같은 기간과 W0: \texttt{scripts/run\_same\_window.py} 스크립트는 모든 방법으로 $T_{\text{obs}}$에서 점수를 매겨요. 그리고 같은 기간 점검, 대비, 규칙~A의 보조 부분을 계산해요. \texttt{scripts/run\_holdout.py} 스크립트는 모든 방법으로 $t_1$에서 점수를 매겨요. 그리고 W0의 스펜더 코호트와 트레이더 코호트를 규칙~A의 주 부분과 함께 평가해요.
\item 등록과 스캔~B: W1의 등록 문서(\texttt{docs/registration\_w1.md})는 \texttt{0d3a057}에 커밋했어요. 이때 설정 파일의 상태는 \texttt{registered}로 정해 두었어요. \texttt{scripts/extract\_w1.py} 스크립트는 두 파일이 로컬 변경 없이 커밋되어 있고 상태가 \texttt{registered}일 때만 돌아가요. 그렇지 않으면 돌기를 거부해요. 이 스크립트는 두 파일의 SHA-256 해시와 커밋을 \texttt{registered-w1/extraction\_manifest.json}에 적어요. 그다음 스캔~B, 7월부터 9월까지의 에고 추출, 그리고 그 기간의 GMX 계정 질의를 돌려요. W1의 라벨과 포지션 닫기 기록은 \texttt{sql/06}과 \texttt{sql/09}에서 나와요. 그리고 \texttt{scripts/run\_registered.py} 스크립트가 규칙~B, 그 규칙의 민감도 분석(조건을 바꿔 결과가 얼마나 움직이는지 보기) 두 개, 그리고 두 화살표(엣지) 종류 어느 쪽으로도 만들지 않은 결과 값에서의 비교를 평가해요. 순서는 등록 문서에 적힌 순서를 따라요.
\item 라벨을 본 뒤: \texttt{scripts/run\_supplementary.py} 스크립트는 라벨을 본 뒤에 정한 스펜더 대비를 계산해요. \texttt{scripts/run\_robustness.py} 스크립트는 상위 토큰 부분 그래프(\texttt{sql/11\_top\_tokens.sql} 질의를 함께 써요)와 표본 정의 단계를 계산해요. \texttt{scripts/sybil\_model.py} 스크립트는 시빌 비용 모형을 계산하고, \texttt{scripts/run\_benchmark.py} 스크립트는 시간 잰 실행을 돌려요. \texttt{scripts/publish\_results.py} 스크립트는 요약 파일들을 \texttt{data/3-published-results-for-thesis/}에 복사해요. 이때 모든 숫자의 출처를 밝힌 요약 목록(digest)과 체크섬(파일이 바뀌지 않았는지 확인하는 값) 목록도 함께 넣어요.
\item 수정: 4단계부터 6단계까지의 결과를 안 뒤에 금액을 고쳤어요(\ref{sec:amount-revision}절, \texttt{docs/revision\_price\_weighting.md}). \texttt{scripts/usd\_inputs.py} 스크립트는 Chainlink 가격 피드(블록체인에 올라오는 가격 정보) 목록에 토큰 규칙을 적용해요. 그리고 목록에 오른 토큰들의 소수 자릿수를 한 블록에서 \texttt{eth\_call}로 읽고, 하루 가격 표를 만들어요. \texttt{scripts/chainlink\_check.py} 스크립트는 그 가격을 Chainlink 가격 피드와 견주어요. 두 스크립트 모두 \texttt{data/0-usd-token-list-and-prices/}에 결과를 적어요. \texttt{scripts/run\_usd\_pipeline.py} 스크립트는 다섯 토큰과 가격 표를 빅쿼리에 올려요(\texttt{sql/19\_usd\_load.sql}). 그리고 \texttt{sql/12}부터 \texttt{sql/18}까지를 차례대로 돌려요. 이 질의들은 수정 코호트와 그 에고 표, $t_1$까지의 이벤트로 구한 금액 기울기와 상한(이보다 큰 금액은 이 값으로 자름), 두 기준 시점의 쌍 표, 대리 지표, W0과 W1의 라벨, 점 번호, 그리고 부록~\ref{ch:appendix-data}에서 설명한 개수를 만들어요. 그다음 \texttt{CONTRACT\_REP\_VARIANT=usd}로 두면, 4단계부터 6단계까지의 스크립트가 내려받은 표에서 수정 분석을 계산해요. 순서는 등록 분석과 같아요(\texttt{scripts/run\_usd\_local.ps1}). \texttt{scripts/run\_registered.py} 스크립트는 규칙~B, 그리고 두 화살표 종류 어느 쪽으로도 만들지 않은 결과 값에서의 비교를 고친 점수에 적용해요. 이것은 등록한 시험이 아니라 재분석이에요. \texttt{scripts/run\_robustness.py} 스크립트는 상위 토큰 부분 그래프 대신 금액 기울기 점검을 돌려요. \texttt{scripts/publish\_results.py} 스크립트는 수정한 요약 파일들을 자기 요약 목록과 함께 \texttt{data/3-published-results-for-thesis/revised-usd/}에 공개해요. 그리고 \texttt{scripts/export\_latex.py} 스크립트는 그 요약 파일들로 \texttt{results/tables/}의 표와 \texttt{Figures/generated/}의 그림을 만들어요. 표~\ref{tab:registered-vs-revised}은 등록 분석의 요약 파일도 읽어요.
\end{enumerate}

이 단계들의 순서는 커밋 기록, 비용 장부(ledger), 그리고 W1의 추출 목록(manifest)에 남아 있어요. 계획은 2026년 10월 9일 13:41~UTC에 커밋했어요. 비용 장부는 작업이 끝날 때마다 그 작업을 적어요. 장부에는 스캔~A의 첫 작업이 13:43~UTC에 적혀 있어요. 2단계부터 4단계까지의 나머지 질의와 스크립트는 이미 돌고 난 뒤에, 등록 문서와 함께 \texttt{0d3a057}로 저장소에 들어왔어요. 이때 같은 기간과 W0의 요약 파일, 그리고 아직 돌지 않은 5단계의 스크립트도 함께 들어왔어요. 그 커밋은 2026년 10월 9일 16:49:32~UTC에 했어요. 커밋에는 이 시각이 한국 표준시로 10월 10일 01:49:32로 적혀 있어요. \texttt{scripts/extract\_w1.py} 스크립트는 그보다 1초 뒤, 첫 질의를 돌리기 전에 추출 목록을 적었어요. 그리고 추출 목록에는 강제 실행 옵션(override)을 쓰지 않았다고 적혀 있어요. 등록 분석의 요약 목록 \texttt{results\_digest.md}에는 요약 파일들을 공개한 코드 버전인 \texttt{8d8b568}이 적혀 있어요. 또 요약 목록을 쓸 때 그 버전과 달랐던 스크립트의 이름도 적혀 있어요. 커밋 \texttt{8d8b568}은 6단계의 스크립트를 더하거나 늘렸어요. 그리고 \texttt{scripts/run\_holdout.py}에서는 W0의 대비를 묶는 방식만 바꾸었어요. W0의 요약 파일에는 이 변경으로 바뀐 값이 하나도 없다고 적혀 있어요. 등록 분석의 벤치마크(속도 재기)는 \texttt{8d8b568}에 커밋된 \texttt{scripts/run\_benchmark.py}로 돌렸어요. 이것은 그 요약 목록에 적힌 버전이에요.

금액 수정은 등록 분석의 요약 파일을 공개한 뒤인 2026년 10월 10일에 커밋했어요. 수정 문서는 07:54~UTC에(\texttt{a897475}), 다섯 토큰의 선택은 08:12~UTC에(\texttt{09dfc32}), 고친 상한은 11:24~UTC에(\texttt{2705b28}) 커밋했어요. 비용 장부에서 수정 코호트를 만든 질의는 이 커밋들 가운데 두 번째 커밋 뒤에 나와요. 그리고 대리 지표와 라벨의 질의는 세 번째 커밋 뒤에 나와요. 점수는 그 뒤에 로컬 컴퓨터에서 계산했어요. 수정 분석의 요약 목록에는 코드 버전 \texttt{2705b28}이 적혀 있고, 그 버전과 달랐던 파일 다섯 개의 이름이 적혀 있어요. 이 가운데 \texttt{scripts/run\_holdout.py}는 송금을 받은 적이 있는 스펜더의 차수(화살표를 센 수) 대비 두 개를 위치가 아니라 이름으로 고른다는 점이 달랐어요. 그리고 \texttt{scripts/chainlink\_check.py}는 가격 비교를 다섯 토큰으로만 좁힌다는 점이 달랐어요. 다섯 파일은 모두 수정한 요약 파일과 함께 \texttt{77e71f5}에 커밋했어요. 토큰 입력 파일도 \texttt{2705b28} 뒤에 다시 썼어요. 다섯 토큰, 소수 자릿수, 가격은 모두 같았어요. 빅쿼리에 올린 소수 자릿수는 블록 $513{,}457{,}143$에서 읽은 것이었고, 블록 $513{,}498{,}909$에서 한 번 더 읽었어요. 공개한 \texttt{usd\_inputs.json}에는 이 두 번째 읽기가 적혀 있어요. 또 토큰 목록과 후보 목록의 근거 칸을 새로 고쳤어요. 고친 것은 목록에 오른 것마다의 확인 결과, 그리고 USDT0과 WBTC가 자기 가격 피드와 맺는 관계예요. 이 관계는 이제 정식 브리지(체인 사이에 토큰을 옮기는 공식 통로)로 적혀 있어요. 이 파일들도 \texttt{77e71f5}에 커밋했어요. 수정 분석의 벤치마크는 \texttt{2705b28} 시점의 \texttt{scripts/run\_benchmark.py}로 돌렸어요. 이 버전은 \texttt{8d8b568}의 버전과 딱 한 가지만 달라요. NumPy와 SciPy 버전도 함께 적는다는 점이에요. 두 분석의 시간 잰 실행은 다른 일을 하지 않는 기계 한 대에서 했어요. 이 기계에는 Intel Core Ultra~9 185H 프로세서와 64~GB 메모리가 있고, Windows~11(빌드 26200)과 Python 3.11.9를 썼어요. \texttt{benchmark/benchmark.json}에는 플랫폼(\texttt{Windows-10-10.0.26200-SP0}, Python이 이 빌드에 대해 알려 주는 문자열), 프로세서(Intel64 Family~6 Model~170 Stepping~4), Python 버전이 적혀 있어요. 수정 분석에서는 NumPy 2.3.2와 SciPy 1.16.1도 적혀 있어요. 등록 분석에서는 프로세서 이름, 메모리(63.4~GB로 기록), 그리고 같은 NumPy와 SciPy 버전을 실행이 끝난 뒤에 같은 기계에서 파일에 더했어요.

공개 저장소 \url{https://github.com/abitocodes/phd_works}의 \texttt{1-Dissertation-Works/ERC20-Allowance-PageRank-Wallet-Reputation/experiments}에는 평가 소스 코드가 들어 있어요. 질의(\texttt{sql/}), 스크립트(\texttt{scripts/}), 테스트(\texttt{tests/}), 버전을 매긴 설정 파일(\texttt{config/contract\_reputation.yaml}), 계획과 등록 문서와 수정 문서(\texttt{docs/}), 수정에 쓴 입력 파일, 뽑아 낸 표, 요약 파일, 그리고 비용 장부예요. 수정에 쓴 입력 파일은 \texttt{data/0-usd-token-list-and-prices/}에 있어요. 소수 자릿수와 가격 피드, 그리고 목록에 오른 토큰마다의 근거가 붙은 토큰 목록(\texttt{listed\_tokens.csv}), DefiLlama의 원래 데이터 점이 붙은 하루 가격(\texttt{daily\_prices\_usd.csv}), 상수(\texttt{constants.json}), 그리고 Chainlink와 견준 결과예요. 뽑아 낸 표는 모두 커밋할 수 있도록 90\,MB보다 작은 parquet 조각으로 나누어 두었어요. 그리고 표마다 옆에 JSON 파일이 있어요. 이 파일에는 그 표가 나온 빅쿼리 표, 행 수, 내려받은 시각이 적혀 있어요(계정 종류 표에는 읽은 접속 주소(블록체인 노드에 질문을 보내는 주소), 블록, 시각이 적혀 있어요). 표의 바탕이 되는 요약 파일은 \texttt{data/3-published-results-for-thesis/revised-usd/}에 있고, 등록 분석의 요약 파일은 \texttt{data/3-published-results-for-thesis/}에 있어요. 폴더마다 설정 파일, 계획, 등록 문서의 복사본이 있어요. 또 모든 숫자에 대해 그 숫자가 든 요약 파일과 키를 알려 주는 요약 목록 \texttt{results\_digest.md}, 그리고 \texttt{SHA256SUMS.txt}가 있어요. 수정 분석 폴더에는 수정 문서, 토큰 목록, 상수, 가격 비교가 더 들어 있어요. 그래서 처리 과정을 다시 돌리지 않고도 표의 어떤 숫자든 그 요약 파일과 맞춰 볼 수 있어요.

점수 계산 코드와 통계 코드는 오프라인(인터넷에 연결하지 않은 상태)으로 점검해요. \texttt{tests/test\_scoring.py}와 \texttt{tests/test\_stats.py}는 작은 인공 그래프 위에서 풀이 프로그램(solver), 화살표 무게(가중치), 짝지은 부트스트랩(다시 뽑기)을 같은 정의로 만든 기준 구현과 견주어요. \texttt{tests/test\_usd.py}는 달러(USD)로 바꾸는 계산, 기울기 $b=\ln 3/m$, 수정한 SQL의 값 변환(금액이 아무리 커도 1을 넘지 않게 바꾸는 식)이 점수 계산 코드와 맞는지, 그리고 가격 규칙을 점검해요. 이 점검은 모두 클라우드에 접속하지 않고 해요. 실제로 추출하려면 구글 클라우드(Google Cloud) 자격 증명과 요금을 낼 프로젝트가 있어야 해요. 그리고 추출의 스캔은 공개된 아비트럼 원(Arbitrum One) 로그 표만 읽어요. 수정 분석은 토큰 소수 자릿수와 Chainlink 피드도 공개 RPC 접속 주소에서 읽었어요. 그리고 하루 가격은 DefiLlama API에서 읽었어요. 이 입력 파일들은 커밋되어 있어요. 그래서 수정 분석의 빅쿼리 단계는 두 출처 어느 쪽도 없이 다시 돌릴 수 있어요.

\dissertationSubheading[sec:bq-extraction]{빅쿼리 추출 설계}

공개 표 \texttt{bigquery-public-data.goog\_blockchain\_arbitrum\_one\_us.logs}는 기간마다 한 번씩만 읽어요. 스캔~A는 2023년 10월 1일부터 2026년 6월 30일까지에 대해 \texttt{sql/01\_materialize\_logs.sql} 질의를 돌려요. 그리고 네 종류의 로그를 달마다 나눈 프로젝트 표 하나(\texttt{contract\_rep.logs\_obs})에 남겨요. 남기는 로그는 이래요. 토픽(topic, 로그에 붙는 꼬리표)이 정확히 세 개인 ERC-20(토큰 공통 규칙) \texttt{Approval} 로그와 \texttt{Transfer} 로그가 있어요. 이 로그들은 표준 topic-0 해시로 알아봐요\parencite{vogelsteller2015erc20}. 또 GMX~V2 서비스(프로토콜)의 이벤트 발신기(이벤트를 내보내는 계약)가 남긴 \texttt{PositionDecrease} 로그가 있어요. 그리고 Aave~V3 풀의 \texttt{Borrow}, \texttt{Repay}, \texttt{LiquidationCall} 로그가 있어요. 이 논문의 어떤 분석도 Aave 로그는 쓰지 않아요. 스캔~B는 같은 질의를 2026년 7월부터 9월까지에 대해 돌려서 \texttt{contract\_rep.logs\_w1}에 넣어요. 그 뒤의 질의는 모두 이 프로젝트 표들을 읽어요. 그래서 체인 전체를 훑는 비용은 기간마다 한 번만 들어요. 수정 분석은 공개 로그를 다시 훑지 않았어요. 부록~\ref{ch:appendix-sql}에는 두 분석의 그래프가 바탕으로 삼는 질의들이 있어요.

에고 추출(\texttt{sql/03\_ego\_events.sql})은 계획의 코호트에 든 계약이 주인이나 스펜더인 허락을 남겨요. 그리고 그런 계약이 보낸 쪽(송신자)이나 받는 쪽(수신자)인 송금을 남겨요. 그래서 각 그래프는 코호트의 에고 네트워크예요. 그리고 뽑아 낸 자료의 크기는 체인이 아니라 코호트를 따라가요. 이 에고 표들은 모든 토큰에 대해 금액을 원래 단위(소수점 조정을 하지 않은 가장 작은 단위)로 담아요. 곧 \texttt{uint256} 값을 배정밀도 실수(double)로 읽은 값이에요. 금액이 0인 송금은 버려요. 반면 금액이 0인 허락은 남겨요. 허락을 지우는 일이기 때문이에요. 관측 기간에 대해 이 표들에는 허락 로그 $389{,}866{,}882$개와 송금 로그 $1{,}538{,}404{,}067$개가 들어 있어요. 수정 분석은 이 표들에서 자기 에고 표를 잘라 내요(\texttt{sql/12\_usd\_ego.sql}). 이벤트를 다섯 토큰의 표와 이어 붙이고, 값마다 10을 그 토큰의 소수 자릿수만큼 거듭제곱한 수로 나눠요. 그다음 송금에는 그 블록이 있는 UTC 날짜의 토큰 가격을 곱해요. 허락은 나중에, 기준 시점 날의 가격으로 금액을 매겨요. 수정 분석은 한쪽 끝에 수정 코호트의 계약이 있는 다섯 토큰의 이벤트를 남겨요. 관측 기간에는 허락 $224{,}954{,}202$개와 송금 $1{,}149{,}168{,}538$개예요. 표~\ref{tab:usd-tokens}에는 이때 빠진 행 수가 있어요.

쌍 표는 한 기준 시점에서 두 그래프의 화살표 목록이에요. 쌍 표는 $T_{\text{obs}}$와 $t_1$에서, 그 기준 시점까지의 이벤트로 만들어요. 로지스틱 감쇠(오래된 일일수록 가볍게 세는 무게)도 그 기준 시점에 맞춰요. 그리고 쌍마다 모든 방법의 무게가 붙어 있어요(부록~\ref{sec:sql-pairs}). 수정 분석에서는 \texttt{sql/14\_usd\_anchor\_pairs.sql} 질의가 쌍 표를 달러로 만들어요. $T_{\text{obs}}$에서는 허락 쌍 $11{,}219{,}116$개와 송금 쌍 $25{,}018{,}892$개이고, $t_1$에서는 $10{,}650{,}258$개와 $23{,}680{,}692$개예요. 엔도스랭크(EndorseRank)와 AWP(적응형 가중 페이지랭크)는 무게가 0보다 큰 쌍만 남겨요. 그래서 금액이 1센트의 아주 작은 조각이라서 배정밀도 계산에서 무게가 0이 되는 쌍은 두 방법의 그래프에서 화살표가 아니에요(그래프 크기는 \ref{sec:dataset}절에 있어요). W1의 동결일(점수를 고정하는 날)은 $T_{\text{obs}}$와 같아요. 그래서 W1은 $T_{\text{obs}}$의 쌍 표를 다시 써요. \texttt{sql/17\_usd\_export\_ids.sql} 질의는 어느 기준 시점이든 수정 쌍 표에 나오거나 수정 코호트에 든 주소에 번호를 매겨요. 그 개수는 점 표의 JSON 기록에 있어요. 등록 분석에서는 \texttt{sql/04\_anchor\_pairs.sql} 질의와 \texttt{sql/08\_export\_ids.sql} 질의가 같은 일을 원래 단위로 했어요. $T_{\text{obs}}$에서 허락 쌍은 $22{,}321{,}470$개, 송금 쌍은 $59{,}737{,}617$개였어요. 번호로 바꾼 쌍 표, 점 표, 대리 지표, 라벨만 내려받아요. 이벤트 표는 빅쿼리에 그대로 남아요.

모든 작업은 \texttt{data/bq\_ledger.json}에 질의, 매개변수, 어림한 바이트 수와 요금이 매겨진 바이트 수, 그리고 비용과 함께 적혀 있어요. 장부에는 작업 32개가 있고, 이 작업들에 요금이 매겨진 양은 $6{,}122{,}812{,}473{,}344$바이트($5.57$\,TiB)예요. 장부에 적힌 쓴 만큼 내는 요금, 곧 TiB당 USD~6.25로 치면 USD~34.80이에요. 설정 파일에 정해 둔 예산은 USD~50이에요. 등록 분석의 작업 17개에는 $5{,}226{,}856{,}710{,}144$바이트($4.75$\,TiB), 곧 USD~29.71이 매겨졌어요. 스캔~A 하나에만 $3.29$\,TiB(USD~20.55)가 매겨졌고, 스캔~B에는 $0.26$\,TiB(USD~1.64)가 매겨졌어요. 등록 문서를 커밋했을 때 장부의 합계는 USD~26.20이었고, 금액 수정을 커밋했을 때는 USD~29.71이었어요. 수정 분석의 작업 15개에는 USD~5.09가 매겨졌어요. 그 가운데 세 개는 \texttt{scripts/run\_bq.py} 밖에서 돌렸고, 장부에 손으로 적어 넣었어요. 상수 표를 읽은 작업 두 개, 그리고 고친 상한의 바탕이 된 허락 금액 분포와 송금 금액 분포를 구한 질의예요.

\dissertationSubheading[sec:scaling-appendix]{규모를 키워 본 표}

\ref{sub:runtime-scaling}절의 표~\ref{tab:benchmark-scaling}는 엔도스랭크와 AWP의 실행 시간을 보여 줘요. 이 시간은 매칭 코호트(모든 점수가 똑같이 점수를 매기는 계약 묶음)에서 SHA256으로 늘 똑같이 뽑은 부분 표본(일부만 뽑은 묶음) 위에서 쟀어요. 시간은 \ref{sec:computational-benchmark}절의 시간 재는 방법대로 쟀어요. 이 절에서는 그 표의 단계들을 어떻게 뽑았는지 설명해요. 먼저 시드 문자열 \texttt{contract-benchmark-v1}에 주소를 이어 붙여요. 그리고 그 SHA256 해시 값의 순서로 계약들을 줄 세워요. 그다음 앞에서부터 $n$개를 한 단계로 삼아요. 단계는 $n=1{,}000$, $2{,}000$, $5{,}000$, $10{,}000$일 때 만들어요. 마지막 행은 전체 코호트($n=15{,}107$)예요. 이 행의 측정값은 표~\ref{tab:benchmark-runtime}의 측정값을 다시 쓴 것이에요. 한 단계는 뽑아 둔 화살표 가운데 한쪽 끝이 그 단계의 계약인 화살표를 모두 남겨요. 그래서 두 그래프의 $|V|$와 $|E|$는 $n$이 커질수록 함께 커져요. $n=1{,}000$에서는 허락 화살표가 $424{,}888$개, 송금 화살표가 $1{,}692{,}513$개예요. 전체 그래프에서는 이 수가 $11{,}210{,}114$개와 $25{,}001{,}216$개가 돼요. 표에는 단계마다 두 수가 모두 나와요. 이 단계들은 표~\ref{tab:robustness-sample-size}의 평가 집합이에요. 그 표는 이 표의 실행 시간을 다시 보여 줘요.

허락 화살표 수에 대한 송금 화살표 수의 비율은 단계마다 같지 않아요. 이 비율은 가장 작은 단계에서 가장 높아요. $n=2{,}000$부터 $n=10{,}000$까지는 거의 같고, 전체 코호트에서 가장 낮아요. 표에는 단계마다 이 비율이 나와요. 그리고 \ref{sub:runtime-scaling}절에서 이 비율을 다뤄요. 그래서 이 표를 보면, 두 실행 시간의 비율이 단계마다 화살표 수의 비율을 따라가는지 알 수 있어요. 어떤 단계도 코호트 그래프보다 크지 않아요. 그래서 이 표는 그보다 큰 그래프에 대해서는 아무것도 말해 주지 않아요.

\dissertationSubheading[sec:account-type-classification]{계정 종류 나누기}

점수를 매기고 평가하는 것은 계약뿐이에요. 그래서 코호트에 들어갈 수 있는 주소는 모두 먼저 그 코드로 종류를 나누었어요. \texttt{scripts/check\_account\_code.py}는 주소마다 한 블록에서 \texttt{eth\_getCode}를 불러요. 이 블록은 한 번 실행할 때 처음에 한 번만 가져와요. 그리고 이 블록을 읽은 시각과 함께 결과 파일에 적어요. 코드가 비어 있지 않은 주소는 계약이에요. 단, 그 코드가 EIP-7702 위임 표시이면 계약이 아니에요. 위임 표시는 \texttt{0xef0100} 바이트 뒤에 20바이트 주소가 붙은 것이에요\parencite{buterin2024eip7702}. 이런 주소는 코드가 없는 주소처럼 외부 소유 계정(EOA, 개인 키로 움직이는 계정)이에요. 표~\ref{tab:account-types}에는 네 번의 실행이 나와요. 네 실행은 모두 2026년 10월 9일에 등록 분석의 주소 집합에 대해 했어요. 등록 분석에서는 모든 토큰의 허락을 셌어요. 계획의 코호트의 후보는 공개 접속 주소 \url{https://arb1.arbitrum.io/rpc}에서 읽었어요. 그런데 $20{,}000$개 주소를 읽은 뒤에 이 접속 주소가 요청 속도를 제한했어요. 그래서 나머지 $9{,}022$개는 같은 블록에서 \url{https://arbitrum-one.public.blastapi.io}로 읽었어요. 뒤의 세 실행도 이 접속 주소를 썼어요. 수정 분석에서는 새로 읽지 않았어요. 수정 코호트는 모두 여기서 종류를 나눈 집합의 부분 집합이기 때문이에요.

\begin{table}[htbp]
\centering
\caption[\texttt{eth\_getCode}로 읽은 계정 종류.]{아비트럼 원에서 \texttt{eth\_getCode}로 읽은 계정 종류예요. 실행마다 한 블록에서 읽었어요. 계약은 코드가 비어 있지 않고, 그 코드가 위임 표시도 아닌 주소예요. 코드 없음과 EIP-7702(위임 표시)는 둘 다 외부 소유 계정이에요. 주소 집합은 등록 분석의 것이에요(어떤 토큰의 허락이든 셈해요). 수정 코호트는 이 집합들의 부분 집합이에요. W0 실행은 코호트 후보가 아닌 스펜더를 읽었어요. W1 실행은 앞의 어느 실행에서도 종류를 나누지 않은 주소를 읽었어요. 출처는 \texttt{data/1-raw-blockchain-extracts/cohort/}와 \texttt{gmx/}에서 parquet 결과 파일 옆에 있는 JSON 파일이에요.}
\label{tab:account-types}
\fitwidth{%
\begin{tabular}{lrrrrrr}
\toprule
주소 집합 & 블록 & 읽은 시각(UTC) & 주소 수 & 계약 & 코드 없음 & EIP-7702 \\
\midrule
코호트 후보: 관측 기간에 주인 $\geq 3$명에게서 0이 아닌 허락을 받음 & 513{,}216{,}506 & 2026년 10월 9일, 13:55 & 29{,}022 & 27{,}844 & 1{,}121 & 57 \\
관측 기간에 닫기를 $\geq 3$번 한 GMX~V2 계정 & 513{,}219{,}999 & 2026년 10월 9일, 14:10 & 40{,}307 & 7{,}633 & 28{,}075 & 4{,}599 \\
$t_1$에 0보다 큰 허락을 가진 W0 스펜더(코호트 후보가 아닌 것) & 513{,}233{,}527 & 2026년 10월 9일, 15:08 & 8{,}224 & 7{,}792 & 413 & 19 \\
W1 스펜더와 W1에서 닫기를 $\geq 3$번 한 GMX~V2 계정(앞에서 읽지 않은 것) & 513{,}257{,}476 & 2026년 10월 9일, 16:52 & 1{,}153 & 350 & 568 & 235 \\
\bottomrule
\end{tabular}
}
\end{table}

계획의 코호트는 후보 가운데 계약인 것들로 이루어져요. 수정 매칭 코호트는 그 가운데 다섯 토큰에서 주인 세 명 이상이 허락한 부분이에요(\ref{sec:amount-revision}절). 그래서 그 계약은 모두 첫 번째 실행에서 종류를 나누었어요. W0 스펜더 코호트는 첫 번째와 세 번째 실행을 써요. $t_1$에 다섯 토큰 가운데 하나의 0보다 큰 마지막 허락(가장 최근 허락)을 가진 스펜더 가운데, $17{,}050$개는 계약이에요. $104$개는 코드가 없고, $13$개는 위임 표시를 가지고 있어요. 이런 스펜더는 등록 분석의 에고 자료에서도 그 허락을 가지고 있었어요. 마지막 허락은 그 토큰, 주인, 스펜더의 이벤트에만 달려 있기 때문이에요. 그래서 이런 스펜더는 이 두 실행 가운데 하나에서 읽었어요. W1 스펜더 코호트와 두 트레이더 코호트는 네 실행 모두에서 계정 종류를 가져와요. 모든 읽기는 2026년 10월 9일의 체인을 기준으로 해요. 그래서 주소는 그날 가지고 있던 코드로 종류가 나뉘어요.

\dissertationSubheading[sec:benchmark-run-logs]{벤치마크 실행 기록}

표~\ref{tab:benchmark-runs}에는 수정 분석의 그래프에서 한 벤치마크의 기록이 있어요. 시간을 잰 풀이 프로그램 호출마다 실제 걸린 시간(초)을 실행한 순서대로 적었어요. 그리고 그 평균과 표준편차(값이 퍼진 정도)도 함께 적었어요. 들어 있는 것은 다음과 같아요. $d=0.85$일 때 전체 그래프에서의 네 점수, $d=0.75$와 $d=0.95$일 때의 엔도스랭크와 AWP, 그리고 전체 코호트보다 작은 모든 규모 단계에서의 두 점수예요. 설정마다, 시간 잰 실행 다섯 번 전에 시간을 재지 않는 준비 실행을 한 번 했어요. 가장 많이 쓴 메모리는 준비 실행에서만 쟀어요. 그래프와 설정마다 시간은 한 번만 쟀어요. 그리고 같은 그래프와 설정이 나오는 곳마다 그 측정값을 다시 써요. 표~\ref{tab:robustness-damping}의 $d=0.85$ 행과 표~\ref{tab:benchmark-scaling}의 전체 코호트 행은 표~\ref{tab:benchmark-runs}의 첫 번째 덩어리를 그대로 다시 보여 줘요. 모든 설정에서, 실행 다섯 번의 표준편차는 평균의 0.7\%에서 21.5\% 사이예요. 이 비율이 가장 큰 것은 $n=1{,}000$에서의 엔도스랭크예요. 이때 호출 한 번은 반 초쯤 걸려요. 등록 분석의 벤치마크는 계획의 코호트의 그래프에서 했어요. 이 벤치마크는 등록 분석의 요약과 함께 공개되어 있어요(부록~\ref{sec:map-registered}).

% Generated by experiments/scripts/export_latex.py -- do not edit by hand.
% Source: experiments/data/2-processed-tables-and-evaluations/revised-usd/benchmark/benchmark.json
\begin{table}[htbp]
\centering
\caption[벤치마크에서 시간을 잰 풀이 프로그램 호출 전부.]{벤치마크에서 시간을 잰 풀이 프로그램(solver) 호출마다 실제 걸린 시간을 초 단위로 보여 줘요(매칭 코호트, 계약 $n=15{,}107$개). 실행한 순서대로 적었고, 평균과 표준편차도 함께 적었어요. 설정마다 시간을 재지 않는 준비 실행을 한 번 한 뒤에 시간 잰 실행을 했어요. 그래프와 설정마다 시간은 한 번씩만 쟀어요. 표~\ref{tab:robustness-damping}의 $d=0.85$ 행과 표~\ref{tab:benchmark-scaling}의 코호트 전체 행은 첫 번째 묶음을 다시 쓴 것이에요. 단계는 표~\ref{tab:benchmark-scaling}의 부분 표본이에요.}
\label{tab:benchmark-runs}
\fitwidth{%
\begin{tabular}{lrrrrrrr}
\toprule
방법 & 실행 1 & 실행 2 & 실행 3 & 실행 4 & 실행 5 & 평균 & 표준편차 \\
\midrule
\multicolumn{8}{l}{\emph{전체 그래프, $d=0.85$}} \\
EndorseRank & 9.570 & 11.063 & 10.497 & 9.811 & 10.530 & 10.294 & 0.601 \\
AWP & 21.139 & 20.446 & 24.122 & 24.492 & 21.637 & 22.367 & 1.825 \\
C-PR ($\lambda=0.5$) & 33.813 & 39.576 & 35.439 & 35.793 & 36.187 & 36.162 & 2.112 \\
S-PR & 38.769 & 36.970 & 29.310 & 35.214 & 31.848 & 34.422 & 3.833 \\
\midrule
\multicolumn{8}{l}{\emph{전체 그래프, $d=0.75$}} \\
EndorseRank & 6.045 & 5.963 & 6.003 & 5.984 & 5.927 & 5.984 & 0.044 \\
AWP & 13.208 & 16.968 & 16.667 & 15.870 & 13.469 & 15.236 & 1.781 \\
\midrule
\multicolumn{8}{l}{\emph{전체 그래프, $d=0.95$}} \\
EndorseRank & 31.701 & 32.171 & 30.886 & 30.574 & 30.836 & 31.233 & 0.673 \\
AWP & 58.498 & 56.896 & 59.713 & 64.490 & 65.212 & 60.962 & 3.697 \\
\midrule
\multicolumn{8}{l}{\emph{단계 $n=1{,}000$}} \\
EndorseRank & 0.425 & 0.489 & 0.585 & 0.683 & 0.719 & 0.580 & 0.125 \\
AWP & 1.892 & 1.408 & 1.870 & 1.758 & 1.458 & 1.677 & 0.229 \\
\midrule
\multicolumn{8}{l}{\emph{단계 $n=2{,}000$}} \\
EndorseRank & 0.944 & 0.956 & 0.925 & 0.932 & 0.937 & 0.939 & 0.012 \\
AWP & 3.748 & 4.876 & 4.133 & 4.811 & 3.792 & 4.272 & 0.543 \\
\midrule
\multicolumn{8}{l}{\emph{단계 $n=5{,}000$}} \\
EndorseRank & 2.293 & 2.420 & 2.489 & 2.555 & 2.723 & 2.496 & 0.159 \\
AWP & 8.713 & 9.861 & 9.895 & 8.051 & 8.453 & 8.995 & 0.840 \\
\midrule
\multicolumn{8}{l}{\emph{단계 $n=10{,}000$}} \\
EndorseRank & 6.253 & 6.375 & 6.065 & 6.285 & 6.119 & 6.220 & 0.126 \\
AWP & 17.640 & 17.362 & 17.160 & 16.124 & 15.816 & 16.821 & 0.802 \\
\bottomrule
\end{tabular}
}
\end{table}


\dissertationSubheading[sec:parameter-summary]{매개변수 요약}

표~\ref{tab:parameter-summary}에는 수정 분석의 기본 매개변수를 모아 두었어요. 대부분은 설정의 키예요(부록~\ref{sec:config-keys}). 이것들은 로그를 하나라도 뽑기 전에 계획과 함께 정했어요. W1의 등록 문서는 규칙~B의 허용 폭(이 폭보다 더 나쁘지 않으면 된다고 정한 값)을 보탰어요. 또 트레이더 비교에 쓰는 다시 뽑은 표본 $2{,}000$개도 보탰어요. 코드 점검의 블록은 점검을 돌릴 때 정했어요. 수정은 등록 분석의 결과를 안 뒤에 했어요. 수정에서는 토큰, 가격, 금액 기울기들, 그리고 C-PR(결합 페이지랭크, 두 그래프를 함께 걷는 점수) 허락 층의 상한을 정했어요. 또 토큰 점검을 금액 기울기 점검으로 바꾸었어요(\ref{sec:amount-revision}절). 수정이 계획의 설정을 바꾼 곳에는 계획의 설정을 표에 괄호로 적었어요. 규칙~A와~B는 그 규칙에 정해 둔 조건과 허용 폭을 그대로 지켜요. 하지만 고친 점수에는 등록 분석의 결과를 안 뒤에 이 규칙들을 적용해요. 그래서 이것은 등록한 검정이 아니라 재분석이에요.

\begin{table}[htbp]
\centering
\caption{평가에 쓴 기본 매개변수(수정 분석; 수정이 바꾼 곳에는 계획의 설정을 괄호에 적음).}
\label{tab:parameter-summary}
\fitwidth{%
\begin{tabular}{lll}
\toprule
매개변수 & 값 & 하는 일 \\
\midrule
페이지랭크 감쇠 계수 $d$ & $0.85$; 바꿔 본 값 $0.75$, $0.95$ & 순간 이동 섞기 \\
수렴 허용 오차 $\varepsilon$ & $L_1$ 노름으로 $10^{-8}$ & 거듭제곱 반복 멈춤 \\
최대 반복 수 $I_{\max}$ & $300$ & 거듭제곱 반복 상한 \\
로지스틱 $k$ & $0.01$ & 로지스틱 감쇠 기울기(AWP, EndorseRank, C-PR의 송금 층) \\
로지스틱 $t_0$ & $180$일 & 로지스틱 감쇠 가운데 점(AWP, EndorseRank, C-PR의 송금 층) \\
값 변환 $b$ & 달러당 $\ln 3/m$, $m_T=200$, $m_A=62$ (원래 단위에서는 $1$) & AWP($m_T$)와 EndorseRank($m_A$)의 금액 무게 $V(z)$ \\
허락 상한 $U$ & $20{,}000$달러(없음) & C-PR 허락 층; 받은 허락 금액 \\
토큰 & WETH, USDC, USDT0, ARB, WBTC (모든 ERC-20 토큰) & 그래프, 대리 지표, 스펜더 라벨에 쓰는 이벤트(표~\ref{tab:usd-tokens}) \\
가격 & DefiLlama 하루 종가; 블록 513{,}457{,}143에서의 소수 자릿수(없음) & 달러 금액 \\
재시작 분포 & $\propto X_u$; 균등 & AWP; EndorseRank, C-PR \\
결합 무게 $\lambda$ & 주된 값 $0.5$; $0.25$, $0.75$ & C-PR 허락 층의 무게, 계획에서 정함 \\
S-PR 재시작 벡터 & 송금 층 점 위의 C-PR ($\lambda=1$), 합이 1이 되게 다시 맞춤 & 보증 씨앗 순간 이동 \\
코호트 규칙 & 다섯 토큰에서 주인 $\geq 3$명에게 받은 0이 아닌 허락; 계약 코드(어떤 토큰의 허락이든 셈) & 매칭 코호트 \\
계정 종류 & \texttt{eth\_getCode}; 코호트는 블록 513{,}216{,}506 & 계약 또는 EOA \\
관측 기간 & 2023-10-01 -- 2026-06-30 & UTC, 양 끝 포함; $T_{\text{obs}}$가 로지스틱 감쇠의 기준점 \\
W0 동결 / 라벨 & 2026-03-31 / 2026-04-01 -- 2026-06-30 & 스펜더 코호트와 트레이더 코호트; 규칙~A(재분석) \\
W1 동결 / 라벨 & 2026-06-30 / 2026-07-01 -- 2026-09-30 & 등록 기간; 규칙~B, 허용 폭 $\delta=0.01$ (재분석) \\
트레이더 규칙 & 계약 코드, 동결일에 그래프의 점, 라벨 기간에 닫기 $\geq 3$번 & GMX~V2 트레이더 코호트 \\
부트스트랩 다시 뽑기 / 시드 & $400$ / $42$; $2{,}000$ / $42$ & 짝지은 계약 부트스트랩; 표~\ref{tab:neutral-label-grid}의 트레이더 라벨 \\
시간 잰 실행 / 준비 실행 & $5$ / $1$ & 실행 시간 재는 방법 \\
규모 단계 시드 & \texttt{contract-benchmark-v1} & SHA256으로 늘 똑같이 뽑는 부분 표본 \\
금액 기울기 점검 & $b/10$과 $10b$ ($T_{\mathrm{obs}}$에서 매칭 코호트의 그래프) & EndorseRank와 AWP의 민감도(표~\ref{tab:robustness-slope}) \\
빅쿼리 예산 & USD~50 & 모든 질의 전에 확인 \\
\bottomrule
\end{tabular}
}
\end{table}

\dissertationSubheading[sec:proxy-appendix]{대리 지표를 실제로 계산하는 방법(참고용)}

대리 지표는 \ref{sec:proxy-formulas}절에서 정의했어요. 아래 목록은 대리 지표마다 이름과 식을 적어요. 그리고 코호트 계약 $v$에 대해 무엇을 세는지 한 줄로 말해요. 모든 대리 지표는 다섯 토큰의 이벤트만 세요.

\begin{itemize}
\item 송금 가족(비슷한 지표 묶음): 들어오는 화살표 수(입차수, 식~\eqref{eq:in-degree})는 $v$로 송금을 보낸 서로 다른 보낸 쪽의 수를 세요. 자기 자신에게 보낸 송금도 넣어요. 받은 양의 합(식~\eqref{eq:in-value})은 그 송금들의 달러 금액을 토큰을 가리지 않고 모두 더해요. 금액은 송금마다 그날의 가격으로 매겨요.
\item 허락 가족: 받은 허락 수(허락 입차수, 식~\eqref{eq:in-approve-degree})는 스펜더 $v$에게 0보다 큰 마지막 허락을 준 주인의 수를 세요. 받은 허락 금액(식~\eqref{eq:in-approve-value})은 그 허락들을 주인과 토큰을 가리지 않고 달러로 모두 더해요. 금액은 기준 시점(앵커)의 가격으로 매기고, 허락 하나하나를 $U=20{,}000$달러에서 잘라요.
\item 시빌 보정 안정성 가족: 들어오는 거래 상대 비율(식~\eqref{eq:counterparty-ratio})은 자기 자신이 아닌 서로 다른 보낸 쪽의 수를, 그들에게서 받은 송금 수로 나눈 값이에요. 송금 기간(식~\eqref{eq:tenure})은 $v$의 것으로 치는 송금들이 걸쳐 있는 날 수예요. 활동한 달 수(식~\eqref{eq:active-months})는 그런 송금이 있었던 달력의 달을 세요.
\end{itemize}

\dissertationSubheading[sec:glossary]{용어 풀이}

아래 용어들은 이 논문 전체에서 여기 적은 뜻으로 써요.

\begin{description}
\item[계정 종류] 계약인지, 외부 소유 계정인지를 말해요. 한 블록에서 \texttt{eth\_getCode}로 읽어요(부록~\ref{sec:account-type-classification}).
\item[활발함] AWP에서 한 주소가 갖는 다시 출발(재시작) 무게예요. 그 주소가 송금을 보낸 주소마다, 그 주소로 보낸 가장 최근 송금의 시간 무게를 구해요. 그리고 이것들을 모두 더한 값이에요.
\item[허락] ERC-20의 허가 하나예요. 이 허가가 있으면 스펜더가 주인 대신 토큰을 옮길 수 있어요.
\item[일치도] 점수와 대리 지표 사이에 순위가 얼마나 맞는지를 말해요. 중간 순위(동점끼리 나눠 갖는 가운데 순위)에 대한 스피어먼의 로($\rho$)와 켄달의 타우($\tau$, 두 줄 세우기가 얼마나 같은 순서인지 나타내는 수)로 재요. 보고한 $\tau$는 모두 동점을 바로잡은 $\tau_b$예요.
\item[구성개념 겹침] 순위 매기는 방법과 대리 지표가 같은 자료를 함께 쓰는 것이에요. 그러면 둘의 일치가 높아져요. 이런 일치는 의도한 구성개념(점수가 재려고 하는 생각) 점검이에요. 그 구성개념 바깥에서 온 증거가 아니에요.
\item[C-PR 끝점] $\lambda=1$과 $\lambda=0$일 때의 C-PR이에요. 앞의 것은 달러 허락 층 하나만 걸어요(허락 하나하나는 $U$에서 잘라요). 뒤의 것은 달러 송금 층 하나만 걸어요. 둘 다 고르게 다시 출발해요(균등 재시작).
\item[에고 네트워크] 다섯 토큰의 이벤트 가운데, 한쪽 끝에 매칭 코호트의 계약이 있는 것들이에요. 한쪽 끝이란 허락의 주인이나 스펜더, 송금의 보낸 쪽이나 받는 쪽을 말해요. 등록 분석에서는 모든 토큰의 이벤트 가운데, 한쪽 끝에 계획의 코호트의 계약이 있는 것들이에요.
\item[보증 그래프] 다섯 토큰의 0보다 큰 마지막 허락을 주인에서 스펜더로 이은 방향이 있는 그래프예요. 엔도스랭크의 입력이고, C-PR의 허락 층이에요.
\item[외부 소유 계정(EOA)] 코드가 없거나 EIP-7702 위임 표시를 가진 주소예요. EOA는 그래프에 주인, 보낸 쪽, 받는 쪽으로 나와요. 하지만 점수를 받거나 평가되는 일은 절대 없어요.
\item[매칭 코호트] 주된 표본이에요. 관측 기간에 서로 다른 주인 세 명 이상에게서 WETH, USDC, USDT0, ARB, WBTC의 0이 아닌 \texttt{Approval}을 받은 계약들이에요(\ref{sec:amount-revision}절과 \ref{sub:cohorts}절). 이 코호트는 등록 분석의 결과를 안 뒤에 정했어요. 다섯 토큰의 허락과 코드 점검으로 정했고, 송금은 보지 않았어요. 등록 분석에서 계획의 코호트는 같은 규칙을 모든 토큰의 허락에 적용했어요. 그리고 송금을 하나라도 모으기 전에 정했어요. 코호트에 들어가는지는 $T_{\text{obs}}$의 어느 그래프에도 달려 있지 않아요. 그래서 코호트 계약이 한 그래프 밖에, 또는 두 그래프 모두 밖에 있을 수 있어요.
\item[짝지은 부트스트랩] 계약을 다시 뽑는 방법이에요. 뽑은 것을 도로 넣고 다시 뽑아요. 이때 모든 계수가 하나의 번호 행렬을 함께 써요. 그래서 계수 사이의 차이를 같은 다시 뽑은 표본에서 검사할 수 있어요.
\item[페이지랭크] 점에 순위를 매기는 스펙트럼 방법(행렬로 순위를 구하는 방법)이에요. 감쇠 계수가 있는 무작위 걷기의 정상 분포(오래 걸었을 때 각 점에 머무는 비율)로 순위를 매겨요.
\item[대리 지표] 계약에 대해 관찰할 수 있는 통계값이에요. 어떤 순위가 재려고 만든 구성개념을 정말 재는지, 이 값에 대어 점검해요.
\item[등록 분석] 계획과 W1의 등록 문서가 정한 분석이에요. 금액은 원래 단위로 쓰고, 모든 ERC-20 토큰을 써요. 그 결과는 얻은 그대로 남겨 두었어요(부록~\ref{sec:map-registered}).
\item[등록 기간(W1)] 2026년 6월 30일에 점수를 동결하고, 2026년 7월부터 9월까지 라벨을 세는 기간이에요. 규칙~B, 그리고 두 화살표 종류 어느 것으로도 만들지 않은 결과에 대한 비교는 그 달들의 로그를 뽑기 전에 등록했어요(\ref{sub:fresh-holdout}절). 고친 점수에는 등록 분석의 결과를 안 뒤에 이 규칙들을 재분석으로 적용해요.
\item[수정 분석] 이 논문의 주된 분석이에요. 등록 분석의 방법과 규칙을 그대로 쓰되, 금액은 달러로 쓰고, 다섯 토큰을 쓰고, 코호트를 그 토큰들로 다시 만들었어요. 이것은 등록 분석의 결과를 안 뒤에 정했어요(\ref{sec:amount-revision}절).
\item[규모 단계] 매칭 코호트에서 SHA256으로 늘 똑같이 뽑은 부분 표본이에요. 그 계약에 닿는 뽑아 둔 화살표도 함께 들어가요.
\item[스펜더 코호트] 뽑아 둔 허락 자료에서, 동결일에 다섯 토큰 가운데 하나의 0보다 큰 마지막 허락을 하나 이상 가진 모든 스펜더 계약이에요. 동결일은 W0에서는 $t_1$(2026년 3월 31일, 23:59:59 UTC)이고, W1에서는 2026년 6월 30일 23:59:59 UTC예요.
\item[순간 이동] 페이지랭크가 $1-d$의 확률로 하는 점프예요. 재시작 분포에서 뽑은 점으로 뛰어요. 재시작 분포는 엔도스랭크와 C-PR에서는 균등하고, AWP에서는 보내는 활동에 따라 무게를 달아요.
\item[트레이더 코호트] 계약 코드를 가지고 있고, 동결일에 그래프의 점이며, 라벨 기간에 포지션을 세 번 이상 닫은 GMX~V2 계정들이에요.
\item[달러 금액] 송금이나 허락의 원래 값을 10을 토큰의 소수 자릿수만큼 거듭제곱한 수로 나누고, 토큰의 DefiLlama 하루 종가를 곱한 값이에요. 송금이면 그 송금의 UTC 날짜의 종가를 쓰고, 마지막 허락이면 기준 시점 날짜의 종가를 써요(\ref{sec:amount-revision}절). 등록 분석은 원래 단위를 썼어요.
\item[금액 기울기] 금액 무게 $V$의 기울기 $b=\ln 3/m$이에요. $m$달러인 금액의 무게가 $1/2$이 되도록 골랐어요($V$는 1을 넘지 않아요). 송금에서는 $m_T=200$이고, 허락에서는 $m_A=62$예요.
\end{description}

% !TEX root = ../main.tex
\section{SQL 틀, 설정, 그리고 더 자세한 설명}
\label{ch:appendix-sql}

다른 사람이 똑같이 다시 해 볼 수 있도록, 이 부록은 몇 가지를 보여 줘요. 두 분석의 그래프가 바탕으로 삼는 SQL 틀, 설정 키, 그리고 계산량과 결과 표의 약속에 대한 설명이에요. 등록 분석(미리 등록한 그대로 한 분석)의 그래프는 분석 계획에 있는 질의 세 개가 만들어요. 수정 분석은 그 질의들이 내놓은 결과에서 자기 표를 잘라 내고, 금액을 달러(USD)로 매겨요(부록~\ref{sec:sql-pairs}). 아래 목록은 줄인 것이에요. 전체 파일은 실험 폴더의 \texttt{sql/}에 있어요(부록~\ref{ch:appendix-eval}). \texttt{\_\_DEST\_\_}처럼 밑줄 두 개로 감싼 자리 표시는 질의를 돌리기 전에 \texttt{scripts/run\_bq.py}가 다른 값으로 바꿔 넣어요. 그리고 \texttt{@}로 시작하는 이름은 질의 매개변수(질의에 밖에서 넣어 주는 값)예요.

\dissertationSubheading[sec:sql-materialize]{로그를 미리 표로 만들어 두는 틀}

스캔~A와 스캔~B는 같은 질의(\texttt{sql/01\_materialize\_logs.sql})를 서로 다른 기간에 대해 돌려요. 이 질의는 공개 로그 표를 한 번만 읽어요. 그리고 로그마다 종류 꼬리표를 붙이고, 이 연구에 필요한 네 종류만 남겨요. 남긴 로그는 달마다 나누고 종류별로 모아 둔 프로젝트 표에 넣어요. GMX 갈래와 Aave 갈래를 줄여서 보이면 다음과 같아요:

\begin{verbatim}
CREATE TABLE `__DEST__`
PARTITION BY TIMESTAMP_TRUNC(block_timestamp, MONTH)
CLUSTER BY kind
AS
WITH src AS (
  SELECT block_timestamp, block_number, log_index,
         LOWER(address) AS address, topics, data
  FROM `bigquery-public-data.goog_blockchain_arbitrum_one_us.logs`
  WHERE block_timestamp >= TIMESTAMP(@start_ts)
    AND block_timestamp < TIMESTAMP(@end_ts)
),
tagged AS (
  SELECT *,
    CASE
      WHEN ARRAY_LENGTH(topics) = 3 AND topics[OFFSET(0)] =
        '0x8c5be1e5ebec7d5bd14f71427d1e84f3dd0314c0f7b2291e5b200ac8c7c3b925'
        THEN 'A'                                  -- ERC-20 Approval
      WHEN ARRAY_LENGTH(topics) = 3 AND topics[OFFSET(0)] =
        '0xddf252ad1be2c89b69c2b068fc378daa952ba7f163c4a11628f55a4df523b3ef'
        THEN 'T'                                  -- ERC-20 Transfer
      WHEN address = '0xc8ee91a54287db53897056e12d9819156d3822fb'
           AND ... THEN 'G'                       -- GMX V2 PositionDecrease
      WHEN address = '0x794a61358d6845594f94dc1db02a252b5b4814ad'
           AND ... THEN 'V'                       -- Aave V3 pool events
    END AS kind
  FROM src
)
SELECT
  block_timestamp, block_number, log_index, kind,
  FROM_HEX(SUBSTR(address, 3)) AS emitter,
  IF(kind IN ('A', 'T'), FROM_HEX(SUBSTR(topics[OFFSET(1)], 27)), NULL) AS a,
  IF(kind IN ('A', 'T'), FROM_HEX(SUBSTR(topics[OFFSET(2)], 27)), NULL) AS b,
  IF(kind IN ('A', 'T'), SAFE.FROM_HEX(SUBSTR(data, 3)), NULL) AS val,
  IF(kind IN ('G', 'V'), topics, NULL) AS raw_topics,
  IF(kind IN ('G', 'V'), data, NULL) AS raw_data
FROM tagged
WHERE kind IS NOT NULL
\end{verbatim}

허락에서는 \texttt{a}가 주인이고 \texttt{b}가 스펜더예요. 송금에서는 \texttt{a}가 보낸 쪽이고 \texttt{b}가 받는 쪽이에요. 두 경우 모두 \texttt{emitter}는 토큰 계약이고, \texttt{val}은 32바이트 금액이에요. 두 topic-0 해시는 \texttt{Approval(address,address,uint256)}의 해시와 \texttt{Transfer(address,address,uint256)}의 해시를 keccak256으로 계산한 값이에요\parencite{vogelsteller2015erc20}. 토픽이 정확히 세 개인 로그만 남기면, ERC-20 로그는 남고 ERC-721 로그는 빠져요. ERC-721 로그는 두 해시를 똑같이 쓰지만, 토큰 번호를 네 번째 토픽에 넣기 때문이에요. 주소는 16진수 글자열이 아니라 20바이트 값으로 저장해요.

\dissertationSubheading[sec:sql-ego]{에고 네트워크 뽑아내기 틀}

\texttt{sql/03\_ego\_events.sql} 질의는 한 기간에 대해 미리 만들어 둔 로그 표를 읽어요. 그리고 한쪽 끝에 코호트 계약이 있는 허락과 송금만 남겨요. \texttt{contract\_rep.cohort} 표에는 계획의 코호트의 20바이트 주소가 들어 있어요. 이 주소들은 \texttt{scripts/build\_cohort.py}가 넣어요. 줄여서 보이면 다음과 같아요:

\begin{verbatim}
CREATE TEMP FUNCTION u256(v BYTES) AS (
  IFNULL((SELECT SUM(c * POW(256.0, LENGTH(v) - 1 - o))
          FROM UNNEST(TO_CODE_POINTS(v)) AS c WITH OFFSET o), 0.0)
);

CREATE OR REPLACE TABLE `dissertation-bq.contract_rep.approvals___SUFFIX__`
PARTITION BY TIMESTAMP_TRUNC(block_timestamp, MONTH)
AS
SELECT
  l.block_timestamp, l.block_number, l.log_index,
  l.emitter AS token_address,
  l.a AS owner,
  l.b AS spender,
  u256(l.val) AS value
FROM `dissertation-bq.contract_rep.__SRC__` AS l
WHERE l.kind = 'A'
  AND l.block_timestamp >= TIMESTAMP(@start_ts)
  AND l.block_timestamp < TIMESTAMP(@end_ts)
  AND (l.a IN (SELECT address FROM `dissertation-bq.contract_rep.cohort`)
       OR l.b IN (SELECT address FROM `dissertation-bq.contract_rep.cohort`));

-- transfers___SUFFIX__: the same with kind = 'T', from_address = a and
-- to_address = b, and zero amounts dropped:
--   AND LTRIM(TO_HEX(l.val), '0') != ''
\end{verbatim}

함수 \texttt{u256}은 32바이트 금액을 원래 단위의 배정밀도 실수(double)로 읽어요. 계획의 질의들은 토큰의 소수 자릿수를 적용하지 않아요. 금액이 0인 허락은 남겨요. 금액이 0인 허락은 허락을 지우는 일이기 때문이에요. 반면 금액이 0인 송금은 버려요. 뒤에 붙는 이름(접미사)은 관측 기간이면 \texttt{ego}, 2026년 7월부터 9월까지이면 \texttt{w1ego}예요. 수정 분석은 이 질의를 다시 돌리지 않아요. 수정 분석은 모든 토큰에 대한 이 표들과 계획의 코호트를 읽고, 그 가운데 다섯 토큰과 수정 코호트에 관한 부분만 남겨요. 수정 코호트는 계획의 코호트에 모두 들어 있어요(부분 집합이에요). 다섯 토큰에서 주인 세 명이 허락한 계약이라면, 어떤 토큰에서 주인 세 명이 허락한 계약이기도 하기 때문이에요.

\dissertationSubheading[sec:sql-pairs]{마지막 허락과 쌍 표를 만드는 틀}

등록 분석에서는 \texttt{sql/04\_anchor\_pairs.sql} 질의가 한 기준 시점(앵커)까지의 에고 이벤트를 두 그래프의 화살표 목록으로 바꿔요. 줄여서 보이면 다음과 같아요:

\begin{verbatim}
CREATE TEMP FUNCTION sig(age_days FLOAT64) AS
  (1.0 / (1.0 + EXP(0.01 * (age_days - 180.0))));
CREATE TEMP FUNCTION vt(x FLOAT64) AS
  (IF(x >= 40.0, 1.0, 2.0 / (1.0 + EXP(-x)) - 1.0));

CREATE OR REPLACE TABLE `dissertation-bq.contract_rep.allow_triples___TAG__` AS
SELECT token_address, owner, spender, value, block_timestamp,
       TIMESTAMP_DIFF(TIMESTAMP(@anchor_ts), block_timestamp, MICROSECOND)
         / 8.64e10 AS age_days
FROM (
  SELECT *, ROW_NUMBER() OVER (
           PARTITION BY token_address, owner, spender
           ORDER BY block_number DESC, log_index DESC) AS rn
  FROM `dissertation-bq.contract_rep.approvals_ego`
  WHERE block_timestamp >= TIMESTAMP(@start_ts)
    AND block_timestamp <= TIMESTAMP(@anchor_ts)
)
WHERE rn = 1 AND value > 0;

CREATE OR REPLACE TABLE `dissertation-bq.contract_rep.allow_pairs___TAG__` AS
SELECT owner, spender,
       COUNT(*) AS n_tokens,
       SUM(sig(age_days) * vt(value)) AS w_er,   -- EndorseRank
       SUM(value) AS w_raw,                      -- C-PR allowance layer
       MAX(sig(age_days)) AS max_sig,
       MAX(block_timestamp) AS last_ts
FROM `dissertation-bq.contract_rep.allow_triples___TAG__`
GROUP BY owner, spender;

-- transfer_pairs___TAG__: per (from_address, to_address) over transfers_ego
-- up to the anchor, with
--   SUM(sig(age_days) * vt(value)) AS w_awp    -- AWP
--   SUM(value * sig(age_days))     AS w_raw    -- C-PR transfer layer
--   MAX(sig(age_days))             AS max_sig  -- AWP activeness
\end{verbatim}

함수 \texttt{sig}는 식~\eqref{eq:logistic-decay}의 로지스틱 감쇠(오래된 일일수록 가볍게 세는 무게)예요. 이때 $k=0.01$이고 $t_0=180$일이에요. \texttt{vt}는 식~\eqref{eq:value-transform}의 값 변환(금액이 아무리 커도 1을 넘지 않게 바꾸는 식)이에요. 여기서는 계획에서 정한 기울기 $b=1$을 원래 단위에 그대로 써요. 그리고 $z=40$부터는 값을 1로 정해 두어요. 그 자리부터는 배정밀도로 계산해도 값이 1과 같기 때문이에요. 한 토큰, 한 주인, 한 스펜더에 대한 허락 가운데 기준 시점까지의 마지막 허락은 블록 번호가 가장 큰 것이에요. 같은 블록 안에서는 로그 번호(log index)가 가장 큰 것이에요. 이 마지막 허락은 금액이 0보다 클 때만 남겨요. 그래서 허락 취소(허락을 거두어들임)가 있으면 그 토큰은 그 쌍에서 빠져요. 쌍 표의 각 행에는 모든 방법의 무게가 함께 들어 있어요. 그래서 한 기준 시점의 모든 점수는 같은 행들에서 나와요. 꼬리표는 $T_{\text{obs}}$이면 \texttt{tobs}, W0의 동결일(점수를 고정하는 날)이면 \texttt{t1}이에요.

수정 분석은 이 틀들에 달러로 값을 매기는 단계를 더해요(\ref{sec:amount-revision}절). \texttt{scripts/run\_usd\_pipeline.py}는 다섯 토큰과 그 소수 자릿수, 그리고 하루 가격 표를 \texttt{data/0-usd-token-list-and-prices/}에서 읽어 \texttt{u\_tokens} 표와 \texttt{u\_prices} 표에 넣어요. \texttt{sql/19\_usd\_load.sql}은 이 표들에 20바이트 주소를 키로 붙여요. 그다음 \texttt{sql/12\_usd\_ego.sql}이 수정 코호트를 만들고, 에고 이벤트를 이 표들과 이어 붙여요. 그리고 \texttt{sql/14\_usd\_anchor\_pairs.sql}이 기울기와 상한을 질의 매개변수로 받아 쌍 표를 만들어요. 줄여서 보이면 다음과 같아요:

\begin{verbatim}
-- sql/12_usd_ego.sql
CREATE OR REPLACE TABLE `dissertation-bq.contract_rep.u_cohort` AS
SELECT spender AS address, COUNT(DISTINCT owner) AS owners_nonzero
FROM `dissertation-bq.contract_rep.approvals_ego`
WHERE block_timestamp >= TIMESTAMP(@start_ts)
  AND block_timestamp < TIMESTAMP(@end_ts)
  AND value > 0
  AND token_address IN
        (SELECT token FROM `dissertation-bq.contract_rep.u_tokens`)
  AND spender IN (SELECT address FROM `dissertation-bq.contract_rep.cohort`)
GROUP BY spender
HAVING COUNT(DISTINCT owner) >= 3;

CREATE OR REPLACE TABLE `dissertation-bq.contract_rep.u_transfers` AS
SELECT x.block_timestamp, x.token_address, x.from_address, x.to_address,
       x.value / POW(10.0, t.decimals) * p.price_usd AS usd
FROM `dissertation-bq.contract_rep.transfers_ego` AS x
JOIN `dissertation-bq.contract_rep.u_tokens` AS t
  ON t.token = x.token_address
LEFT JOIN `dissertation-bq.contract_rep.u_prices` AS p
  ON p.token = x.token_address AND p.day = DATE(x.block_timestamp)
WHERE x.from_address IN
        (SELECT address FROM `dissertation-bq.contract_rep.u_cohort`)
   OR x.to_address IN
        (SELECT address FROM `dissertation-bq.contract_rep.u_cohort`);
-- u_approvals: the same filter on approvals_ego, amount = value / 10^decimals

-- sql/14_usd_anchor_pairs.sql
CREATE TEMP FUNCTION vt(z FLOAT64, b FLOAT64) AS
  (IF(b * z >= 40.0, 1.0, 2.0 / (1.0 + EXP(-b * z)) - 1.0));
-- u_allow_triples___TAG__: the latest approvals of 04_anchor_pairs.sql, taken
-- from u_approvals and valued at the price of the anchor's day:
--   r.amount * p.price_usd AS usd ... JOIN u_prices AS p
--   ON p.token = r.token_address AND p.day = DATE(TIMESTAMP(@anchor_ts))
CREATE OR REPLACE TABLE `dissertation-bq.contract_rep.u_allow_pairs___TAG__` AS
SELECT owner, spender,
       COUNT(*) AS n_tokens,
       SUM(sig(age_days) * vt(usd, @b_allowance)) AS w_er,  -- EndorseRank
       SUM(LEAST(usd, @cap_allowance)) AS w_raw,   -- C-PR allowance layer
       MAX(sig(age_days)) AS max_sig,
       MAX(block_timestamp) AS last_ts
FROM `dissertation-bq.contract_rep.u_allow_triples___TAG__`
GROUP BY owner, spender;
-- u_transfer_pairs___TAG__: per (from_address, to_address) over u_transfers
-- up to the anchor, rows with usd IS NULL left out, with
--   SUM(sig(age_days) * vt(usd, @b_transfer)) AS w_awp   -- AWP
--   SUM(usd * sig(age_days))                 AS w_raw   -- C-PR transfer layer
--   MAX(sig(age_days))                       AS max_sig -- AWP activeness
\end{verbatim}

송금은 그 블록이 있는 UTC 날짜의 가격으로 값을 매겨요. 마지막 허락은 기준 시점 날짜의 가격으로 값을 매겨요. 화살표가 기준 시점에 살아 있는 허락을 나타내기 때문이에요. 매개변수는 \texttt{constants.json}에서 가져와요. 1달러당 $b_A=\ln 3/62$와 $b_T=\ln 3/200$이고, 상한은 $U=20{,}000$달러예요. 이 값들은 $t_1$까지의 이벤트로 계산했어요(\texttt{sql/13\_usd\_constants.sql}). 무제한 허락(금액 한도가 없는 허락)은 $bz\geq 40$에 이르러요. 그래서 엔도스랭크에는 $V=1$로, C-PR의 허락 층에는 $U$로 들어가요. 토큰의 첫 가격보다 앞선 날의 송금은 값이 비어(null) 있게 되어 빠질 거예요. 하지만 남긴 송금 가운데 가격이 없는 것은 하나도 없어요(표~\ref{tab:usd-tokens}). 같은 질의는 기울기를 10배로 키운 경우와 10분의 1로 줄인 경우의 엔도스랭크 무게와 AWP 무게도 써 둬요. 표~\ref{tab:robustness-slope}의 금액 기울기 점검은 이것으로 계산해요. \texttt{sql/15}부터 \texttt{sql/18}까지의 질의는 대리 지표, 라벨, 점 번호, 설명을 수정한 표에서 계산해요. 등록한 표에서 \texttt{sql/05}, \texttt{sql/06}, \texttt{sql/08}, \texttt{sql/10} 질의가 하는 것과 같은 방식이에요. 다만 받은 양의 합은 달러로 더해요.

\dissertationSubheading[sec:config-keys]{설정 키}

처리 과정의 설정은 버전을 관리하는 파일 하나(\texttt{config/contract\_reputation.yaml})에 들어 있어요. 스크립트들은 시작할 때 이 파일을 읽어요. 환경 변수 \texttt{CONTRACT\_REP\_VARIANT}가 어느 분석을 할지 골라요. 이 변수를 정하지 않으면, 스크립트들은 아래 키로 등록 분석을 계산해요. 이 변수를 \texttt{usd}로 정하면, 스크립트들은 \texttt{data/1-raw-blockchain-extracts/graph-tables-usd/}에 있는 수정한 표를 읽고 모든 결과를 \texttt{revised-usd/} 아래에 써요. 수정 분석은 토큰, 소수 자릿수, 가격, 금액 기울기, 상한을 설정 파일에서 가져오지 않아요. 이것들은 \texttt{data/0-usd-token-list-and-prices/}(\texttt{listed\_tokens.csv}, \texttt{daily\_prices\_usd.csv}, \texttt{constants.json})에서 가져와요. 설정의 키 묶음은 다음과 같아요:

\begin{itemize}
\item \texttt{chain\_id}, \texttt{chain\_name}, \texttt{bigquery}: 체인 42161(Arbitrum One), 프로젝트, 데이터셋, 위치, 공개 로그 표, 달러로 정한 예산, 그리고 스캔~A와 스캔~B의 기간이 들어 있어요.
\item \texttt{period}, \texttt{holdout}, \texttt{registered\_window}: 관측 기간과 그 기준 시점 $T_{\text{obs}}$, W0의 동결일과 라벨 기간, 그리고 W1의 동결일, 라벨 기간, 상태, 허용 폭 $\delta$가 들어 있어요.
\item \texttt{events}, \texttt{gmx\_arbitrum}, \texttt{aave\_arbitrum}: \texttt{Approval}과 \texttt{Transfer}의 topic-0 해시, GMX~V2의 이벤트를 내보내는 계약, \texttt{EventLog1} 토픽, \texttt{PositionDecrease} 해시, 청산을 처리하는 계약, 청산의 주문 종류, 최소 닫기 횟수 세 번, 그리고 Aave~V3 풀과 그 이벤트 토픽들이 들어 있어요.
\item \texttt{accounts}, \texttt{cohort}: 코드 점검에 쓰는 RPC 접속 주소(블록체인 노드에 질문을 보내는 주소)와 위임 접두어 \texttt{0xef0100}, 그리고 서로 다른 주인의 최소 수인 세 명이 들어 있어요.
\item \texttt{reputation}: 감쇠 계수(화살표를 따라갈 확률), 허용 오차(멈출 때 쓰는 작은 차이), 반복 횟수 상한, 로지스틱 감쇠의 $k$와 $t_0$, 등록 분석의 값 매개변수 $b=1$, $\lambda$의 주된 값과 민감도 분석용 값, 그리고 규칙~A의 두 부분이 들어 있어요.
\item \texttt{alignment}, \texttt{benchmark}, \texttt{robustness}: 부트스트랩의 다시 뽑기 횟수와 시드, 준비 실행과 시간 잰 실행, 규모 실험의 시드, 그리고 감쇠 계수 값들과 등록한 토큰 점검의 상위 토큰 수가 들어 있어요.
\item \texttt{storage}: 데이터 파일 조각 하나의 최대 크기인 90\,MB가 들어 있어요.
\end{itemize}

감쇠 계수, 로지스틱 감쇠, 금액 기울기, 상한, 토큰 가운데 하나라도 바꾸면 평가를 다시 돌려야 해요. 감쇠 계수 점검, 금액 기울기 점검, 표본 정의 점검도 함께 다시 해야 해요. 그리고 같은 요약에서 표와 그림을 다시 뽑아내야 해요. 그래야 보고한 모든 숫자가 기계가 읽을 수 있는 원본과 맞아요. 수정 분석이 바로 이렇게 다시 돌린 것이에요. \texttt{scripts/run\_usd\_pipeline.py}가 빅쿼리 단계를 순서대로 돌렸고, 모든 로컬 스크립트가 \texttt{CONTRACT\_REP\_VARIANT=usd}로 다시 돌았어요(부록~\ref{sec:pipeline-stages}). \texttt{scripts/extract\_w1.py}는 스캔~B를 시작할 때 이 설정 파일의 SHA-256 해시를 기록해요(부록~\ref{sec:map-artifacts}).

\dissertationSubheading[sec:complexity-notes]{계산량에 대한 설명}

화살표 수를 $m=|E|$라고 써요. 빅쿼리에서 마지막 허락은 $R_a$개의 허락 행에 윈도 함수를 써서 구해요. 이 함수는 토큰, 주인, 스펜더마다 그 안에서 줄을 세워요. 그래서 시간이 많아야 $O(R_a\log R_a)$만큼 걸려요. 송금 쌍은 $R_t$개의 송금 행을 한 번 모아 더하면 되므로 $O(R_t)$만큼 걸려요. 수정 분석에서 토큰 표, 가격 표와 이어 붙이는 일은 각 행을 토큰 하나, 날짜 하나와 맞추는 것이에요. 그래서 이 일로 늘어나는 계산은 $R_a$와 $R_t$에 정비례해요. 로컬 컴퓨터에서는 페이지랭크를 $I$번 반복하면 $O(m+mI)$만큼 걸려요. 앞의 항은 듬성듬성한 연산자(0이 아닌 칸만 적어 둔 행렬)를 만드는 몫이에요. 뒤의 항은 반복 한 바퀴마다 듬성듬성한 곱셈을 한 번 하는 몫이에요(부록~\ref{sec:pr-sparsity}). C-PR은 두 층을 합친 것을 걸어요. 그래서 한 번 반복하는 데 각 층에서 한 번씩 반복하는 것만큼 걸려요. S-PR은 허락 층을 먼저 풀고, 그다음 송금 층을 풀어요. 메모리는 대부분 화살표 목록과 점수 벡터가 차지해요. 그래서 한 번 풀 때 가장 많이 쓰는 메모리는 $m$과 함께 커져요. 표~\ref{tab:benchmark-runtime}에서 엔도스랭크에 대한 AWP의 최대 메모리 비율은 $1.95$예요($1{,}712.6$ 대 $879.7$~MB). 이때 화살표 수의 비율은 $2.23$이에요.

\dissertationSubheading[sec:interpretation-checklist]{결과 표의 약속}

\ref{ch:results}장과 이 부록의 표들은 다음 약속을 따라요:

\begin{enumerate}
\item 모든 $\tau$는 원래 점수와 대리 지표에 대해 구한 켄달의 $\tau_b$예요. 각 값은 95\% 짝지은 부트스트랩 구간과 함께 보고해요. 계약을 다시 뽑는 횟수는 $400$번이고, 시드는 42예요. 다만 두 화살표 종류 어느 것으로도 만들지 않은 결과에 대한 트레이더 비교에서는 $2{,}000$번이에요. 이 비교의 판정 대비 P와 Q에는 97.5\% 구간도 붙어 있어요. 계수 사이의 차이는 점 추정값을 빼서 구하지 않아요. 짝지은 대비 표에서 가져와요.
\item 어떤 방법의 그래프에 없는 코호트 계약은 그 방법에서 0점을 받아요. 그 대신 그 계약을 외톨이 점(화살표가 하나도 없는 점)으로 넣어 두는 것은 표~\ref{tab:tie-sensitivity}의 민감도 분석이에요.
\item 각 대비 표는 그 안의 어떤 대비가 규칙을 판정하는지 밝혀요. 계획이 정한 규칙~A일 수도 있어요. 또는 W1의 등록 문서가 정한 규칙~B와, 두 화살표 종류 어느 것으로도 만들지 않은 결과에 대한 비교일 수도 있어요. 수정 분석의 표에서는 이 규칙들을 등록 분석의 결과를 안 뒤에 정한 점수에 적용해요. 그래서 그 표의 어떤 대비도 등록한 검정이 아니에요. 등록한 판정 대비는 표~\ref{tab:registered-vs-revised}에 있어요. 나머지 대비는 모두 라벨을 본 뒤에 계산했어요.
\item 실행 시간은 시간을 재지 않는 준비 실행을 한 번 한 뒤, 시간을 잰 풀이 프로그램 호출 다섯 번의 실제 걸린 시간을 평균한 값이에요. 호출 한 번은 미리 만들어 둔 화살표 목록을 점수로 바꿔요. 여기에는 점에 번호 붙이기, 듬성듬성한 행렬 조립, 행 정규화(각 행의 합을 1로 맞추기), 재시작 벡터 만들기, 거듭제곱 반복이 들어가요. 뽑아내기와 전처리는 들어가지 않아요. 모든 실행 시간은 그 풀이 그래프의 $|V|$, $|E|$와 함께 나와요. 그리고 같은 설정의 같은 그래프는 모든 표에서 같은 측정값을 보여요(\ref{sec:computational-benchmark}절).
\item 표본 정의 단계(표~\ref{tab:robustness-sample-size})는 평가 집합에 따라 통계량이 어떻게 움직이는지 보여 줘요. 이것은 강건성의 증거가 아니에요. 그리고 그 값은 전체 코호트 표와 비교하지 않아요.
\end{enumerate}

\dissertationSubheading[sec:limitations-reprise]{한계}

위의 질의들은 아비트럼 원(Arbitrum One)에서만 돌렸어요. 그 결과는 스캔~A와 스캔~B를 돌린 2026년 10월 9일 당시의 공개 로그 표에 달려 있어요. 수정 분석은 공개 로그를 다시 읽지 않았어요. 하지만 그 결과는 다른 것에도 달려 있어요. 2026년 10월 10일에 읽은 Chainlink 목록, 블록 $513{,}457{,}143$에서 읽은 토큰 소수 자릿수(블록 $513{,}498{,}909$에서 다시 읽었고 값이 같았어요), 그리고 수정 분석을 위해 가져온 DefiLlama 하루 가격이에요. 에고 필터는 한쪽 끝에 코호트 계약이 있는 이벤트만 남겨요. 그래서 그래프에는 코호트 밖의 두 주소 사이 화살표가 하나도 없어요. 수정한 그래프는 다섯 토큰만 남겨요. 다른 토큰의 행은 개수를 센 뒤 뺐어요(표~\ref{tab:usd-tokens}). 그리고 \texttt{Transfer} 로그 없이 움직이는 기본 이더(native ether)는 WETH로만 들어가요. \ref{ch:results}장의 숫자는 수정 분석의 기계가 읽을 수 있는 요약에서 나와요. 매칭 코호트, W0와 W1의 스펜더 코호트와 트레이더 코호트, 규모 실험 단계, 시빌 모델의 숫자예요. 그 가운데 7월부터 9월까지의 결과는 등록 분석의 결과를 안 뒤에 등록한 규칙을 고친 점수에 적용한 것이에요. 곧 재분석이에요(\ref{sub:fresh-results}절). 등록 분석의 결과는 등록 분석의 요약에서 나와요(부록~\ref{sec:map-registered}).

% !TEX root = ../main.tex
\section{페이지랭크와 순위 상관의 기본 성질}
\label{ch:appendix-theory}

\ref{ch:literature}--\ref{ch:results}장은 여기에 정리한 기본 성질에 기대고 있어요. 이 내용은 널리 알려진 표준적인 것이에요. 페이지랭크의 성질은 \textcite{page1999pagerank}의 연구와 \textcite{langville2006pagerank}의 연구를 따라요. 순위 상관(두 줄 세우기가 얼마나 비슷한지 나타내는 값)은 \textcite{spearman1904proof}의 연구와 \textcite{kendall1938rank}의 연구를 따라요.

\dissertationSubheading[sec:pr-existence]{페이지랭크의 존재와 유일성}

$P$를 $|V|$개의 점 위에 있는 행 확률 행렬(각 행의 수를 모두 더하면 1인 행렬)이라고 해요. 필요하면 식~\eqref{eq:transition}처럼 매달린 점(나가는 화살표가 없는 점)을 고친 뒤의 행렬이에요. $d\in(0,1)$일 때 식~\eqref{eq:google-matrix}의 구글 행렬은 다음과 같아요.
\[
G_{\mathrm{PR}}=d P^{\top}+\frac{1-d}{|V|}\mathbf{1}\mathbf{1}^{\top},
\]
이 행렬은 모든 칸이 양수이고, 열 확률 행렬(각 열의 수를 모두 더하면 1인 행렬)이에요. 양수 행렬에 대한 페론--프로베니우스 정리(Perron--Frobenius theorem)에 따르면, $G_{\mathrm{PR}}$에는 크기(절댓값)가 1인 고윳값(행렬을 곱해도 방향이 바뀌지 않는 벡터가 늘어나는 배수)이 하나뿐이에요. 그리고 그 고윳값에 맞는 고유 벡터 $\mathbf{r}$은 모든 칸이 0보다 커요. 우리는 이 벡터를 $\sum_i r_i=1$이 되도록 크기를 맞출 수 있어요. 따라서 감쇠 계수(화살표를 따라갈 확률)와 매달린 점 규칙을 정해 두면, 방향이 있는 모든 그래프는 페이지랭크 점수를 단 하나만 가져요\parencite{langville2006pagerank}.

이 논문의 재시작 벡터 가운데 두 개는 고르지 않아요. AWP의 재시작 벡터는 활발함에 비례해요. 그래서 아무것도 보내지 않은 주소에서는 0이에요. S-PR의 재시작 벡터는 허락 걷기의 점수를 정규화(합이 1이 되게 맞춤)한 것이에요. 그래서 송금 층의 점 가운데 허락 점수가 없는 점에서는 0이에요. 따라서 두 구글 행렬 모두, 모든 칸이 양수인 행렬이 아니에요. 그래도 고정점(계산을 한 번 더 해도 바뀌지 않는 값)은 여전히 하나뿐이에요. 두 확률 벡터의 차이를 구하면 다시 출발 항이 서로 지워져요. 그리고 이 연산자는 두 벡터 사이의 $\|\cdot\|_1$ 거리를 많아야 $d$배로 만들어요. 하지만 일반적으로 고정점의 모든 칸이 0보다 크지는 않아요. 다시 출발 몫이 없고, 다시 출발하는 점에서 출발한 어떤 길로도 닿지 않는 점은 점수가 0으로 남기 때문이에요. 그래서 표~\ref{tab:tie-sensitivity}의 두 규칙 모두에서, AWP는 송금 그래프 밖에 있는 모든 코호트 계약에 0점을 줘요.

점수가 하나뿐이라는 성질(유일성) 덕분에 엔도스랭크와 AWP를 분명하게 비교할 수 있어요. 매개변수를 정하고 나면, 각 방법은 주어진 화살표 집합에 점수 벡터를 정확히 하나만 줘요. 그래서 $\mathbf{s}_{\mathrm{ER}}$와 $\mathbf{s}_{\mathrm{AWP}}$ 사이의 차이는 모두 두 방법의 화살표와 재시작 분포에서 나와야 해요. 그 차이를 풀이 프로그램(solver) 탓으로 돌릴 수는 없어요.

\dissertationSubheading[sec:pr-continuity]{화살표 무게에 대한 연속성}

$d\in(0,1)$을 고정하면, 페이지랭크는 $P$의 칸 값에 따라 연속적으로(끊기지 않고 조금씩) 변해요. 같은 점들 위에 두 번째 행 확률 행렬 $\tilde P$가 있고, 그 페이지랭크 벡터를 $\tilde{\mathbf{r}}$로 써요. 그러면 다음이 성립해요.
\[
\|\mathbf{r}-\tilde{\mathbf{r}}\|_1\le\frac{d}{1-d}\,\max_i\bigl\|P_{i\cdot}-\tilde P_{i\cdot}\bigr\|_1,
\]
이 상한(넘지 않는 값)은 $|V|$에 따라 달라지지 않아요\parencite{langville2006pagerank}.

$P$에 대한 연속성은 화살표 무게에 대한 연속성과 같지 않아요. 행 $P_{i\cdot}=w_{i\cdot}/o_i$는 나가는 무게의 합 $o_i$가 0보다 큰 동안에만 무게에 따라 연속적으로 움직여요. 어떤 점이 첫 번째 나가는 화살표를 얻거나 마지막 나가는 화살표를 잃으면, 그 점의 행은 갑자기 뛰어요. 고르게 퍼진 매달린 점의 행과, 화살표가 가리키는 이웃들에 몰린 행 사이를 오가요. 관련된 무게가 아무리 작아도 그래요. \ref{sec:pagerank-worked-example}절의 계산 예에서 무게가 $10^{-12}$인 화살표 $D\to A$를 더하면, $\mathbf{r}$이 $\|\cdot\|_1$ 노름(차이의 절댓값을 모두 더한 크기)으로 $0.37$만큼 움직여요. 그리고 $r_A$는 $0.1375$에서 $0.3202$로 올라가요.

감쇠 계수를 여러 값으로 바꿔 보는 실험은 연속성으로 설명할 수 있어요. 이 실험은 $P$를 그대로 두고 $d$만 매끄럽게 바꿔요. 그리고 $d<1$이면 페이지랭크 벡터는 $d$에 따라 연속적으로 변해요. 금액 기울기 점검은 일부만 설명할 수 있어요. $V$의 기울기 $b$를 바꾸면 유한한 모든 금액의 무게가 바뀌어요. 반면 무제한 허락(금액 한도가 없는 허락)은 $V=1$을 그대로 유지해요. 그리고 정확한 계산에서는 화살표가 하나도 더해지거나 빠지지 않아요. 모든 $z>0$에 대해 $V(z)>0$이기 때문이에요. 그래서 $P$는 $b$에 따라 연속적으로 움직여요. 하지만 컴퓨터의 부동소수점 계산에서는 아주 작은 금액만 실은 화살표의 무게가 반올림되어 0이 될 수 있어요(\ref{sec:rank-divergence}절). 그리고 $b$를 바꾸면 그런 화살표가 0이 되거나 0에서 벗어날 수 있어요. 그런데 $b$를 10배 바꾸는 것은 작은 변화가 아니에요. 위의 상한은 이런 변화에 대해 별로 말해 주지 않아요. 그래서 그 점검이 안정적인지는 자료에 달려 있어요(표~\ref{tab:robustness-slope}). 실제 자료에서는 $d\in\{0.75,0.85,0.95\}$에 걸쳐 엔도스랭크와 AWP의 지표 가족(비슷한 지표 묶음) 평균이 조금만 바뀌었어요. 가장 많이 바뀐 것은 AWP의 송금 평균이었어요(\ref{sub:damping}절; 표~\ref{tab:robustness-damping}). 다만 반복 횟수는 $d$가 클수록 늘어났어요.

\dissertationSubheading[sec:pr-sparsity]{희소성과 반복 한 바퀴의 비용}

0이 아닌 화살표가 $m=|E|$개 있으면, 듬성듬성한 곱 $P^{\top}\mathbf{x}$에는 $O(m)$번의 산술 연산이 들어요. 화살표 목록으로 듬성듬성한 연산자를 만드는 일도 $O(m)$이에요. 그래서 $I$번 반복하는 거듭제곱 반복(같은 계산을 되풀이하기)은 $O(m+mI)$예요. 매칭 코호트에서 송금 그래프의 화살표 수는 허락 그래프의 $2.23$배예요(\ref{sec:benchmark-results}절). 그래서 반복 횟수가 같다면, 송금 방법은 시간이 약 $2.23$배 걸릴 거예요. 두 점수는 같은 풀이 프로그램을 써요. 그래서 그 밖의 시간 차이는 알고리즘이 아니라, 반복 횟수와 $|V|$에 걸친 벡터 연산에서 나와요. 실제로 잰 호출에서 엔도스랭크에 대한 AWP의 시간 비율은 $2.17$이고, 반복 횟수는 각각 $92$번과 $69$번이에요(\ref{sec:benchmark-results}절). 그래서 이 그래프들에서는 화살표 하나, 반복 한 바퀴당 걸린 시간이 송금 방법 쪽이 엔도스랭크보다 짧아요.

\dissertationSubheading[sec:spearman-properties]{스피어먼의 \texorpdfstring{$\rho$}{rho}가 가진 성질}

스피어먼의 로($\rho$)는 순위에 대해 계산한 피어슨 상관(두 값이 함께 커지고 작아지는 정도)이에요. 이 값은 $[-1,1]$ 안에 있어요. 그리고 두 변수 가운데 어느 쪽을 엄격하게 증가하는 변환(순서를 그대로 지키는 바꾸기)으로 바꿔도 값이 바뀌지 않아요. $\rho=1$이면 완벽한 단조 일치(한쪽이 커지면 다른 쪽도 늘 커짐)를 뜻해요. $\rho=-1$이면 완벽하게 거꾸로인 일치를 뜻해요. 값이 같은 것(동점)에는 가운데 순위를 줘요. 평가 처리 과정은 동점을 제대로 다루는 표준 구현을 써요. 평판 점수와 대리 지표는 연속적인 값이고, 꼬리가 두꺼워요(아주 큰 값이 가끔 나와요). 그래서 원래 값에 대한 피어슨 상관보다 순위에 바탕을 둔 연관성이 더 잘 맞아요. AWP 기준선 논문도 마찬가지로 순위 통계로 자기 순위를 평가해요\parencite{do2023awp}.

\dissertationSubheading[sec:kendall-properties]{켄달의 \texorpdfstring{$\tau$}{tau}가 가진 성질}

켄달의 타우($\tau$)는 순서가 맞는 쌍의 수 $n_c$와 순서가 어긋나는 쌍의 수 $n_d$를 비교해요. 동점 보정을 하지 않으면 $\tau_a=(n_c-n_d)/n_0$이고, 여기서 $n_0=n(n-1)/2$예요. 이 값은 무작위로 뽑은 두 계약을 두 순위가 같은 순서로 세울 확률에서, 반대 순서로 세울 확률을 뺀 값이에요\parencite{kendall1938rank}. 보고한 모든 $\tau$는 식~\eqref{eq:kendall}의 동점 보정 $\tau_b$예요. 이 값의 분모는 $n_0$에서 어느 한 순위에서라도 동점인 쌍을 뺀 것이에요. 동점이 없으면 두 값은 같아요. 그리고 이 논문의 점수들에서는 동점이 드물어요. 매칭 코호트에서 어떤 점수도 계약 쌍의 $3.2\%$보다 많이 동점을 만들지 않아요. 엔도스랭크와 AWP의 동점 비율은 \ref{sec:rank-divergence}절에 있어요. 그러므로 엔도스랭크와 AWP를 비교할 때 $\tau_b$의 분모는 $n_0$과 2퍼센트 미만으로 달라요. 그리고 \ref{sec:rank-divergence}절의 쌍 개수로 계산한 $\tau_a$는 보고한 $\tau_b$와 2퍼센트 안쪽으로만 달라요(표~\ref{tab:tie-sensitivity}).

대리 지표에 있는 동점도 $\tau_b$에 천장(얻을 수 있는 가장 큰 값)을 만들어요. 들어오는 화살표 수(입차수)나 받은 허락 수(허락 입차수) 같은 정수 대리 지표는 많은 쌍을 동점으로 만들어요. 그래서 그런 대리 지표에 대해 어떤 점수가 다다를 수 있는 가장 큰 $\tau_b$는 동점 쌍의 비율이 클수록 낮아져요. 여기서는 어떤 점수도 쌍의 $3.2\%$보다 많이 동점을 만들지 않아요. 그래서 한 대리 지표에서 비교하는 점수들 사이에서는 이 천장이 거의 다르지 않아요. 따라서 한 대리 지표 안의 짝지은 대비는 이 천장 때문에 비뚤어지지 않아요. 하지만 서로 다른 대리 지표 사이의 대비는 비뚤어질 수 있어요. 대리 지표들의 동점이 서로 다른 만큼 그래요.

$n$이 크면 $\tau$와 $\rho$는 보통 부호가 같지만, 크기는 달라요. 엔도스랭크나 AWP를 대리 지표 하나와 짝지은 14개 쌍에서(표~\ref{tab:alignment-transfer}부터~\ref{tab:alignment-sybil-stability}까지), 모든 $|\tau|$는 $0.07$을 넘어요. 비율 $\rho/\tau$는 약 $1.2$와 $1.8$ 사이에 있어요. 그리고 둘 가운데 늘 $|\rho|$가 더 커요. 계수의 부호가 서로 반대인 쌍은 하나도 없어요. 우리는 AWP 기준선처럼 둘 다 보고해요\parencite{do2023awp}. 우리의 주된 요약 통계량은 짝지은 부트스트랩 구간을 붙인 지표 가족 평균 $\tau$예요\parencite{efron1979bootstrap}.

\dissertationSubheading[sec:construct-validity-note]{구성 타당도에 대한 설명}

\textcite{cronbach1955construct}의 연구는 기준 관련 타당도(정답 기준과 얼마나 맞는지)와 구성 타당도(점수가 재려는 것을 정말 재는지)를 구별해요. 계약 평판에 대한 확실한 정답 라벨은 없어요. 그래서 이 연구는 구성개념(점수가 재려고 하는 생각)을 중심에 둔 방법을 택해요. 각 점수는 그 개념 정의에서 나오는 통계와 맞춰 봐요. 예를 들어 엔도스랭크는 허락을 받는 것과, AWP는 들어오는 송금을 받는 것과 맞춰 봐요. 점수 자신의 화살표로 만든 통계와 맞는지 보는 것은 의도한 구성개념 점검이에요. 이 점검은 만든 방식 때문에 일부는 저절로 맞게 되어 있어요. 각 점수는 같은 구성개념이 나중에 새로 생긴 것과도 맞춰 봐요. 곧 새 허락 쌍(새로 생긴 주인–스펜더 허락)과 새 송금자(새로 송금을 보내온 주소)예요. 또 그 점수가 매끄럽게 다듬는 원시 차수(화살표를 그냥 센 수)와도 맞춰 봐요. 그리고 W1과 W0 기간에서는 두 화살표 종류 어느 것으로도 만들지 않은 트레이더 결과와도 맞춰 봐요. 한 구성개념 안에서 잘 맞으면, 점수를 의도한 대로 읽어도 된다는 근거가 돼요. 다만 앞에서 말한 겹침은 생각해야 해요. 서로 다른 구성개념 사이에서 잘 맞지 않는다고 해서, 그것만으로 실패라는 뜻은 아니에요. 한 구성개념을 위해 만든 점수가 다른 구성개념을 따라갈 필요는 없기 때문이에요.

\dissertationSubheading[sec:numerical-summary-sheet]{주요 값을 보고한 곳}

아래 찾아보기는 수정 분석의 주요 값마다 그 값을 밝힌 표나 절을 알려 줘요. 마지막 행들은 금액 수정과 등록 분석(미리 등록한 그대로 한 분석)을 가리켜요.

\begin{center}
\fitwidth{%
\begin{tabular}{ll}
\toprule
값 & 보고한 곳 \\
\midrule
실행 시간, 표준편차(SD), 메모리, 반복 횟수, $|V|$, $|E|$ (점수 네 개, 전체 그래프) & 표~\ref{tab:benchmark-runtime} \\
규모 실험 단계(매칭 코호트의 부분 표본) & 표~\ref{tab:benchmark-scaling} \\
벤치마크에서 시간을 잰 모든 실행 & 표~\ref{tab:benchmark-runs} \\
지표 가족 평균 $\tau$와 95\% 신뢰구간, 모든 페이지랭크 & 표~\ref{tab:alignment-hybrid}과~\ref{tab:endorserank-restarts} \\
C-PR과 S-PR의 짝지은 $\Delta\tau$ 대비 & 표~\ref{tab:tau-diff-hybrid} \\
방법 사이의 $\tau$; 그래프 밖의 계약을 외톨이 점으로 둔 경우 & 표~\ref{tab:tie-sensitivity} \\
감쇠 계수, 금액 기울기, 표본 정의 점검 & 표~\ref{tab:robustness-damping}, \ref{tab:robustness-slope}, \ref{tab:robustness-sample-size} \\
고르게 다시 출발하기와 활동 재시작을 쓴 EndorseRank & 표~\ref{tab:endorserank-restarts} \\
스펜더 홀드아웃 W0 & 표~\ref{tab:holdout-spenders}, \ref{tab:holdout-diff}과~\ref{tab:holdout-receiving} \\
트레이더 실행 W0 (이익 라벨) & 표~\ref{tab:holdout-traders} \\
기간 W1, 고친 점수 & 표~\ref{tab:fresh-holdout}, \ref{tab:fresh-contrasts}, \ref{tab:fresh-traders}과~\ref{tab:fresh-posthoc} \\
두 화살표 종류 어느 것으로도 만들지 않은 결과 & 표~\ref{tab:neutral-label-contrasts}, \ref{tab:neutral-label-recipients}와~\ref{tab:neutral-label-grid} \\
시빌 비용 모델 & \ref{sec:sybil-cost-model}절 \\
EndorseRank와 AWP 사이의 비교, 대리 지표별 결과 & 부록~\ref{ch:appendix-er-awp} \\
코호트의 계정 종류 & 표~\ref{tab:account-types} \\
뽑아낸 이벤트, 토큰, 분기별 GMX~V2 활동 & 부록~\ref{ch:appendix-data} \\
다섯 토큰, 남긴 행과 뺀 행, 가격 & 표~\ref{tab:usd-tokens} \\
관측 기간 & 2023-10-01 -- 2026-06-30 (표~\ref{tab:parameter-summary}) \\
\midrule
금액 수정(달러, 다섯 토큰, 기울기, 상한) & \ref{sec:amount-revision}절 \\
등록 분석: 판정 대비와 판정 & 부록~\ref{sec:map-registered}, 표~\ref{tab:registered-vs-revised} \\
\bottomrule
\end{tabular}
}
\end{center}

여기 적은 값은 모두 기계가 읽을 수 있는 요약에 들어 있고, 옆에 적은 표나 절에도 나와요. 등록 분석의 결과는 따로 마련한 요약에 들어 있어요.

% !TEX root = ../main.tex
\section{자료 설명}
\label{ch:appendix-data}

이 부록은 \ref{ch:results}장의 바탕이 되는, 뽑아낸 자료를 설명해요. 다루는 것은 다음과 같아요. 에고 네트워크(코호트 계약에 닿는 화살표만 모은 그래프)를 뽑아낸 자료의 양을 다뤄요. 수정 분석(등록 뒤에 금액을 달러로 바꿔 다시 한 분석)이 남기는 행과, 그 분석의 두 그래프 크기를 다뤄요. 다섯 토큰과 그 가격을 다뤄요. 송금(transfer)과 허락(allowance, 토큰을 얼마까지 대신 써도 된다는 허락)의 미국 달러 금액, 그리고 주인(토큰 주인)이 무제한 허락(금액 한도가 없는 허락)을 얼마나 자주 해 주는지를 다뤄요. 마지막으로 GMX~V2 스마트 계약(블록체인 위에서 저절로 실행되는 프로그램) 계정의 분기별 활동을 다뤄요. 코호트(묶음)들의 계정 종류는 표~\ref{tab:account-types}에 있어요. 자료를 뽑아낸 과정 자체는 부록~\ref{sec:bq-extraction}에서 설명해요. 금액을 고친 일은 \ref{sec:amount-revision}절에서 설명해요.

\dissertationSubheading[sec:data-events]{뽑아낸 이벤트와 그래프}

관측 기간의 에고 자료는 계약 $27{,}844$개로 이루어진 계획의 코호트를 위해 뽑아냈어요. 이 자료에는 모든 ERC-20 토큰의 허락 로그(계약이 남기는 기록) $389{,}866{,}882$개와 송금 로그 $1{,}538{,}404{,}067$개가 들어 있어요. 허락 로그는 금액이 0인 것도 포함해요. 송금 로그는 금액이 0이 아닌 것만이에요. 로그마다 한쪽 끝에 그 코호트의 계약이 있어요. 수정 분석은 이 가운데 허락 $224{,}954{,}202$개와 송금 $1{,}149{,}168{,}538$개를 남겨요. 다섯 토큰의 것이면서, 한쪽 끝에 수정 코호트의 계약이 있는 것들이에요. 뽑아낸 로그의 $57.7\%$와 $74.7\%$예요. W1의 7월부터 9월까지의 로그 가운데서는 허락 $12{,}209{,}894$개와 송금 $60{,}182{,}192$개를 남겨요. 표~\ref{tab:usd-tokens}에 토큰마다의 행 수와 빠진 행 수가 있어요. 남긴 이벤트로 만든 $T_{\text{obs}}$ 시점의 허락 그래프(점과 화살표로 이은 그림)에는 점(노드) $6{,}221{,}665$개와 화살표(엣지) $11{,}210{,}114$개가 있어요. 송금 그래프에는 점 $7{,}884{,}602$개와 화살표 $25{,}001{,}216$개가 있어요. 이 화살표 수는 엔도스랭크(EndorseRank)나 AWP(적응형 가중 페이지랭크)의 무게(가중치)가 0보다 큰 쌍만 센 것이에요. 금액으로 구한 금액 무게가 배정밀도(double) 계산에서 0이 되는 쌍은 세지 않았어요. 그런 쌍은 허락 쌍 $9{,}002$개와 송금 쌍 $17{,}676$개예요. 반면 점 수에는 그런 쌍만 닿는 주소 $2{,}602$개와 $2{,}246$개가 들어 있어요. W0의 동결일(점수를 고정하는 날)에 두 그래프의 쌍 표에는 쌍이 $10{,}650{,}258$개와 $23{,}680{,}692$개 들어 있어요. 이 수는 그 거르기를 하지 않고 센 것이에요. 허락 그래프의 화살표 수는 송금 그래프의 절반도 안 돼요. 허락 그래프는 토큰, 주인, 스펜더(spender, 허락을 받아 대신 쓰는 쪽)의 조합마다 마지막 허락(가장 최근 허락)을 하나만 남기고, 거둔 허락은 버려요. 반면 송금 그래프는 그 기간에 다섯 토큰의 송금이 오간 모든 보낸 쪽(송신자)과 받는 쪽(수신자)의 쌍으로 만들어요.

어떤 계약이 코호트에 드는지는 $T_{\text{obs}}$ 시점의 두 그래프 어느 쪽에도 달려 있지 않아요. 코호트의 모든 계약은 그 기간에 주인 세 명 이상에게서 다섯 토큰의 0이 아닌 허락을 받았어요. 하지만 $15{,}107$개 가운데 $523$개는 $T_{\text{obs}}$에 0보다 큰 마지막 허락을 받아 둔 것도, 해 준 것도 없어요. 그래서 이 계약들은 허락 그래프 밖에 있어요. $2{,}670$개는 송금 그래프 밖에 있고, $80$개는 두 그래프 모두의 밖에 있어요. \ref{ch:results}장은 이런 계약이 어떤 방법의 그래프에 없으면, 그 방법에서 0점을 줘요. 표~\ref{tab:tie-sensitivity}는 그 대신 이런 계약을 외톨이 점(화살표가 하나도 없는 점)으로 남기는 규칙을 보여 줘요.

\dissertationSubheading[sec:data-tokens]{토큰과 허락 금액}

수정 분석은 다섯 토큰의 허락과 송금을 남겨요(\ref{sec:amount-revision}절). 토큰은 두 조건을 채우면 자격을 얻었어요. 첫째, 아비트럼 원(Arbitrum One)에 대한 Chainlink의 데이터 피드 목록에, 시장 위험이 낮음이나 보통으로 매겨진 시장 가격 피드(블록체인에 올라오는 가격 정보)가 있어야 해요. 그 피드는 다음 가운데 하나의 가격을 매겨야 해요. 토큰 자신, 그 토큰이 일대일로 감싼 자산, 또는 그 토큰이 공식 브리지(체인 사이에 토큰을 옮겨 주는 다리)로 옮겨 온 복사본일 때 원래의 이더리움 토큰이에요\parencite{chainlinknddatafeeds}. 둘째, 그 토큰이 관측 기간 자료에 나와야 해요. 그 목록이 짝지어 둔 아비트럼 원 토큰 $133$개 가운데 $64$개가 자격을 얻었고, 관측 기간 자료에도 나와요. 이 가운데 $46$개는 DefiLlama에서 하루 가격을 얻을 수 있어요\parencite{defillamanddocs}. 분석은 이 $46$개 가운데 행이 가장 많은 다섯 개를 남겨요. 행 수는 계획의 코호트의 에고 표에서 $T_{\text{obs}}$까지의 송금과 $T_{\text{obs}}$ 시점의 0보다 큰 마지막 허락을 세어 구해요. 그 다섯은 WETH, 네이티브(아비트럼 원에서 바로 발행한) USDC, USDT0, ARB, WBTC예요. 이 다섯 토큰은 모두 합쳐 전체 토큰 행의 $75.75\%$를 차지해요. 그리고 다섯 토큰 모두 위험이 낮음으로 매겨진 피드로 가격이 매겨져요. 표~\ref{tab:usd-tokens}는 각 토큰에 대해 다음을 보여 줘요. 주소, 소수 자릿수, 피드가 있어요. 계획의 코호트의 에고 표에 든 허락 로그와 송금 로그의 수, 그리고 수정 분석이 남긴 로그의 수가 있어요. 이 로그 수는 허락과 송금을 함께 센 것이에요. 또 남긴 송금의 미국 달러 거래량과, 하루 가격의 바탕이 된 자료가 있어요.

\IfFileExists{results/tables/usd-tokens.tex}{% Generated by experiments/scripts/export_latex.py -- do not edit by hand.
% Source: experiments/data/0-usd-token-list-and-prices/listed_tokens.csv
% Source: experiments/data/2-processed-tables-and-evaluations/revised-usd/describe/describe.json
% Source: experiments/data/0-usd-token-list-and-prices/chainlink_check.json
\begin{table}[htbp]
\centering
\caption[수정 분석의 다섯 토큰과 남긴 행, 그리고 가격.]{수정 분석에 쓴 다섯 토큰이에요. 토큰마다 다음을 적었어요. 먼저 체인링크(Chainlink) 가격 피드(시장 위험 등급)를 적었어요. 계획의 코호트가 가진 에고 표(에고 네트워크를 담은 표)에 관측 기간 동안 그 토큰의 행(허락과 송금)이 몇 개 있는지, 그리고 수정 분석이 그 가운데 몇 개를 남겼는지도 적었어요. 남긴 송금의 달러 거래 금액도 적었어요. 가격 시계열에서 두 번째 DefiLlama 키로 채운 날, 전날 값으로 채운 날, 하루짜리 오류로 보고 바꾼 날의 수도 적었어요(1{,}096일 가운데). 마지막으로, 38개 날짜에 블록체인에서 읽은 Chainlink 피드와 DefiLlama 가격 사이 절대 차이의 중앙값과 가장 큰 값을 퍼센트로 적었어요. 계획의 코호트가 가진 에고 표에는 관측 기간에 허락 389{,}866{,}882건과 송금 1{,}538{,}404{,}067건이 있어요. 이 가운데 허락 164{,}202{,}859건과 송금 364{,}468{,}917건은 다른 토큰의 것이라서 뺐어요(관측 기간과 W1의 에고 표에 있는 다른 토큰은 142{,}958개예요). 다섯 토큰의 허락 709{,}821건과 송금 24{,}766{,}612건은 양 끝 어느 쪽도 수정 코호트의 계약이 아니어서 뺐어요. 남긴 송금 가운데 그 토큰의 첫 가격보다 앞선 것은 0건이에요. $T_{\mathrm{obs}}$에서 0보다 큰 마지막 허락 12{,}590{,}613개 가운데 5{,}763{,}757개는 무제한 허락이에요. 그리고 6{,}776{,}289개는 C-PR 허락 층의 상한 $U=20{,}000$달러에 닿아요.}
\label{tab:usd-tokens}
\fitwidth{%
\begin{tabular}{llrlrrrrrrr}
\toprule
토큰 & 주소 & 소수 자릿수 & Chainlink 피드 & 행 수, 계획 & 행 수, 남김 & 송금 USD & 2번째 키 & 채움 & 오류 & CL 차이 (\%) \\
\midrule
WETH & \texttt{0x82af\ldots{}bab1} & 18 & ETH / USD (낮음) & 586{,}587{,}366 & 569{,}011{,}073 & $8.4\times10^{11}$ & 319 & 0 & 0 & 0.06 / 0.72 \\
USDC & \texttt{0xaf88\ldots{}5831} & 6 & USDC / USD (낮음) & 400{,}716{,}805 & 396{,}280{,}336 & $9.9\times10^{11}$ & 0 & 0 & 0 & 0.01 / 0.26 \\
USDT0 & \texttt{0xfd08\ldots{}cbb9} & 6 & USDT / USD (낮음) & 233{,}074{,}004 & 231{,}428{,}275 & $2.7\times10^{11}$ & 481 & 0 & 0 & 0.03 / 0.27 \\
ARB & \texttt{0x912c\ldots{}6548} & 18 & ARB / USD (낮음) & 90{,}752{,}065 & 89{,}229{,}800 & $8.5\times10^{10}$ & 0 & 0 & 0 & 0.08 / 0.32 \\
WBTC & \texttt{0x2f2a\ldots{}5b0f} & 8 & WBTC / USD (낮음) & 88{,}468{,}933 & 88{,}173{,}256 & $2.2\times10^{11}$ & 308 & 0 & 0 & 0.06 / 0.24 \\
\bottomrule
\end{tabular}
}
\end{table}
}{}

UTC 기준 하루의 토큰 가격은 DefiLlama의 하루 종가(그날 마지막 가격)예요(\ref{sec:amount-revision}절). 아비트럼 원 주소로 찾은 DefiLlama의 가격 기록은 WETH와 WBTC는 2024년 8월에, USDT0는 2025년 1월에 시작해요. 그래서 세 토큰의 $1{,}096$일 가운데 각각 $319$일, $308$일, $481$일은 같은 자산의 CoinGecko 식별자, 곧 두 번째 키로 찾은 가격을 써요. 두 기록이 함께 있는 날들에서, 이 가격은 아비트럼 기록과 중앙값으로 $0.1\%$보다 적게 달라요. USDC와 ARB는 두 번째 키가 필요 없어요. 다섯 토큰 가운데 어느 것에도 전날 값으로 채운 날이나, 자료 오류로 보고 바꾼 값이 없어요. 그리고 남긴 송금 가운데 그 토큰의 첫 가격보다 앞선 것은 없어요. 토큰마다 $38$개 날짜에 체인 위에서 읽은 Chainlink 피드와 비교하면, DefiLlama 가격은 토큰마다 중앙값으로 $0.01\%$에서 $0.08\%$만큼 달라요. 가장 크게 다른 것은 $0.72\%$로, 2024년 1월 15일의 WETH예요.

표~\ref{tab:top-tokens}에는 토큰 $27$개가 있어요. 행이 가장 많은 스무 개에, 등록한 토큰 점검의 금액순 목록에서 온 일곱 개를 더한 것이에요. 수정 분석은 앞의 다섯 개를 남겨요. 여섯 번째인 브리지로 옮겨 온 USDC는 자격도 있고 가격도 있지만 빠져요. 스무 개 가운데 일곱 개는 피드 규칙에 따른 자격이 없어요. NPC, UXUY, AIUS가 그 가운데 들어요. 마지막 두 열은 등록 분석(미리 등록한 그대로 한 분석)의 토큰 점검에 쓴 두 목록을 표시해요. 그 점검은 토큰 스무 개로 두 그래프를 다시 만들었어요. 스무 개는 한 번은 행 수로, 한 번은 원래 금액을 더한 값으로 골랐어요(\texttt{sql/11\_top\_tokens.sql}). 수정 분석에서는 이 점검을 금액 기울기 점검으로 바꿨어요(\ref{sub:top-tokens}절). 두 목록에 함께 든 토큰은 $13$개예요. 그 분석에서는 소수 자릿수를 생각하지 않고 원래 금액을 더했어요. 그리고 무제한 허락은 원래 단위(소수점 조정을 하지 않은 가장 작은 단위)로 $2^{256}-1$로 셌어요. 그래서 금액순 목록에는 행이 비교적 적은 토큰 일곱 개가 들어갔어요. 그 가운데에는 WETH를 Aave에 맡기고 받는 영수증인 aArbWETH와, 소수 자릿수가 두 자리인 Fiat24 토큰 두 개가 있어요.

\begin{table}[htbp]
\centering
\caption[행이 가장 많은 토큰과, 등록한 토큰 점검의 두 상위 20개 목록.]{계획의 코호트의 에고 표에서 행이 가장 많은 토큰과, 등록 분석의 토큰 점검에 쓴 두 상위 20개 목록의 토큰이에요. 행 수로 줄 세웠어요. 행 수는 $T_{\text{obs}}$까지 에고 표에 있는 송금에 $T_{\text{obs}}$ 시점의 0보다 큰 마지막 허락을 더한 것이에요. 수정 분석은 앞의 다섯 개를 남겨요(표~\ref{tab:usd-tokens}). 금액순: 같은 행들의 원래 금액을 더한 값이 가장 큰 20개 토큰에 들어요. 행 수순: 행이 가장 많은 20개 토큰에 들어요. 기호와 이름은 2026년 10월 10일에 \texttt{symbol()}과 \texttt{name()}을 불렀을 때 토큰 계약이 돌려준 값이에요. 다만 USDT0의 계약은 투그릭 기호(몽골 돈의 기호)를 돌려주는데, 여기서는 그 자리에 T를 썼어요. 뽑아낸 자료에는 주소만 남아 있어요. 전체 주소는 \texttt{data/1-raw-blockchain-extracts/graph-tables/top\_tokens.csv}에 있어요.}
\label{tab:top-tokens}
\fitwidth{%
\begin{tabular}{lllrcc}
\toprule
기호 & 이름 & 주소 & 행 수 & 금액순 & 행 수순 \\
\midrule
WETH & Wrapped Ether & \texttt{0x82af49}\ldots\texttt{3fbab1} & 509{,}394{,}850 & \checkmark & \checkmark \\
USDC & USD Coin & \texttt{0xaf88d0}\ldots\texttt{8e5831} & 366{,}294{,}721 & \checkmark & \checkmark \\
USDT0 & USDT0 & \texttt{0xfd086b}\ldots\texttt{9fcbb9} & 163{,}462{,}893 & \checkmark & \checkmark \\
ARB & Arbitrum & \texttt{0x912ce5}\ldots\texttt{9e6548} & 75{,}552{,}592 & \checkmark & \checkmark \\
WBTC & Wrapped BTC & \texttt{0x2f2a25}\ldots\texttt{fc5b0f} & 71{,}824{,}604 & \checkmark & \checkmark \\
USDC & USD Coin (Arb1) & \texttt{0xff970a}\ldots\texttt{db5cc8} & 71{,}329{,}831 & \checkmark & \checkmark \\
PENDLE & Pendle & \texttt{0x0c880f}\ldots\texttt{a8c9e8} & 19{,}516{,}891 & \checkmark & \checkmark \\
WBTC\_ETH\_FS\_LT & WBTC ETH Futureswap Liquidity Token & \texttt{0xa522a5}\ldots\texttt{0d8d98} & 17{,}735{,}463 &  & \checkmark \\
NPC & NPC & \texttt{0xe1a498}\ldots\texttt{0bb8b0} & 12{,}163{,}764 &  & \checkmark \\
UXUY & UXUY Token & \texttt{0xe2035f}\ldots\texttt{7c79f9} & 9{,}761{,}886 & \checkmark & \checkmark \\
AIUS & Arbius & \texttt{0x4a24b1}\ldots\texttt{0cabc3} & 8{,}804{,}452 &  & \checkmark \\
DAI & Dai Stablecoin & \texttt{0xda1000}\ldots\texttt{000da1} & 8{,}677{,}480 & \checkmark & \checkmark \\
LINK & ChainLink Token & \texttt{0xf97f4d}\ldots\texttt{539fb4} & 8{,}675{,}046 &  & \checkmark \\
xUSD & xUSD & \texttt{0xe80772}\ldots\texttt{afca65} & 6{,}437{,}746 &  & \checkmark \\
ZRO & LayerZero & \texttt{0x698588}\ldots\texttt{3271cd} & 6{,}318{,}171 & \checkmark & \checkmark \\
UXLINK & UXLINK Token & \texttt{0x1a6b3a}\ldots\texttt{c8cff1} & 6{,}187{,}243 &  & \checkmark \\
VSN & Vision & \texttt{0x6fbbbd}\ldots\texttt{7db74b} & 6{,}052{,}221 & \checkmark & \checkmark \\
wstETH & Wrapped liquid staked Ether 2.0 & \texttt{0x5979d7}\ldots\texttt{800529} & 5{,}553{,}387 &  & \checkmark \\
GRT & Graph Token & \texttt{0x962306}\ldots\texttt{7e88c7} & 4{,}761{,}427 & \checkmark & \checkmark \\
GMX & GMX & \texttt{0xfc5a1a}\ldots\texttt{35ad0a} & 4{,}378{,}643 & \checkmark & \checkmark \\
BEBE & BEBE & \texttt{0x406d59}\ldots\texttt{64f97b} & 2{,}720{,}903 & \checkmark &  \\
STG & StargateToken & \texttt{0x669434}\ldots\texttt{ce1eb6} & 1{,}647{,}816 & \checkmark &  \\
aArbWETH & Aave Arbitrum WETH & \texttt{0xe50fa9}\ldots\texttt{8128c8} & 1{,}343{,}622 & \checkmark &  \\
OTMT & OTMAirdrop & \texttt{0x6b708d}\ldots\texttt{37f770} & 493{,}833 & \checkmark &  \\
BPET & BPET & \texttt{0x6daf58}\ldots\texttt{862ac5} & 389{,}567 & \checkmark &  \\
USD24 & Fiat24 USD & \texttt{0xbe00f3}\ldots\texttt{cc0d62} & 207{,}237 & \checkmark &  \\
CNH24 & Fiat24 CNH & \texttt{0x7288ac}\ldots\texttt{0144fd} & 102{,}957 & \checkmark &  \\
\bottomrule
\end{tabular}
}
\end{table}

미국 달러로 따지면, 관측 기간에 남긴 송금의 중앙값은 $193.57$달러예요. 10분의 1은 $0.14$달러보다 적고, 4분의 1은 $12.19$달러보다 적어요. 4분의 3은 $891.94$달러보다 적어요. 그리고 90번째 백분위(작은 값부터 줄 세웠을 때 그 비율의 자리에 있는 값)와 99번째 백분위는 $2{,}967.6$달러와 $19{,}999.9$달러예요. $T_{\text{obs}}$ 시점의 0보다 크고 유한한(한도가 있는) 마지막 허락은 2026년 6월 30일의 가격으로 값을 매겨요. 그 중앙값은 $60.58$달러예요. 10분의 1은 $0.33$달러보다 적고, 4분의 1은 $5.44$달러보다 적고, 4분의 3은 $1{,}198.4$달러보다 적어요. 90번째 백분위는 약 $2.1\times10^{9}$달러이고, 99번째 백분위는 약 $5.9\times10^{45}$달러예요. $2^{255}$나 $2^{128}$처럼 $2^{256}-1$ 원래 단위보다 작은 허락 가운데 많은 것이 사실상 무제한이기 때문이에요. \ref{sec:amount-revision}절의 상수들은 $t_1$까지의 같은 값들에서 나와요. 송금의 중앙값 $203.71$달러와 유한한 허락의 중앙값 $61.51$달러에서 $m_T=200$달러와 $m_A=62$달러가 나와요. 그리고 송금 하나하나의 99번째 백분위인 $20{,}266.4$달러에서 상한(이보다 큰 금액은 이 값으로 자름) $U=20{,}000$달러가 나와요. $T_{\text{obs}}$에서는 두 중앙값이 $m_T$와 $m_A$보다 조금 아래에 있어요. 그래서 송금과 유한한 허락 가운데 절반이 넘는 것이, 절반보다 작은 금액 무게로 AWP와 엔도스랭크에 들어가요. 그리고 두 금액 모두, 아래 사분위수(아래에서 4분의 1 자리)와 위 사분위수(위에서 4분의 1 자리) 사이에서 무게가 0에 가까운 값에서 1에 가까운 값까지 올라가요.

무제한 허락은 흔해요. $T_{\text{obs}}$ 시점의 0보다 큰 마지막 허락 $12{,}590{,}613$개 가운데 $5{,}763{,}757$개($45.8\%$)는 원래 단위로 $1.15\times10^{77}$ 이상이에요. 이것은 \texttt{uint256}의 가장 큰 값인 $2^{256}-1\approx1.16\times10^{77}$에 가까워요. 그리고 무제한 허락을 포함해 $6{,}776{,}289$개($53.8\%$)는 값이 상한 $U=20{,}000$달러 이상이에요. 이 허락 $12{,}590{,}613$개는 주인 $6{,}208{,}769$명의 것이에요. 그 가운데 $3{,}294{,}370$명은 무제한 허락을 적어도 하나 해 두었어요. 그리고 $797{,}852$명은 무제한 허락과 유한한 허락을 둘 다 해 두었어요(\texttt{sql/18\_usd\_describe.sql}). 엔도스랭크는 $m_A=62$달러짜리 허락에 금액 무게 절반을 줘요. $124$달러짜리 허락에는 $0.8$을 주고, 무제한 허락에는 1을 줘요. 1은 허락 하나가 가질 수 있는 가장 큰 무게예요. C-PR(결합 페이지랭크, 두 그래프를 함께 걷는 점수)의 허락 층에서는 $U$ 이상인 허락을 모두 $U$로 세요. 그래서 두 종류의 허락을 다 해 둔 주인의 행(이동 행렬의 한 줄)에서는, 무제한 허락이 $20{,}000$달러짜리 유한한 허락과 같은 무게를 가져요. 그리고 더 작은 허락들은 금액에 비례해서 그 행의 몫을 나눠 가져요(\ref{sub:coupled-graph}절). 등록 분석은 모든 토큰을 원래 단위 그대로 두고 기울기 $b=1$을 썼어요. 그 분석에서는 에고 송금의 $99.52\%$와, $T_{\text{obs}}$ 시점의 마지막 허락 $28{,}069{,}418$개 가운데 $99.78\%$가 원래 단위로 10 이상이었어요. 그 범위에서는 $V$가 1과 $10^{-4}$ 안쪽으로만 차이 나요. 그래서 엔도스랭크와 AWP에서 금액은 아무 역할도 하지 않았어요. 그 분석에서는 0보다 큰 마지막 허락을 해 둔 주인 $13{,}029{,}306$명 가운데 $1{,}532{,}301$명이 무제한 허락과 그보다 작은 허락을 둘 다 해 두었어요. 그리고 원래 금액을 쓰는 C-PR의 허락 층에서는, 무제한 허락이 그런 주인의 행을 거의 모두 차지했어요(부록~\ref{sec:map-registered}).

\dissertationSubheading[sec:data-gmx]{분기별 GMX~V2 계약 트레이더}

트레이더(거래하는 사람) 코호트는 계약 코드를 가진 GMX~V2 계정들이에요(\ref{sub:cohorts}절). 표~\ref{tab:gmx-quarters}은 이 계정들이 포지션을 닫은 횟수를 달력 분기마다 세요. 포지션 닫기는 이 서비스(프로토콜)의 이벤트 발신기(이벤트를 내보내는 계약)가 남기는 \texttt{PositionDecrease} 이벤트예요\parencite{gmxnddocs}. 표는 그 가운데 청산(손실이 너무 커져 거래가 강제로 끝나는 일)의 수도 함께 세요. 또 그 분기에 포지션을 적어도 하나 닫은 계약 계정 수와, 적어도 세 개 닫은 계약 계정 수도 세요.

\begin{table}[htbp]
\centering
\caption[계약 계정의 분기별 GMX~V2 포지션 닫기.]{계약 계정의 분기별 GMX~V2 포지션 닫기예요. 2023년 10월부터 2026년 6월까지: 관측 기간에 포지션을 적어도 세 번 닫은 GMX~V2 계정 $40{,}307$개 가운데 계약 $7{,}633$개가 닫은 것이에요(표~\ref{tab:account-types}). 2026년 7월부터 9월까지(W1의 라벨 기간): 관측 기간의 계약 계정과 W1 실행의 계약 계정이 닫은 것이에요. 청산: 주문 종류가 청산이거나, 청산 처리기를 거쳐 나온 닫기예요. 출처: \texttt{data/2-processed-tables-and-evaluations/gmx/}에 있는 해독한 닫기 기록이에요.}
\label{tab:gmx-quarters}
\fitwidth{%
\begin{tabular}{lrrrr}
\toprule
분기 & 닫기 & 청산 & 닫기가 있는 계정 & 닫기가 $\geq 3$번인 계정 \\
\midrule
2023년 4분기 & 1{,}119 & 104 & 180 & 105 \\
2024년 1분기 & 7{,}956 & 692 & 1{,}247 & 961 \\
2024년 2분기 & 11{,}509 & 589 & 2{,}120 & 1{,}165 \\
2024년 3분기 & 7{,}947 & 456 & 1{,}807 & 724 \\
2024년 4분기 & 7{,}343 & 904 & 1{,}321 & 698 \\
2025년 1분기 & 6{,}929 & 498 & 497 & 273 \\
2025년 2분기 & 12{,}080 & 379 & 368 & 250 \\
2025년 3분기 & 28{,}753 & 897 & 1{,}750 & 1{,}567 \\
2025년 4분기 & 40{,}240 & 1{,}596 & 1{,}915 & 1{,}744 \\
2026년 1분기 & 6{,}026 & 233 & 269 & 171 \\
2026년 2분기 & 3{,}783 & 125 & 169 & 127 \\
\midrule
2026년 3분기 (W1) & 6{,}805 & 182 & 165 & 131 \\
\midrule
2023년 10월 -- 2026년 6월 & 133{,}685 & 6{,}473 & & \\
\bottomrule
\end{tabular}
}
\end{table}

계약 계정은 2025년 하반기에 포지션을 가장 많이 닫았어요. 3분기에 $28{,}753$번, 4분기에 $40{,}240$번 닫았어요. 그때 포지션을 적어도 세 번 닫은 계정은 $1{,}567$개와 $1{,}744$개였어요. 2026년 상반기에는 닫기가 $6{,}026$번과 $3{,}783$번으로 줄었어요. 적어도 세 번 닫은 계정도 $171$개와 $127$개로 줄었어요. 트레이더 코호트에 들려면, 계정이 동결일에 그래프의 점이기도 해야 해요. 2026년 4월부터 6월까지 적어도 세 번 닫은 계약 계정 $127$개 가운데 $88$개가 W0의 트레이더 코호트를 이뤄요. 그리고 7월부터 9월까지의 $131$개 가운데 $70$개가 W1의 트레이더 코호트를 이뤄요. \ref{ch:results}장의 모든 트레이더 결과가 넓은 구간(믿을 만한 범위)을 가지는 것은 코호트가 이렇게 작기 때문이에요.

% !TEX root = ../main.tex
\section{엔도스랭크와 AWP를 나란히 놓고 보기}
\label{ch:appendix-er-awp}

\ref{ch:results}장은 이 논문의 모든 점수를 함께 보고해요. 그리고 각 페이지랭크(PageRank)를, 그 페이지랭크가 매끄럽게 다듬는(평활하는) 원시 차수(화살표를 그냥 센 수)와 견주어 읽어요. 이 부록은 한 층만 걷는 두 점수인 엔도스랭크(EndorseRank)와 AWP(적응형 가중 페이지랭크)를 바로 견준 결과를 모아요. 같은 기간에서의 지표 가족(비슷한 지표 묶음) 평균과 짝지은 차이, 대리 지표(대신 재는 값) 하나하나에 대한 계수, 기간 밖에서의 짝지은 차이, 그리고 모든 트레이더(거래하는 사람) 라벨(나중에 채점에 쓰는 정답 값)에서의 대비예요. 여기 있는 모든 수는 수정 분석(등록 뒤에 금액을 달러로 바꿔 다시 한 분석)에서 나와요. 금액은 다섯 토큰에 걸쳐 달러(USD)로 매겼어요(\ref{sec:amount-revision}절). 등록 분석(미리 등록한 그대로 한 분석)은 부록~\ref{sec:map-registered}에 있어요. 매칭 코호트(모든 점수가 똑같이 점수를 매기는 묶음)나 스펜더(spender, 허락을 받아 대신 쓰는 쪽) 코호트에서 한 비교는 어느 것도 미리 정한 것이 아니에요. 그리고 어느 것도 이 논문의 주장(C)을 담고 있지 않아요. 트레이더 쌍은 등록 분석을 위해 W1과 함께 등록했어요. 여기서는 그 쌍을 고친 점수로 다시 계산해요(\ref{sec:er-awp-neutral-grid}절). 두 점수는 서로 다른 화살표(엣지) 위를 걷고, 서로 다른 구성개념(점수가 재려고 하는 생각)을 위해 만들었어요. 그래서 코호트 전체에서 하는 점검마다, 앞서는 점수는 그 점검이 세는 화살표를 걷는 점수예요. 두 화살표 종류 어느 쪽으로도 만들지 않은 결과는 예외예요. 어느 점수도 그 결과를 위해 만들지 않았기 때문이에요. 그 결과는 \ref{sub:neutral-label-results}절에서 보고해요.

두 점수가 점수를 매기는 범위는 서로 달라요. 매칭 코호트 가운데 스마트 계약(블록체인 위에서 저절로 실행되는 프로그램) $523$개는 허락(allowance, 토큰을 얼마까지 대신 써도 된다는 허락) 그래프(점과 화살표로 이은 그림) 밖에 있고, $2{,}670$개는 송금(transfer) 그래프 밖에 있어요. 또 W0의 스펜더 $17{,}050$개 가운데 $2{,}891$개는 동결(점수를 고정하는 일) 전에 송금 화살표가 하나도 없었어요. 그래서 AWP는 이 계약들에 0점을 줘요. 이와 달리 엔도스랭크는 코호트 계약 쌍을 거의 하나도 동점으로 남기지 않아요(\ref{sec:rank-divergence}절). 기간 밖에서는 두 점수가 자기가 다듬는 원시 차수와 견주어 서로 다른 자리에 서요. AWP는 두 기간 모두에서 나중의 새 송금자(새로 송금을 보내온 주소)를 송금 입차수(송금을 보내온 서로 다른 주소 수)보다 더 잘 줄 세워요. 반면에 엔도스랭크는 두 기간 모두에서 나중의 새 허락 쌍(새로 생긴 주인–스펜더 허락)을 받은 허락 수(허락 입차수)보다 못하게 줄 세워요(\ref{sec:holdout-results}절). 그러므로 엔도스랭크가 AWP보다 앞선다고 해서, 허락을 따라 퍼뜨리는 것(전파)이 허락한 주인(허락을 해 주는 쪽)의 수에 무언가를 보탠다는 것을 보여 주지는 않아요.

\dissertationSubheading[sec:er-awp-same-window]{같은 기간}

표~\ref{tab:alignment-hybrid}은 두 점수의 지표 가족 평균을 구간(믿을 만한 범위)과 함께 보여 줘요. 표~\ref{tab:tie-sensitivity}의 마지막 행은 두 점수 사이의 일치를 보여 줘요. 두 모양은 서로 엇갈려요. AWP는 송금 가족을 가장 가깝게 따라가고($0.572$, 구간 $[0.564,0.580]$), 허락 가족을 가장 덜 따라가요($0.193$, $[0.182,0.203]$). 반면에 엔도스랭크는 허락 가족을 가장 가깝게 따라가고($0.585$, $[0.578,0.592]$), 송금 가족에서는 $0.158$($[0.147,0.169]$)에 머물러요. 표~\ref{tab:tau-diff}은 계약을 다시 뽑은 같은 표본에서 짝지은 차이 다섯 개를 검사해요(\ref{sec:statistical-tests}절). 이 차이들에는 표의 행 순서대로 i부터 v까지 번호를 붙였어요.

% Generated by experiments/scripts/export_latex.py -- do not edit by hand.
% Source: experiments/data/2-processed-tables-and-evaluations/revised-usd/same-window/eval_summary.json
\begin{table}[htbp]
\centering
\caption[같은 기간에서 엔도스랭크와 AWP 사이의 짝지은 대비.]{같은 기간에서 엔도스랭크와 AWP의 켄달의 $\tau$ 계수 사이 차이를 짝지은 부트스트랩으로 검사한 결과예요(매칭 코호트, 계약 $n=15{,}107$개). 지표 가족 항목에는 지표 가족 평균 $\tau$를 써요. $a$와 $b$에 같은 계약 다시 뽑기를 써요. 그래서 구간은 짝지은 차이에 대한 구간이에요. 구간이 0을 벗어나면, 5\% 수준에서 표본을 뽑을 때 생기는 흔들림만으로는 설명할 수 없는 차이라는 뜻이에요. 신뢰구간(CI)은 계약을 다시 뽑아 구한 95\% 백분위 부트스트랩 구간이에요(짝지은 다시 뽑기 400번, 시드 42).}
\label{tab:tau-diff}
\fitwidth{%
\begin{tabular}{lccccc}
\toprule
대비 ($a$ 빼기 $b$) & $\tau_a$ & $\tau_b$ & $\Delta\tau$ & $\Delta\tau$의 95\% 신뢰구간 & $P(\Delta\tau>0)$ \\
\midrule
EndorseRank 허락 빼기 EndorseRank 송금 & 0.585 & 0.158 & +0.427 & [0.416, 0.440] & 1.000 \\
AWP 송금 빼기 EndorseRank 송금 & 0.572 & 0.158 & +0.414 & [0.401, 0.427] & 1.000 \\
EndorseRank 허락 빼기 AWP 허락 & 0.585 & 0.193 & +0.392 & [0.380, 0.405] & 1.000 \\
AWP 시빌 안정성 빼기 EndorseRank 시빌 안정성 & 0.395 & 0.007 & +0.388 & [0.376, 0.401] & 1.000 \\
EndorseRank 받은 허락 수 빼기 EndorseRank 입차수 & 0.655 & 0.182 & +0.473 & [0.461, 0.486] & 1.000 \\
\bottomrule
\end{tabular}
}
\end{table}


다섯 구간은 모두 0보다 위에 있어요. 엔도스랭크의 허락 계수는 자기의 송금 계수보다 커요(대비~i, $+0.427$, 구간 $[0.416,0.440]$). AWP의 송금 계수는 엔도스랭크의 송금 계수보다 커요(대비~ii, $+0.414$, $[0.401,0.427]$). 엔도스랭크의 허락 계수는 AWP의 허락 계수보다 커요(대비~iii, $+0.392$, $[0.380,0.405]$). AWP는 시빌 안정성 가족(시빌 공격, 곧 가짜 주소를 많이 만들어 속이는 공격을 생각해 고른 지표 묶음)에서 앞서요(대비~iv, $+0.388$, $[0.376,0.401]$). 그리고 엔도스랭크는 들어오는 화살표 수(입차수)보다 받은 허락 수를 더 가깝게 따라가요(대비~v, $+0.473$, $[0.461,0.486]$). 대비~i와~v는 대리 지표만 바꾸어 엔도스랭크를 자기 자신과 견줘요. 엔도스랭크는 계약 쌍의 $0.17\%$만 동점으로 둬요. 그래서 엔도스랭크 자신의 동점은 이 계수들 어느 것에도 천장(얻을 수 있는 가장 큰 값)을 씌우지 않아요. 그리고 대리 지표의 동점으로는 이만큼 큰 차이를 설명할 수 없어요. 입차수에서 AWP는 $0.578$($[0.569,0.587]$)에 이르지만, 엔도스랭크는 $0.182$($[0.169,0.194]$)에 이르러요. 허락 그래프 밖의 계약들을 외톨이 점(화살표가 하나도 없는 점)으로 남겨 두어도, 지표 가족 평균은 하나도 바뀌지 않아요(표~\ref{tab:tie-sensitivity}).

대비~ii와~iii의 부호가 서로 반대인 것은 구성개념에서 나와요. 각 점수는 자기 화살표로 만든 가족에서 앞서요. 그리고 이 가족들은 저마다, 그에 맞는 페이지랭크가 다듬는 입차수를 담고 있어요. 이 일치는 일부가 만든 방식 때문에 저절로 생기고, 각 점수의 의도한 구성개념을 점검해요. 이것은 점수 바깥에서 온 증거가 아니에요. 코호트의 모든 계약은 적어도 세 주인에게서 WETH, USDC, USDT0, ARB 또는 WBTC의 허락을 받았어요. 그리고 AWP의 허락 계수는 0보다 위에 있어요(표~\ref{tab:alignment-allowance}). 곧, 많은 주인이 허락한 계약은 송금 그래프에서도 중심에 있는 경향이 있어요. 어느 차이도 어느 점수가 평판을 더 잘 재는지는 말해 주지 않아요.

표~\ref{tab:robustness-sample-size}의 규모 키우기 단계들의 값은 코호트 전체의 값과 견줄 수 없어요. 이 단계들에서도 각 가족 안의 두 점수 순서는 모든 단계에서 같아요. AWP는 송금 가족과 시빌 안정성 가족에서 앞서고, 엔도스랭크는 허락 가족에서 앞서요.

\dissertationSubheading[sec:proxy-family-detail]{대리 지표 하나하나 보기}

지표 가족 평균 하나는 대리 지표마다 구한 계수 여러 개를 한데 줄인 값이에요. 표~\ref{tab:alignment-transfer}부터~\ref{tab:alignment-sybil-stability}까지는 엔도스랭크와 AWP의 계수를 대리 지표 하나씩 보고해요. 스피어먼의 로 $\rho$(두 줄 세우기가 얼마나 비슷한지 나타내는 수), 켄달의 타우 $\tau$(두 줄 세우기가 얼마나 같은 순서인지 나타내는 수), 그리고 $\tau$의 95\% 구간이에요. 모든 대리 지표는 다섯 토큰의 이벤트(계약이 남기는 기록)를 세고, 금액은 달러로 매겨요(\ref{sec:proxy-formulas}절). 상관을 구하기 전에, 모든 대리 지표는 값이 클수록 기대되는 평판이 높다는 뜻이 되도록 방향을 맞춰요. 그리고 송금 대리 지표는 AWP를 만든 사람들이 AWP를 평가할 때 쓴 바로 그 지표예요\parencite{do2023awp}.

\dissertationSubsubheading[sub:transfer-family]{송금 가족}

표~\ref{tab:alignment-transfer}은 두 점수와 입차수, 그리고 받은 양의 합 사이의 계수를 보여 줘요.

% Generated by experiments/scripts/export_latex.py -- do not edit by hand.
% Source: experiments/data/2-processed-tables-and-evaluations/revised-usd/same-window/eval_summary.json
\begin{table}[htbp]
\centering
\caption[송금 대리 지표에 대한 엔도스랭크와 AWP의 분야 일치도.]{엔도스랭크와 AWP가 송금 대리 지표(각 계약의 입차수와 받은 양의 합)와 얼마나 잘 맞는지 나타낸 분야 일치도예요. 같은 기간에서 쟀어요(매칭 코호트, 계약 $n=15{,}107$개). 점수의 순위와 대리 지표의 순위 사이의 켄달의 $\tau$와 스피어먼의 $\rho$를 보여 줘요. 신뢰구간(CI)은 계약을 다시 뽑아 구한 95\% 백분위 부트스트랩 구간이에요(짝지은 다시 뽑기 400번, 시드 42).}
\label{tab:alignment-transfer}
\fitwidth{%
\begin{tabular}{lccc}
\toprule
대리 지표 & 스피어먼의 $\rho$ & 켄달의 $\tau$ & $\tau$의 95\% 신뢰구간 \\
\midrule
\multicolumn{4}{l}{\textit{EndorseRank}} \\
입차수 & 0.261 & 0.182 & [0.169, 0.194] \\
받은 양의 합 (들어온 양을 모두 더함) & 0.205 & 0.134 & [0.123, 0.145] \\
\addlinespace
\multicolumn{4}{l}{\textit{AWP}} \\
입차수 & 0.744 & 0.578 & [0.569, 0.587] \\
받은 양의 합 (들어온 양을 모두 더함) & 0.741 & 0.566 & [0.558, 0.575] \\
\bottomrule
\end{tabular}
}
\end{table}


AWP는 입차수와 $0.578$(구간 $[0.569,0.587]$)만큼 맞고, 받은 양의 합과 $0.566$($[0.558,0.575]$)만큼 맞아요. AWP의 금액 무게는 중앙값 송금인 200달러에서 2분의 1이에요. 그리고 그보다 몇 배 큰 금액에서만 1에 가까워져요(\ref{sub:awp-graph}절). 그래서 큰 송금일수록 더 크게 세요. 그리고 AWP는 한 계약이 받은 송금의 달러 합계를, 그 송금의 수와 거의 같은 만큼 가깝게 따라가요. 엔도스랭크의 송금 계수는 $0.182$($[0.169,0.194]$)와 $0.134$($[0.123,0.145]$)로 0보다 위에 있어요. 하지만 그 구간은 AWP의 구간보다 훨씬 아래에 있어요. 곧, 많은 주인이 허락한 계약은 코호트의 나머지 계약보다 더 많은 보낸 쪽(송신자)에게서 송금을 받아요. 기간 밖에서 AWP는 두 기간 모두에서 나중의 새 송금자를 원시 입차수보다 더 잘 줄 세워요(\ref{sec:holdout-results}절). 그러니 송금 쪽에서는 걷기가 자기가 다듬는 수에 무언가를 보태요.

\dissertationSubsubheading[sub:allowance-family]{허락 가족}

표~\ref{tab:alignment-allowance}는 두 점수와 받은 허락 수, 그리고 받은 허락 금액 사이의 계수를 보여 줘요.

% Generated by experiments/scripts/export_latex.py -- do not edit by hand.
% Source: experiments/data/2-processed-tables-and-evaluations/revised-usd/same-window/eval_summary.json
\begin{table}[htbp]
\centering
\caption[허락 대리 지표에 대한 엔도스랭크와 AWP의 분야 일치도.]{엔도스랭크와 AWP가 허락 대리 지표(각 스펜더 계약의 받은 허락 수와 받은 허락의 양)와 얼마나 잘 맞는지 나타낸 분야 일치도예요. 같은 기간에서 쟀어요(매칭 코호트, 계약 $n=15{,}107$개). 점수의 순위와 대리 지표의 순위 사이의 켄달의 $\tau$와 스피어먼의 $\rho$를 보여 줘요. 신뢰구간(CI)은 계약을 다시 뽑아 구한 95\% 백분위 부트스트랩 구간이에요(짝지은 다시 뽑기 400번, 시드 42).}
\label{tab:alignment-allowance}
\fitwidth{%
\begin{tabular}{lccc}
\toprule
대리 지표 & 스피어먼의 $\rho$ & 켄달의 $\tau$ & $\tau$의 95\% 신뢰구간 \\
\midrule
\multicolumn{4}{l}{\textit{EndorseRank}} \\
받은 허락 수 & 0.811 & 0.655 & [0.648, 0.662] \\
받은 허락의 양 (합) & 0.681 & 0.515 & [0.506, 0.524] \\
\addlinespace
\multicolumn{4}{l}{\textit{AWP}} \\
받은 허락 수 & 0.286 & 0.206 & [0.193, 0.217] \\
받은 허락의 양 (합) & 0.260 & 0.180 & [0.169, 0.191] \\
\bottomrule
\end{tabular}
}
\end{table}


엔도스랭크는 받은 허락 수와 $0.655$(구간 $[0.648,0.662]$)만큼 맞아요. 그리고 받은 허락 금액과는 $0.515$($[0.506,0.524]$)만큼 맞아요. 받은 허락 금액은 상한(이보다 큰 금액은 이 값으로 자름)으로 자른 허락들을 더한 값이에요. 금액 무게는 끝이 있는 금액의 중앙값인 62달러짜리 허락에 2분의 1을 줘요. 하지만 0보다 큰 마지막 허락(가장 최근 허락) 가운데 절반이 넘는 것이 무제한 허락이거나 20{,}000달러에 이르러요(표~\ref{tab:usd-tokens}). 이런 허락은 꽉 채워 세요. 그래서 엔도스랭크는 주인들이 허락한 금액의 합보다, 허락한 주인의 수를 더 가깝게 따라가요. 주인들이 허락을 하나도 받지 않은 스펜더에서는, 엔도스랭크가 \ref{sub:theory-closed-form}절의 닫힌 식(바로 계산할 수 있는 식)과 같아요. 이 식은 주인 정규화 입강도(각 주인이 준 무게를 그 주인이 허락한 모든 스펜더에게 나누어 더한 값)에 따라 정해져요. AWP는 두 대리 지표와 $0.206$($[0.193,0.217]$), $0.180$($[0.169,0.191]$)만큼 맞아요. 이 비교는 일부가 만든 방식 때문에 저절로 그렇게 돼요. 엔도스랭크가 대리 지표와 같은 허락들로 계산되기 때문이에요.

\dissertationSubsubheading[sub:sybil-family]{시빌 보정 안정성 가족}

표~\ref{tab:alignment-sybil-stability}은 두 점수와 들어오는 상대 비율, 송금 활동 기간, 활동한 달 수 사이의 계수를 보여 줘요.

% Generated by experiments/scripts/export_latex.py -- do not edit by hand.
% Source: experiments/data/2-processed-tables-and-evaluations/revised-usd/same-window/eval_summary.json
\begin{table}[htbp]
\centering
\caption[시빌 보정 안정성 대리 지표에 대한 엔도스랭크와 AWP의 분야 일치도.]{엔도스랭크와 AWP가 송금 활동에서 구한 시빌 보정 안정성 대리 지표와 얼마나 잘 맞는지 나타낸 분야 일치도예요. 같은 기간에서 쟀어요(매칭 코호트, 계약 $n=15{,}107$개). 점수의 순위와 대리 지표의 순위 사이의 켄달의 $\tau$와 스피어먼의 $\rho$를 보여 줘요. 신뢰구간(CI)은 계약을 다시 뽑아 구한 95\% 백분위 부트스트랩 구간이에요(짝지은 다시 뽑기 400번, 시드 42).}
\label{tab:alignment-sybil-stability}
\fitwidth{%
\begin{tabular}{lccc}
\toprule
대리 지표 & 스피어먼의 $\rho$ & 켄달의 $\tau$ & $\tau$의 95\% 신뢰구간 \\
\midrule
\multicolumn{4}{l}{\textit{EndorseRank}} \\
들어오는 상대 비율 & -0.231 & -0.156 & [-0.166, -0.147] \\
송금 활동 기간 (일) & 0.124 & 0.080 & [0.069, 0.091] \\
활동한 달 수 & 0.146 & 0.098 & [0.087, 0.109] \\
\addlinespace
\multicolumn{4}{l}{\textit{AWP}} \\
들어오는 상대 비율 & 0.234 & 0.127 & [0.114, 0.141] \\
송금 활동 기간 (일) & 0.683 & 0.508 & [0.499, 0.517] \\
활동한 달 수 & 0.704 & 0.549 & [0.540, 0.558] \\
\bottomrule
\end{tabular}
}
\end{table}


표를 보면 AWP는 송금 활동 기간에서 엔도스랭크보다 앞서요($0.508$, 구간 $[0.499,0.517]$, 대 $0.080$, $[0.069,0.091]$). 활동한 달 수에서도 앞서요($0.549$, $[0.540,0.558]$, 대 $0.098$, $[0.087,0.109]$). 기간 전체에 퍼진 활동에 점수를 더 주는 송금 점수라면 그럴 것이라고 예상할 수 있어요. 지표 가족 평균에서의 짝지은 차이는 0을 벗어나요(표~\ref{tab:tau-diff}, 대비~iv). 들어오는 상대 비율에서는 두 계수의 부호가 서로 반대예요. 엔도스랭크는 $-0.156$($[-0.166,-0.147]$)이고, AWP는 $0.127$($[0.114,0.141]$)이에요. 이 비율은 자기 자신을 뺀 서로 다른 보낸 쪽의 수를, 그들에게서 받은 송금의 수로 나눈 값이에요. 그래서 같은 보낸 쪽이 다시 오면 비율이 내려가요. 엔도스랭크의 계수가 음수라는 것은, 많은 주인이 허락한 계약이 이 잣대로 보면 보낸 쪽이 다시 찾아오는 계약이라는 뜻이에요. AWP의 계수가 양수인 것은 점수를 매기는 범위 때문이에요. 송금 그래프 밖에 있는 계약 $2{,}670$개는 비율도 0이고 AWP 점수도 0이에요. 그래서 이 계약들이 두 줄 세우기의 맨 아래를 함께 채워요. AWP가 0보다 큰 점수를 준 계약들만 보면, AWP의 순서도 엔도스랭크처럼 비율이 커질수록 내려가요. 이 점검은 결과를 안 뒤에 했고, 구간은 없어요. 이 가족은 각 점수가 시간에 걸친 송금 정보를 얼마나 담고 있는지를 보여 줘요.

\dissertationSubheading[sec:er-awp-holdout]{기간 밖}

표~\ref{tab:endorserank-awp-holdout}은 두 홀드아웃(나중 기간을 떼어 두고 맞혀 보는 시험) 기간에서 두 점수 사이의 짝지은 차이를 모아요. 모든 행에서 $a$는 엔도스랭크이고 $b$는 AWP예요.

% Generated by experiments/scripts/export_latex.py -- do not edit by hand.
% Source: experiments/data/2-processed-tables-and-evaluations/revised-usd/holdout-w0/eval_summary.json
% Source: experiments/data/2-processed-tables-and-evaluations/revised-usd/supplementary/supplementary.json
% Source: experiments/data/2-processed-tables-and-evaluations/revised-usd/window-w1/eval_summary.json
\begin{table}[htbp]
\centering
\caption[기간 밖에서 엔도스랭크와 AWP 사이의 짝지은 부트스트랩 차이.]{기간 밖에서 엔도스랭크와 AWP 사이의 차이를 짝지은 부트스트랩으로 구했어요(W0 스펜더 $n=17{,}050$, W0 트레이더 $n=88$, W1 스펜더 $n=18{,}152$, W1 트레이더 $n=70$). 이 가운데 어느 것도 규칙 A나 규칙 B에 들어가지 않아요. 스펜더 행은 라벨을 본 뒤에 계산했어요. 트레이더 행은 어느 화살표 종류로도 만들지 않은 결과에서 한 비교 가운데, 엔도스랭크에서 AWP를 뺀 쌍이에요(표~\ref{tab:neutral-label-contrasts}과~\ref{tab:neutral-label-grid}). W1 스펜더 행은 W1 평가가 끝난 뒤에, 저장해 둔 W1 점수로 계산했어요. 다시 뽑은 표본은 표~\ref{tab:fresh-holdout}와 같은 것을 썼어요. 각 묶음 안에서 $a$와 $b$는 같은 계약 다시 뽑은 표본을 함께 써요(스펜더는 400번, 트레이더는 2{,}000번, 시드 42).}
\label{tab:endorserank-awp-holdout}
\fitwidth{%
\begin{tabular}{llccccc}
\toprule
대비 ($a$ 빼기 $b$) & 라벨 & $\tau_a$ & $\tau_b$ & $\Delta\tau$ & $\Delta\tau$의 95\% 신뢰구간 & $P(\Delta\tau>0)$ \\
\midrule
\multicolumn{7}{l}{\emph{W0, 스펜더 ($n=17{,}050$)}} \\
EndorseRank 빼기 AWP & 새 허락 쌍 & 0.231 & 0.196 & +0.035 & [0.019, 0.051] & 1.000 \\
EndorseRank 빼기 AWP & 새 송금자 & 0.200 & 0.346 & -0.146 & [-0.156, -0.136] & 0.000 \\
\midrule
\multicolumn{7}{l}{\emph{W0, 세 번 이상 닫은 트레이더 ($n=88$)}} \\
EndorseRank 빼기 AWP & 청산 없이 닫은 비율 & 0.224 & 0.240 & -0.017 & [-0.156, 0.128] & 0.394 \\
\midrule
\multicolumn{7}{l}{\emph{W1, 스펜더 ($n=18{,}152$)}} \\
EndorseRank 빼기 AWP & 새 허락 쌍 & 0.234 & 0.169 & +0.065 & [0.051, 0.081] & 1.000 \\
EndorseRank 빼기 AWP & 새 송금자 & 0.194 & 0.330 & -0.136 & [-0.146, -0.127] & 0.000 \\
\midrule
\multicolumn{7}{l}{\emph{W1, 세 번 이상 닫은 트레이더 ($n=70$)}} \\
EndorseRank 빼기 AWP & 청산 없이 닫은 비율 & 0.297 & 0.329 & -0.032 & [-0.133, 0.066] & 0.279 \\
\bottomrule
\end{tabular}
}
\end{table}


W0의 스펜더에서 엔도스랭크는 새 허락 쌍을 $\tau=0.231$(구간 $[0.220,0.242]$)로 줄 세워요. AWP는 $0.196$($[0.181,0.210]$)이에요. 짝지은 차이는 $+0.035$($[0.019,0.051]$)예요. 새 송금자에서는 순서가 뒤집혀요. AWP는 $0.346$($[0.336,0.356]$)에 이르고, 엔도스랭크는 $0.200$($[0.189,0.210]$)이에요. 차이는 $-0.146$($[-0.156,-0.136]$)이에요. 모든 라벨은 점수에 쓴 모든 이벤트보다 뒤에 와요. 그러니 이 코호트에서는 각 점수가 자기 화살표 종류로 나중에 들어온 것을 다른 점수보다 더 잘 줄 세워요. 첫 번째 차이의 일부는 점수를 매기는 범위 때문이에요. AWP는 동결 전에 송금 화살표가 없던 스펜더 $2{,}891$개에 0점을 줘요. 그런데 이 스펜더의 $13.3\%$가 새 허락 쌍을 얻었어요. 나머지 스펜더는 $11.9\%$였어요. 다른 주소에게서 송금을 받은 적이 있는 스펜더 $14{,}094$개에서는, 새 허락 쌍에서의 순서가 뒤집혀요. AWP는 $0.298$($[0.287,0.308]$)에 이르고, 엔도스랭크는 $0.236$($[0.224,0.249]$)이에요. 여기에는 짝지은 검사를 하지 않았어요(표~\ref{tab:holdout-receiving}). 이 스펜더들에서도 새 송금자로는 AWP가 계속 앞서요. $0.357$($[0.347,0.367]$) 대 $0.231$($[0.219,0.244]$)이에요.

7월부터 9월까지의 기간에서도, 스펜더 코호트 전체에서 같은 모양이 되풀이돼요. 이때 대비는 모든 라벨을 본 뒤에, 수정한 W1 점수로 계산했어요. 엔도스랭크는 새 허락 쌍을 $0.234$($[0.222,0.244]$)로 줄 세우고, AWP는 $0.169$($[0.157,0.183]$)예요($+0.065$, $[0.051,0.081]$). AWP는 새 송금자를 $0.330$($[0.321,0.340]$)으로 줄 세우고, 엔도스랭크는 $0.194$($[0.184,0.205]$)예요($-0.136$, $[-0.146,-0.127]$). 트레이더의 청산 없이 닫은 비율(강제로 끝나지 않고 거래를 닫은 비율)도 보아요. 이 쌍은 W1의 등록 문서가 등록 분석을 위해 정해 둔 쌍이에요. 엔도스랭크는 W0에서 $0.224$($[0.105,0.339]$), W1에서 $0.297$($[0.171,0.410]$)에 이르러요. AWP는 $0.240$($[0.073,0.380]$)과 $0.329$($[0.207,0.440]$)예요. 두 차이 모두 0과 가려낼 수 없어요($-0.017$, $[-0.156,0.128]$, 그리고 $-0.032$, $[-0.133,0.066]$). 트레이더 코호트에는 계약이 각각 $88$개와 $70$개 있어서, 구간이 넓어요.

W0 트레이더의 이익 라벨(표~\ref{tab:holdout-traders})에서, 이익을 낸 닫기와 실현 이익(실제로 얻은 이익)에 대한 엔도스랭크의 구간은 0을 품어요. 값은 $0.150$($[-0.019,0.295]$)과 $0.078$($[-0.094,0.226]$)이에요. 반면에 AWP의 구간은 0보다 위에 있어요. 값은 $0.155$($[0.005,0.297]$)와 $0.196$($[0.048,0.351]$)이에요. 두 라벨은 모두 닫기 횟수가 많을수록 커져요. 두 비율은 그렇지 않은데, 여기서는 두 점수가 모두 음수예요. 그리고 두 점수 모두 손실이 아닌 비율에서만 구간이 0보다 아래에 있어요. AWP는 손실이 아닌 비율에서 $-0.254$($[-0.404,-0.104]$)이고, 이익을 낸 비율에서는 $-0.124$($[-0.265,0.015]$)예요. 엔도스랭크는 $-0.178$($[-0.298,-0.042]$) 대 $-0.055$($[-0.196,0.084]$)예요. 송금 입차수도 두 비율 모두에서 음수예요. 값은 $-0.191$($[-0.329,-0.041]$)과 $-0.197$($[-0.350,-0.035]$)이에요. 어느 점수도 이 계약 트레이더(GMX V2에서 거래하는 계약)들을 성공한 닫기의 비율대로 줄 세우지 못해요.

\dissertationSubheading[sec:er-awp-neutral-grid]{두 기간의 모든 트레이더 라벨}

표~\ref{tab:neutral-label-grid}는 두 기간에서 트레이더 라벨 여섯 개 모두에 대해 허락 쪽과 송금 쪽을 맞세워요. 이 표의 쌍 네 개와 라벨 여섯 개는 등록 분석을 위해 W1과 함께 등록했어요. 7월부터 9월까지의 로그를 뽑기 전이었어요(\ref{sub:neutral-label-protocol}절). 그때까지 등록 분석은 W0 트레이더를 이익 라벨로는 채점했지만, 청산 라벨로는 채점하지 않았어요. 이 표는 같은 쌍을 고친 점수로 보여 줘요. 고친 점수는 청산 라벨을 포함한 모든 등록 분석의 결과를 안 뒤에 정의했어요. 그래서 이 표의 값들은 재분석이지, 등록한 시험이 아니에요(부록~\ref{sec:map-registered}). 청산 없이 닫은 비율에서의 P와 Q는, W1의 값과 W0의 부호로 규칙 R1과 R2에 들어가요. W1의 두 비율에서의 P는 규칙 R3에 들어가요. 이 규칙들은 모두 고친 점수에 적용해요(주장~C6, \ref{sub:neutral-label-results}절). 나머지 값들은 엔도스랭크 빼기 AWP 열까지 포함해서, 그 옆에 함께 보고해요.

% Generated by experiments/scripts/export_latex.py -- do not edit by hand.
% Source: experiments/data/2-processed-tables-and-evaluations/revised-usd/window-w1/eval_summary.json
\begin{table}[htbp]
\centering
\caption[두 기간의 모든 트레이더 라벨에서 허락 쪽 점수 빼기 송금 쪽 점수.]{허락 쪽 점수에서 송금 쪽 점수를 뺀 $\Delta\tau$와 95\% 짝지은 부트스트랩 구간이에요. 기간 W1(GMX~V2 계약 트레이더 $n=70$)과 W0($n=88$)의 트레이더 라벨 여섯 개 모두에서 구했어요. P는 받은 허락 수 빼기 송금 입차수예요. Q는 활동 재시작을 쓰는 엔도스랭크 빼기 AWP예요(화살표 무게와 다시 출발 규칙이 같음). C-PR은 $\lambda=1$ 빼기 $\lambda=0$이에요(같은 달러 방식으로 만듦). ER은 엔도스랭크 빼기 AWP예요. 이익을 낸 닫기 수와 실현 이익은 닫은 횟수가 많을수록 커져요. 다른 네 라벨은 그렇지 않아요. 신뢰구간(CI)은 계약을 다시 뽑아 구한 95\% 백분위 부트스트랩 구간이에요(짝지은 다시 뽑기 2{,}000번, 시드 42).}
\label{tab:neutral-label-grid}
\fitwidth{%
\begin{tabular}{lcccc}
\toprule
라벨 & P (개수) & Q (페이지랭크) & C-PR 층 & ER 빼기 AWP \\
\midrule
\multicolumn{5}{l}{\emph{기간 W1, 2026년 7월--9월 (트레이더 $n=70$)}} \\
청산 없이 닫은 비율 & \begin{tabular}[c]{@{}c@{}}0.084\\[-1pt]{\footnotesize [-0.013, 0.176]}\end{tabular} & \begin{tabular}[c]{@{}c@{}}-0.041\\[-1pt]{\footnotesize [-0.135, 0.050]}\end{tabular} & \begin{tabular}[c]{@{}c@{}}-0.049\\[-1pt]{\footnotesize [-0.169, 0.060]}\end{tabular} & \begin{tabular}[c]{@{}c@{}}-0.032\\[-1pt]{\footnotesize [-0.133, 0.066]}\end{tabular} \\
청산 없음 & \begin{tabular}[c]{@{}c@{}}0.078\\[-1pt]{\footnotesize [-0.012, 0.166]}\end{tabular} & \begin{tabular}[c]{@{}c@{}}-0.039\\[-1pt]{\footnotesize [-0.138, 0.062]}\end{tabular} & \begin{tabular}[c]{@{}c@{}}-0.051\\[-1pt]{\footnotesize [-0.170, 0.064]}\end{tabular} & \begin{tabular}[c]{@{}c@{}}-0.029\\[-1pt]{\footnotesize [-0.136, 0.080]}\end{tabular} \\
이익을 낸 닫기 수 & \begin{tabular}[c]{@{}c@{}}-0.143\\[-1pt]{\footnotesize [-0.318, 0.013]}\end{tabular} & \begin{tabular}[c]{@{}c@{}}-0.231\\[-1pt]{\footnotesize [-0.429, -0.027]}\end{tabular} & \begin{tabular}[c]{@{}c@{}}-0.190\\[-1pt]{\footnotesize [-0.411, 0.037]}\end{tabular} & \begin{tabular}[c]{@{}c@{}}-0.214\\[-1pt]{\footnotesize [-0.423, -0.012]}\end{tabular} \\
실현 이익 & \begin{tabular}[c]{@{}c@{}}-0.135\\[-1pt]{\footnotesize [-0.308, 0.025]}\end{tabular} & \begin{tabular}[c]{@{}c@{}}-0.350\\[-1pt]{\footnotesize [-0.536, -0.157]}\end{tabular} & \begin{tabular}[c]{@{}c@{}}-0.404\\[-1pt]{\footnotesize [-0.624, -0.182]}\end{tabular} & \begin{tabular}[c]{@{}c@{}}-0.346\\[-1pt]{\footnotesize [-0.547, -0.144]}\end{tabular} \\
이익을 낸 비율 & \begin{tabular}[c]{@{}c@{}}-0.196\\[-1pt]{\footnotesize [-0.364, -0.046]}\end{tabular} & \begin{tabular}[c]{@{}c@{}}-0.336\\[-1pt]{\footnotesize [-0.523, -0.141]}\end{tabular} & \begin{tabular}[c]{@{}c@{}}-0.344\\[-1pt]{\footnotesize [-0.551, -0.120]}\end{tabular} & \begin{tabular}[c]{@{}c@{}}-0.316\\[-1pt]{\footnotesize [-0.517, -0.112]}\end{tabular} \\
손실이 아닌 비율 & \begin{tabular}[c]{@{}c@{}}-0.252\\[-1pt]{\footnotesize [-0.391, -0.120]}\end{tabular} & \begin{tabular}[c]{@{}c@{}}-0.350\\[-1pt]{\footnotesize [-0.526, -0.175]}\end{tabular} & \begin{tabular}[c]{@{}c@{}}-0.324\\[-1pt]{\footnotesize [-0.531, -0.109]}\end{tabular} & \begin{tabular}[c]{@{}c@{}}-0.346\\[-1pt]{\footnotesize [-0.526, -0.161]}\end{tabular} \\
\midrule
\multicolumn{5}{l}{\emph{기간 W0, 2026년 4월--6월 (트레이더 $n=88$)}} \\
청산 없이 닫은 비율 & \begin{tabular}[c]{@{}c@{}}0.182\\[-1pt]{\footnotesize [0.052, 0.314]}\end{tabular} & \begin{tabular}[c]{@{}c@{}}-0.014\\[-1pt]{\footnotesize [-0.147, 0.130]}\end{tabular} & \begin{tabular}[c]{@{}c@{}}-0.041\\[-1pt]{\footnotesize [-0.187, 0.105]}\end{tabular} & \begin{tabular}[c]{@{}c@{}}-0.017\\[-1pt]{\footnotesize [-0.156, 0.128]}\end{tabular} \\
청산 없음 & \begin{tabular}[c]{@{}c@{}}0.197\\[-1pt]{\footnotesize [0.059, 0.333]}\end{tabular} & \begin{tabular}[c]{@{}c@{}}-0.010\\[-1pt]{\footnotesize [-0.152, 0.146]}\end{tabular} & \begin{tabular}[c]{@{}c@{}}-0.051\\[-1pt]{\footnotesize [-0.198, 0.102]}\end{tabular} & \begin{tabular}[c]{@{}c@{}}-0.012\\[-1pt]{\footnotesize [-0.154, 0.137]}\end{tabular} \\
이익을 낸 닫기 수 & \begin{tabular}[c]{@{}c@{}}0.075\\[-1pt]{\footnotesize [-0.080, 0.222]}\end{tabular} & \begin{tabular}[c]{@{}c@{}}-0.016\\[-1pt]{\footnotesize [-0.150, 0.118]}\end{tabular} & \begin{tabular}[c]{@{}c@{}}-0.067\\[-1pt]{\footnotesize [-0.222, 0.087]}\end{tabular} & \begin{tabular}[c]{@{}c@{}}-0.005\\[-1pt]{\footnotesize [-0.163, 0.155]}\end{tabular} \\
실현 이익 & \begin{tabular}[c]{@{}c@{}}0.017\\[-1pt]{\footnotesize [-0.150, 0.176]}\end{tabular} & \begin{tabular}[c]{@{}c@{}}-0.107\\[-1pt]{\footnotesize [-0.252, 0.044]}\end{tabular} & \begin{tabular}[c]{@{}c@{}}-0.248\\[-1pt]{\footnotesize [-0.433, -0.067]}\end{tabular} & \begin{tabular}[c]{@{}c@{}}-0.118\\[-1pt]{\footnotesize [-0.286, 0.046]}\end{tabular} \\
이익을 낸 비율 & \begin{tabular}[c]{@{}c@{}}0.045\\[-1pt]{\footnotesize [-0.113, 0.184]}\end{tabular} & \begin{tabular}[c]{@{}c@{}}0.046\\[-1pt]{\footnotesize [-0.071, 0.172]}\end{tabular} & \begin{tabular}[c]{@{}c@{}}0.027\\[-1pt]{\footnotesize [-0.112, 0.181]}\end{tabular} & \begin{tabular}[c]{@{}c@{}}0.069\\[-1pt]{\footnotesize [-0.060, 0.208]}\end{tabular} \\
손실이 아닌 비율 & \begin{tabular}[c]{@{}c@{}}-0.080\\[-1pt]{\footnotesize [-0.228, 0.058]}\end{tabular} & \begin{tabular}[c]{@{}c@{}}0.055\\[-1pt]{\footnotesize [-0.070, 0.192]}\end{tabular} & \begin{tabular}[c]{@{}c@{}}0.037\\[-1pt]{\footnotesize [-0.111, 0.189]}\end{tabular} & \begin{tabular}[c]{@{}c@{}}0.076\\[-1pt]{\footnotesize [-0.063, 0.223]}\end{tabular} \\
\bottomrule
\end{tabular}
}
\end{table}


W1에서는 어느 쌍도 두 청산 라벨 어디에서든 두 쪽을 갈라놓지 못해요. 모든 구간이 0을 품어요. 청산 없이 닫은 비율에서 P는 $+0.084$($[-0.013,0.176]$)이고, 엔도스랭크 쌍은 $-0.032$($[-0.133,0.066]$)예요. W0에서는 두 청산 라벨 모두에서 그냥 센 수가 허락 쪽에 유리해요. P는 청산 없이 닫은 비율에서 $+0.182$($[0.052,0.314]$), 청산이 없었는지 나타내는 표시에서 $+0.197$($[0.059,0.333]$)이에요. 반면에 페이지랭크 쌍 세 개는 0보다 아래에 있고, 그 구간은 0을 품어요. W1의 이익 라벨 네 개에서는 모든 쌍이 음수예요. Q와 엔도스랭크 쌍의 구간은 네 라벨 모두에서 0보다 아래에 있어요. P의 구간은 두 비율에서, C-PR(결합 페이지랭크, 두 그래프를 함께 걷는 점수) 층들의 구간은 실현 이익과 두 비율에서 0보다 아래에 있어요. W0에서는 이익 라벨에서의 모든 구간이 0을 품어요. 하나만 예외예요. 실현 이익에서의 C-PR 층들의 구간이고, 값은 $-0.248$($[-0.433,-0.067]$)이에요. 그래서 실현 이익에서의 C-PR 층들의 구간은 두 기간 모두에서 0보다 아래에 있어요. 실현 이익은 닫기 횟수가 많을수록 커지는 라벨이에요. W0의 두 비율에서는 Q, C-PR 층들, 엔도스랭크 쌍이 허락 쪽으로 기울지만, 그 구간은 0을 품어요. 그래서 두 비율은 두 기간 모두에서 유지되는 방향을 보여 주지 않아요. 트레이더가 $70$개와 $88$개라서, 이 차이들 가운데 어느 것도 정밀하지 않아요.

% !TEX root = ../main.tex
\section{주장과 증거를 잇는 지도}
\label{ch:appendix-claims}

이 부록은 각 가설(H), 연구 질문(RQ), 그리고 실제 자료로 확인한 주장(C)을 \ref{ch:results}장에 있는 관련 증거와 이어 줘요. 그리고 그 증거가 보여 주지 않는 것도 적어요. 여기서 말하는 증거는 수정 분석(등록 뒤에 금액을 달러로 바꿔 다시 한 분석)의 증거예요. 수정 분석에서는 다섯 토큰의 금액을 달러(USD)로 쟀어요. 이 방법은 등록 분석(미리 등록한 그대로 한 분석)의 결과를 안 뒤에 정했어요(\ref{sec:amount-revision}절). 등록 분석의 결과는 부록~\ref{sec:map-registered}에 남겨 두었어요.

\dissertationSubheading[sec:map-h1]{H1 / RQ1: 순위 갈라짐}

RQ1은 엔도스랭크(EndorseRank)와 AWP(적응형 가중 페이지랭크)가 같은 스마트 계약(블록체인 위에서 저절로 실행되는 프로그램)들을 서로 다른 순서로 세우는지 물어요.

\begin{itemize}
\item 진술: 매칭 코호트(모든 점수가 똑같이 점수를 매기는 계약 묶음)에서 엔도스랭크와 AWP는 계약들에 서로 다른 순위를 매겨요.
\item 증거: \ref{sec:rank-divergence}절, 표~\ref{tab:tie-sensitivity}, 그리고 그림~\ref{fig:rank-scatter}과~\ref{fig:score-distributions}에 있어요. 그 표의 마지막 행은 두 방법 사이의 켄달의 타우 $\tau_b$(두 줄 세우기가 얼마나 같은 순서인지 나타내는 수)와 그 구간(믿을 만한 범위)을 보여 줘요. 각 그래프(점과 화살표로 이은 그림) 밖에 있는 계약을 외톨이 점(화살표가 하나도 없는 점)으로 남겨 둔 경우의 값도 함께 보여 줘요.
\item 해석: 두 순위는 같은 것이 아니에요. 두 점수가 모두 순서를 매기는 쌍 가운데 많은 몫이 서로 반대 순서로 놓여요. 그리고 그래프 밖의 계약을 다루는 규칙은 두 점수의 일치를 바꾸지 않아요. 목표 O1의 기준은 채워졌어요(표~\ref{tab:objectives-achievement}). 등록 분석에서도 그랬어요(표~\ref{tab:registered-vs-revised}).
\end{itemize}

\dissertationSubheading[sec:map-h2]{H2 / RQ2: 같은 기간 일치도}

RQ2는 각 점수가 같은 기간의 허락(allowance, 토큰을 얼마까지 대신 써도 된다는 허락) 점검과 얼마나 가깝게 맞는지 물어요. 그리고 같은 기간의 송금(transfer) 점검과 얼마나 가깝게 맞는지도 물어요.

\begin{itemize}
\item 진술: 같은 기간 안에서, 각 점수는 자기가 걷는 화살표(엣지) 종류의 대리 지표(대신 재는 값)와 가장 잘 맞아요. 엔도스랭크는 허락 지표 가족(비슷한 지표 묶음)과, AWP는 송금 지표 가족과 가장 잘 맞아요. 그리고 C-PR(결합 페이지랭크, 두 그래프를 함께 걷는 점수)의 혼합 점수들은 모든 가족에서 자기의 두 층 사이에 있어요. 금액을 달러로 재면, AWP는 계약에 들어오는 송금 금액의 합을 따라가요. 그 가까운 정도는, 그 송금을 보낸 쪽(송신자)의 수, 곧 들어오는 화살표 수(입차수)를 따라가는 정도와 거의 같아요.
\item 증거: 모든 점수의 지표 가족 평균과 그 구간은 표~\ref{tab:alignment-hybrid}에 있어요. 송금 대리 지표 하나하나의 값은 표~\ref{tab:alignment-transfer}에 있어요. C-PR과 S-PR(보증 씨앗 페이지랭크)의 짝지은 대비는 표~\ref{tab:tau-diff-hybrid}에 있어요. 그래프 밖의 계약을 외톨이 점으로 남겨 둔 같은 기간 계수는 표~\ref{tab:tie-sensitivity}에 있어요.
\item 참고: 이것은 의도한 구성개념(점수가 재려고 하는 생각) 점검이에요. 대리 지표가 그래프와 같은 이벤트(계약이 남기는 기록)에서 나오기 때문에, 일부는 만든 방식 때문에 저절로 그렇게 돼요. 엔도스랭크는 코호트 안에 동점을 조금만 남겨요. 그리고 그래프 밖의 계약을 다루는 규칙은 어떤 계수도 바꾸지 않아요. 계획에서 정한 것은 규칙~A의 같은 기간 부분뿐이에요. 이 부분은 C-PR에 관한 것이에요(\ref{sec:coupled-results}절). 그런데 여기서는 이 부분을 등록 분석의 결과를 안 뒤에 정의한 점수에 적용해요. 그래서 이 비교들 가운데 어느 것도 자기 자료보다 먼저 정해지지 않았어요. 엔도스랭크와 AWP 사이의 짝지은 차이와 대리 지표별 결과는 부록~\ref{ch:appendix-er-awp}에 있어요.
\end{itemize}

\dissertationSubheading[sec:map-h3]{H3 / RQ3: 계산 비용}

RQ3은 같은 계약 집합에서 각 점수를 계산하는 데 비용이 얼마나 드는지 물어요.

\begin{itemize}
\item 진술: 같은 계약들에서, 풀이 프로그램(solver)을 한 번 부를 때 걸리는 시간은 그 그래프의 화살표 수에 반복(한 바퀴) 횟수를 곱한 값을 따라가요. 그래서 더 듬성듬성한(화살표가 적은) 그래프를 쓰는 엔도스랭크가 시간이 가장 적게 걸려요. 반면 C-PR은 두 층을 합친 더 큰 그래프에서 돌아요. 그리고 S-PR은 다시 출발(재시작)을 위해 허락 층을 한 번 더 풀어요. 그래서 C-PR과 S-PR은 AWP보다 오래 걸려요.
\item 증거: 표~\ref{tab:benchmark-runtime}(실행 시간, 표준편차(SD), 메모리, 반복 횟수, $|V|$, $|E|$), 표~\ref{tab:benchmark-scaling}, 그리고 표~\ref{tab:benchmark-runs}의 모든 시간 잰 실행에 있어요.
\item 참고: 모든 점수는 같은 풀이 프로그램을 써요. 그래서 속도는 $|E|$의 성질이지, 어떤 알고리즘이 가진 장점이 아니에요(부록~\ref{sec:pr-sparsity}). 규모에 따른 변화 단계는 다섯 토큰의 코호트 그래프에서 뽑은 부분 표본이에요. 그래서 이 주장은 그 그래프보다 크지 않은 그래프에만 들어맞아요.
\end{itemize}

\dissertationSubheading[sec:map-rq4]{RQ4: 시간 홀드아웃}

RQ4는 $t_1$에 동결한(고정한) 점수가 스펜더 계약(허락을 받아 대신 쓰는 계약)에 나중에 새로 들어오는 것들을 잘 줄 세우는지 물어요. 그리고 어떤 페이지랭크(PageRank)가 자기가 매끄럽게 다듬는(평활하는) 원시 차수(화살표를 그냥 센 수)보다 이것들을 더 잘 줄 세우는지도 물어요.

\begin{itemize}
\item 진술: 각 라벨(나중에 채점에 쓰는 정답 값) 기간이 시작할 때 점수를 동결했어요(W0은 $t_1$, W1은 $T_{\mathrm{obs}}$). 그러면 각 점수는 스펜더 코호트에서 자기 화살표 종류로 나중에 새로 들어온 것들을 가장 잘 줄 세워요. 동결일의 원시 차수와 견주면 두 걷기는 서로 달라요. AWP는 두 기간 모두에서 나중의 새 송금자(새로 송금을 보내온 주소)를 송금 입차수(송금을 보내온 서로 다른 주소 수)가 줄 세우는 것보다 더 잘 줄 세워요. 반면 계획에서 정한 고르게 다시 출발하기(균등 재시작)를 그대로 쓰는 엔도스랭크는, 두 기간 모두에서 나중의 새 허락 쌍(새로 생긴 주인–스펜더 허락)을 받은 허락 수(허락 입차수)가 줄 세우는 것보다 못하게 줄 세워요. 등록 분석도 이 대비들에 같은 부호를 주었어요(표~\ref{tab:registered-vs-revised}).
\item 증거: 표~\ref{tab:holdout-spenders}과~\ref{tab:holdout-diff}(\ref{sec:holdout-results}절)에 있어요. 기간 W1은 표~\ref{tab:fresh-posthoc}에 있어요. 동결 전에 송금을 받은 적이 있는 스펜더는 표~\ref{tab:holdout-receiving}에 있어요. 활동 재시작(활발한 주소에서 더 자주 다시 출발하기)을 쓰는 엔도스랭크는 표~\ref{tab:endorserank-restarts}에 있어요.
\item 주의할 점: 두 라벨은 나중 기간에 관찰한 구성개념이에요. 곧 다섯 토큰에서 새 허락 쌍의 수와 새 송금자의 수예요. 신용이나 거래의 결과가 아니에요. 등록 분석에서는 규칙~A의 두 대비를 계획에서 정했어요. 그리고 규칙~B의 대비는 F4--F6과 함께 W1 등록 문서에서 정했어요. 둘 다 이 연구가 그 판정에 쓸 자료를 뽑기 전에 정했어요. 하지만 고친 점수는 등록 분석의 결과를 안 뒤에 정의했어요(\ref{sec:amount-revision}절). 그래서 고친 점수에서는 어느 대비도 미리 정한 것이 아니에요. 다음 값들은 라벨을 본 뒤에 계산했어요. 엔도스랭크와 AWP의 차수 대비, C-PR ($\lambda=1$)이 받은 허락 수보다 더 얻은 몫, S-PR의 대비, 그리고 엔도스랭크와 AWP 사이의 차이(부록~\ref{ch:appendix-er-awp})예요. AWP의 활동 재시작을 쓰는 엔도스랭크는 두 기간 모두에서 새 허락 쌍을 받은 허락 수가 줄 세우는 것보다 더 잘 줄 세워요. 하지만 그 비교도 라벨을 본 뒤에 했어요. 그래서 그것은 발견이 아니라, 등록해서 시험할 가설이에요.
\end{itemize}

\dissertationSubheading[sec:map-rq5]{RQ5: 결합 연산자}

RQ5는 C-PR이 송금만 걷는 걷기보다 나은지 물어요. 그리고 C-PR이 어느 한 층만 걷는 걷기보다 계약들을 더 잘 줄 세우는지도 물어요.

\begin{itemize}
\item 진술: 고친 점수에서, 두 화살표 종류를 한 번에 걷는 걷기는 한 층만 걷는 것을 이기지 못해요. 규칙~A는 실패해요. C-PR은 새 허락 쌍에서는 허락 걷기보다 나쁘지 않아요. 하지만 새 송금자에서는 송금 걷기보다 나빠요. 그리고 규칙~A의 같은 기간 부분도 실패해요. 규칙~B의 기준도 채워지지 않아요. C-PR은 새 허락 쌍에서 송금 걷기보다 나아요(F1). 그래서 W0에서 본 이득이 W1에서도 다시 나타나요. 하지만 새 송금자에서는 허용 폭을 채우지 못해요(F2). 그리고 트레이더(거래하는 사람)의 청산 없이 닫은 비율(청산, 곧 손실이 너무 커져 거래가 강제로 끝나는 일 없이 거래를 닫은 비율)에서도 아슬아슬하게 채우지 못해요(F3). 여기서는 두 점 추정값이 같지만, 구간의 아래 끝이 $-\delta$보다 아래에 있어요. 첫 번째 민감도 분석(조건을 바꿔 결과가 얼마나 움직이는지 보기)은 AWP를 비교 대상으로 써요. 여기서는 F2가 분명하게 실패해요. 라벨을 본 뒤에 계산한 W0 대비에서, S-PR은 새 송금자에서는 송금 걷기와 같은 수준에 머물러요. 그리고 새 허락 쌍에서는 허락 걷기보다 아래예요. 곧, 허락 화살표는 걷기가 그 화살표를 따라 움직일 때에만 도움이 돼요. 등록 분석도 같은 판정에 이르렀어요(부록~\ref{sec:map-registered}).
\item 증거, 규칙~A: 주 부분은 표~\ref{tab:holdout-diff}에 있어요. 보조 부분은 표~\ref{tab:alignment-hybrid}과~\ref{tab:tau-diff-hybrid}에 있어요(\ref{sec:coupled-results}절).
\item 증거, 규칙~B: 표~\ref{tab:fresh-contrasts}에 있어요. 이 표에는 F1--F3이 있고, 그 옆에 F4--F6과 두 민감도 분석도 함께 보고해요(\ref{sub:fresh-results}절). W0의 탐색적(미리 정하지 않고 살펴본) 대비는 표~\ref{tab:holdout-diff}에 있어요. 등록 분석의 판정과 수정 분석의 판정을 나란히 놓은 것은 표~\ref{tab:registered-vs-revised}에 있어요.
\item 증거, 비용: 표~\ref{tab:benchmark-runtime}의 C-PR 행과 S-PR 행에 있어요.
\item 주의할 점: 이 시험은 다음 조건에서 했어요. $\lambda\in\{0.25,0.5,0.75\}$로 고정했어요. 층은 다섯 토큰의 금액을 달러로 쟀고, 허락 하나하나는 상한 $U$(이보다 큰 금액은 이 값으로 자름)에서 잘랐어요. 순간 이동(다시 출발)은 고르게 했어요. 섞기는 이동 행렬(전이 행렬) 수준에서 했어요. 그리고 체인은 하나만 썼어요. 규칙~A와 $\lambda$ 값들은 계획과 설정 파일에서 정했어요. 계획과 설정 파일은 이 연구의 로그를 하나도 뽑기 전인 2026년 10월 9일에 \texttt{43d8045}로 커밋했어요(\ref{sec:holdout-protocol}절). 규칙~B, 그 허용 폭, 그리고 그 두 민감도 분석은 W0의 결과를 안 뒤에 썼어요. 그리고 스캔~B 전에 \texttt{0d3a057}로 등록했어요(\ref{sub:fresh-holdout}절). 달러로 잰 층과 상한은 등록 분석의 결과를 안 뒤인 2026년 10월 10일에 정했어요(\ref{sec:amount-revision}절). 그래서 위의 판정들은 이 규칙들을 고친 점수에 재분석으로 적용한 것이에요. 등록한 시험이 아니에요.
\end{itemize}

\dissertationSubheading[sec:map-c1c4]{주장 C1--C5}

\ref{ch:results}장은 주장 C1--C6을 말해요. 처음 다섯 주장은 아래의 증거에 기대요. C6의 증거는 부록~\ref{sec:map-c6}에 이어서 나와요.

\begin{itemize}
\item C1: H3/RQ3 아래에 적은 실행 시간 증거예요. 화살표 수의 비와 시간의 비는 \ref{sec:benchmark-results}절과 부록~\ref{sec:pr-sparsity}에 있어요. 속도는 알고리즘이 아니라 더 듬성듬성한 그래프 덕분이에요. 그리고 C-PR과 S-PR은 AWP보다 비용이 더 들어요.
\item C2: 모든 점수가 자기 화살표로 만든 대리 지표와 같은 기간 안에서 보이는 일치예요(표~\ref{tab:alignment-hybrid}). 이것은 의도한 구성개념 점검이고, 일부는 만든 방식 때문에 저절로 그렇게 돼요. AWP는 들어오는 송금의 달러 금액을, 그 송금을 보낸 쪽의 수를 따라가는 것과 거의 같은 정도로 가깝게 따라가요(표~\ref{tab:alignment-transfer}). 동점과, 그래프 밖의 계약을 다루는 규칙은 표~\ref{tab:tie-sensitivity}에 있어요. 감쇠 계수(화살표를 따라갈 확률) 점검과 금액 기울기 점검은 엔도스랭크와 AWP의 같은 기간 지표 가족 평균을 다시 계산해요. 이 점검들은 표~\ref{tab:robustness-damping}과~\ref{tab:robustness-slope}, 그리고 그림~\ref{fig:damping-sensitivity}에 있어요. 표본 정의 단계(표~\ref{tab:robustness-sample-size})는 평가 집합에 따라 통계량이 어떻게 움직이는지 보여 줘요. 이것은 강건성(조건을 바꿔도 결과가 버티는지)의 증거가 아니에요.
\item C3: H1/RQ1 아래에 적은 순위 갈라짐이에요(표~\ref{tab:tie-sensitivity}, 그림~\ref{fig:rank-scatter}).
\item C4: 각 걷기를 자기 화살표 종류의 원시 차수와 기간 밖에서 견준 비교예요(표~\ref{tab:holdout-spenders}, \ref{tab:holdout-diff}과~\ref{tab:fresh-posthoc}). 포함 범위 점검(표~\ref{tab:holdout-receiving})과 활동 재시작을 쓰는 변형(표~\ref{tab:endorserank-restarts})도 함께 있어요.
\item C5: 고친 점수에서, 결합은 두 층을 각각 한 층만 걷는 것을 이기지 못해요. 이것은 범위가 정해진 부정적인 결과예요. 그리고 송금 걷기 하나와 견주는 규칙~B를 이 점수들에 적용하면, 그 규칙의 기준이 채워지지 않아요(표~\ref{tab:fresh-contrasts}). 등록 분석도 같은 판정에 이르렀어요(표~\ref{tab:registered-vs-revised}). 위의 RQ5와 \ref{sec:coupled-results}절을 보세요.
\end{itemize}

\dissertationSubheading[sec:map-nonclaims]{분명히 밝히는, 주장하지 않는 것들}

이 연구는 다음을 주장하지 \emph{않아요}.

\begin{itemize}
\item 어떤 점수가 일반적으로 평판을 더 잘 재는 척도라는 것, 또는 AWP가 이제 쓸모없어졌다는 것(부록~\ref{ch:appendix-er-awp}의 엔도스랭크와 AWP 비교는 구성개념을 따르고, 포함 범위에 달려 있어요).
\item 어떤 점수가 자기 화살표로 만든 대리 지표와 맞는 것이, 그 점수 바깥에서 온 증거라는 것(이것은 의도한 구성개념을 점검하는 것이고, 일부는 만든 방식 때문에 저절로 그렇게 돼요).
\item 계획의 고르게 다시 출발하기를 쓰는 엔도스랭크가 나중의 새 허락 쌍을 그냥 센 받은 허락 수가 줄 세우는 것보다 더 잘 줄 세운다는 것(두 기간 모두에서 더 못하게 줄 세웠어요).
\item 활동 재시작을 쓰는 엔도스랭크가 받은 허락 수를 이긴다는 것(그 비교는 라벨을 본 뒤에 했고, 등록한 시험을 기다리고 있어요).
\item 허락 그래프에는 여러 걸음 퍼뜨림(전파)이 있는데, 이것 때문에 엔도스랭크가 그냥 센 수를 이겼다는 것(\ref{sub:theory-closed-form}절).
\item 청산을 피하는 데에서, 그냥 센 수의 수준(R1)이나 페이지랭크 수준(R2)에서 허락 화살표가 송금 화살표보다 앞선다는 것, 또는 거기서 엔도스랭크가 AWP를 이긴다는 것(고친 점수에서도, 등록 분석에서도 어느 규칙도 채워지지 않았어요).
\item 어떤 점수가 신용도(돈을 갚을 만한 정도), 채무 불이행(빚을 못 갚는 일) 위험, 또는 거래 실력을 잰다는 것(W1에서 허락한 주인 수(허락을 해 주는 서로 다른 주인의 수)는 이익을 낸 비율과 손실이 아닌 비율에서 트레이더를 송금 수보다 못하게 줄 세웠어요. 그리고 자료에는 채무 불이행 라벨이 없어요).
\item 허락을 받는 것이 계약 트레이더(GMX V2에서 거래하는 계약)로 하여금 청산을 피하게 만든다는 것(표본이 작고, 허락을 받은 트레이더는 포지션을 더 많이 닫아요. 그리고 송금 수도 이 라벨을 줄 세워요).
\item 금액 기울기 $b=\ln 3/m$, C-PR 허락 층의 상한 $U$, 또는 엔도스랭크의 고르게 다시 출발하기가 가장 좋은 선택이라는 것(기울기와 상한은 등록 분석의 결과를 안 뒤에 수정 과정에서 정했어요. 그리고 표~\ref{tab:robustness-slope}의 금액 기울기 점검은 기울기는 바꾸지만 상한은 바꾸지 않아요. 고르게 다시 출발하기는 계획에서 정했고, 활동 재시작과의 비교는 라벨을 본 뒤에 했어요).
\item 고친 점수에서 나온 규칙~A와~B, 그리고 R1--R3의 판정이 등록한 시험이라는 것(이 규칙들은 등록 분석을 위해 정했고, 그 결과를 안 뒤에 정의한 점수에 적용했어요. 등록 분석의 판정은 부록~\ref{sec:map-registered}에 있어요).
\item C-PR이 송금만 걷는 걷기보다 낫다는 것(규칙~B는 새 송금자에서도, 트레이더에서도 실패해요), 또는 AWP보다 낫다는 것(두 기간 모두에서 새 송금자에서 AWP에 뒤져요).
\item 두 화살표 종류를 어떻게 합쳐도 한 층 방법들을 이길 수 없다는 것. 이 연구가 보인 것은 여기서 시험한 방법, 곧 $\lambda$를 고정하고 점마다 볼록 혼합(두 값을 정한 비율로 섞기)을 하는 방법이 이기지 못한다는 것뿐이에요.
\item 엔도스랭크의 더 짧은 푸는 시간이 그 알고리즘의 장점이라는 것(이것은 더 적은 화살표 수에서 나와요), 또는 코호트의 에고 네트워크(코호트 계약에 닿는 화살표만 모은 그래프)보다 큰 그래프에서 실행 시간이 어떻다는 것.
\item 엔도스랭크가 시빌 농장(공격자가 만든 가짜 주소 무리)을 견뎌 낸다는 것(평범한 점수나 중앙값보다 높은 점수에 이르는 데는 비용이 적게 들어요. 그리고 다섯 토큰 가운데 하나의 허락에는 잔고가 필요 없어요. 그래서 토큰 목록이 농장의 비용을 올리지 않아요. \ref{sec:sybil-cost-model}절).
\item 결과가, 살펴본 몇 달 동안 아비트럼 원(Arbitrum One)에서 WETH, USDC, USDT0, ARB, WBTC로 서로 다른 주인 셋 이상에게 허락을 받은 계약을 넘어서도 들어맞는다는 것. 또는 다른 토큰의 허락과 송금에도 들어맞는다는 것(수정 분석은 이것들을 빼 두었어요. 표~\ref{tab:usd-tokens}).
\item 온체인(블록체인 위의) 점수가 담보, 감사(코드 점검), 신원 시스템을 대신한다는 것. 또는 이 점수가 오프체인(블록체인 밖의) 정보를 쓴다는 것(이 연구는 그런 정보를 쓰지 않아요).
\end{itemize}

\dissertationSubheading[sec:map-artifacts]{재현 자료}

아래의 경로는 부록~\ref{ch:appendix-eval}에서 이름을 밝힌 실험 폴더를 기준으로 한 상대 경로예요. 본문의 장들이 보고하는 수정 분석은 자기 폴더가 따로 있어요. 등록 분석의 폴더는 얻었을 때 그대로 남겨 두었어요.

\begin{enumerate}
\item 모든 방법의 계약별 점수예요. 같은 기간의 코호트 점수, 그리고 W0과 W1의 스펜더 점수와 트레이더 점수가 parquet 파일에 들어 있어요. 수정 분석의 것은 \texttt{data/2-processed-tables-and-evaluations/revised-usd/}(\texttt{same-window/}, \texttt{holdout-w0/}, \texttt{window-w1/})에 있어요. 등록 분석의 것은 \texttt{data/2-processed-tables-and-evaluations/}의 \texttt{same-window/}, \texttt{holdout-w0/}, \texttt{registered-w1/}에 있어요.
\item 기계가 읽을 수 있는 요약이에요. 같은 기간(감쇠 계수를 바꿔 본 실험 포함), W0, W1, 강건성 점검, 보충 대비, 시빌 모형, 벤치마크(속도 재기), 자료 설명의 요약이에요. 여기에는 모든 계수, 구간, 대비, 판정이 들어 있어요. 수정 분석의 요약은 금액 기울기 점검과 함께 \texttt{data/3-published-results-for-thesis/revised-usd/}에 있어요. 등록 분석의 요약은 상위 토큰 점검과 함께 \texttt{data/3-published-results-for-thesis/}에 있어요. 폴더마다 자기 \texttt{results\_digest.md}와 \texttt{SHA256SUMS.txt}가 있어요.
\item 분석 계획(\texttt{docs/analysis\_plan.md})과 설정 파일이에요. 둘 다 이 연구의 로그를 하나도 뽑기 전인 2026년 10월 9일에 \texttt{43d8045}로 커밋한 그대로예요.
\item W1 등록 문서(\texttt{docs/registration\_w1.md})예요. \texttt{0d3a057}로 커밋한 그대로예요. 그리고 추출 목록 \texttt{data/2-processed-tables-and-evaluations/registered-w1/extraction\_manifest.json}도 있어요. 이 목록에는 커밋, W1 추출을 시작한 시각과 끝낸 시각, 강제 실행 옵션(override)을 쓰지 않았다는 것, 그리고 SHA-256 해시가 적혀 있어요. 등록 문서의 해시는 \texttt{2126c8dc3fe3c0f9273aa4ecfc1e15bddc1f5a6f65e98548b9528f2b18ca2a6b}이고, 설정 파일의 해시는 \texttt{d6a30adbcabcf511d6892332da73639efdbece25891b392e13e90faf09b7a833}이에요.
\item 금액 수정 문서(\texttt{docs/revision\_price\_weighting.md})예요. 2026년 10월 10일 07:54~UTC에 \texttt{a897475}로 커밋했어요. 다섯 토큰의 선택은 08:12~UTC에 \texttt{09dfc32}로, 고친 허락 상한과 상수들은 11:24~UTC에 \texttt{2705b28}로 커밋했어요. 모두 등록 분석의 결과를 안 뒤이고, 고친 점수를 하나라도 계산하기 전이에요.
\item \texttt{data/0-usd-token-list-and-prices/}에 있는 토큰 목록과 가격이에요. 여기에는 체인링크(Chainlink)의 가격 피드(블록체인에 올라오는 가격 정보) 목록을 찍어 둔 것과 검토한 후보 토큰들, 조건을 채운 토큰들, 하루 달러 가격과 그 가격이 나온 DefiLlama 데이터 점들, 상수 $m_T$, $m_A$, $U$, 그리고 이 가격을 체인에서 읽은 Chainlink 피드와 견준 결과가 있어요.
\item 비용 장부 \texttt{data/bq\_ledger.json}이에요. 여기에는 모든 빅쿼리(BigQuery) 작업과 그 매개변수, 그리고 요금이 매겨진 바이트 수가 있어요.
\item \texttt{data/1-raw-blockchain-extracts/}에 있는 뽑은 표들이에요. 계정 종류가 붙은 코호트 후보, GMX~V2 계정과 닫기 기록이 있어요. \texttt{graph-tables/}에는 등록 분석의 쌍 표, 점 표, 대리 지표, 라벨, 상위 토큰 부분 그래프가 있어요. \texttt{graph-tables-usd/}에는 다섯 토큰에 대한 수정 코호트와 그 쌍 표, 점 표, 대리 지표, 라벨이 있어요. 모두 90\,MB보다 작은 parquet 조각으로 나뉘어 있고, 조각마다 원본 표를 적은 JSON 기록이 붙어 있어요.
\item 이 논문 폴더에 있는 내보낸 표 조각(\texttt{results/tables/})과 벡터 그림(\texttt{Figures/generated/})이에요. 이것들은 수정 분석의 요약으로 만들었어요. 표~\ref{tab:registered-vs-revised}은 등록 분석의 요약도 읽어요.
\end{enumerate}

표와 그림은 뽑은 표 위에서 평가 스크립트를 돌린 다음, 같은 요약에서 표와 그림을 내보내면 다시 만들어져요(부록~\ref{sec:pipeline-stages}). \texttt{CONTRACT\_REP\_VARIANT=usd}로 두면 스크립트는 다섯 토큰의 표를 읽어요. 이 표는 \texttt{scripts/run\_usd\_pipeline.py}가 빅쿼리에서 만들어요. 그리고 스크립트는 수정한 결과를 써요. 이 변수를 설정하지 않으면 스크립트는 등록 분석의 결과를 다시 만들어요.

\dissertationSubheading[sec:map-c6]{C6: 어느 화살표 종류로도 만들지 않은 결과 값}

C6은 두 그래프 어디에도 들어 있지 않은 트레이더 결과 값을 놓고 두 화살표 종류를 견줘요.

\begin{itemize}
\item 진술: GMX~V2 계약 트레이더의 청산 없이 닫은 비율을 봤어요. 이때 W1과 함께 등록한 규칙을 고친 점수에 적용했어요. 그랬더니 허락한 주인 수가 송금 수보다 그냥 센 수의 수준에서 앞선다는 것(R1)은 W1에서 뒷받침되지 않아요. W1은 규칙이 판정을 내리는 기간이에요. W0에서는 P가 양수인데도 그래요. 페이지랭크 수준에서 앞선다는 것(R2)도 뒷받침되지 않아요. 그리고 R3이 적용돼요. 등록할 때 정한 문구로 말하면 이래요. 나중에 청산을 피하는 데에서, 그냥 센 수의 수준으로는 허락이 송금보다 앞서지 않아요. 페이지랭크 버전들은 어느 쪽이 나은지 가려지지 않았어요. 그리고 이 신호는 강제 청산을 피하는 것에 관한 것이지, 거래 성공에 관한 것이 아니에요. 이 판정들은 재분석이지 등록한 시험이 아니에요. 등록 분석도 같은 판정에 이르렀어요(표~\ref{tab:registered-vs-revised}).
\item 증거: 표~\ref{tab:neutral-label-contrasts}과~\ref{tab:neutral-label-recipients}(\ref{sub:neutral-label-results}절)에 있어요. 두 기간의 트레이더 라벨 여섯 개 모두는 표~\ref{tab:neutral-label-grid}에 있어요. 절차는 \ref{sub:neutral-label-protocol}절에 있어요. 등록 내용은 \texttt{docs/registration\_w1.md}의 3절에 있어요. 이 문서는 스캔~B 전에 \texttt{0d3a057}로 커밋했어요. 이 등록 내용을 고친 점수에 적용한 것은 \texttt{docs/revision\_price\_weighting.md}에 있어요(\ref{sec:amount-revision}절).
\item 주의할 점: 트레이더 코호트는 작아요. W1에서는 계약 $70$개, W0에서는 $88$개예요. 2026년에 계약 계정이 GMX~V2에서 거래를 적게 했기 때문이에요(표~\ref{tab:gmx-quarters}). 다섯 토큰 가운데 하나의 허락을 받은 적이 있는 트레이더는 두 기간 모두에서 전부 청산을 피했어요. 다른 트레이더는 약 네 개 중 세 개꼴로 피했어요. 하지만 허락을 받은 트레이더는 포지션을 더 많이 닫았어요. 그리고 송금 수도 이 라벨을 줄 세워요. 모든 트레이더가 계약이에요. 그래서 계정 종류는 고정되어 있어요. 이 점수를 걸러 내는 데 쓰려면, 먼저 더 큰 표본에서 등록한 시험을 해야 해요(\ref{sec:future-work}절).
\end{itemize}

\dissertationSubheading[sec:map-registered]{등록 분석}

분석 계획과 W1 등록 문서는 금액을 토큰의 원래 단위(소수점 조정을 하지 않은 가장 작은 단위)로 재고, 모든 ERC-20(토큰 공통 규칙) 토큰을 쓰고, 값 변환(금액이 아무리 커도 1을 넘지 않게 바꾸는 식)에 $b=1$을 쓰기로 정했어요. 금액은 그 분석의 결과를 안 뒤에 고쳤어요(\ref{sec:amount-revision}절). 본문의 장들은 수정 분석을 보고해요. 등록 분석은 얻은 그대로 여기에 남겨 두었어요. 표~\ref{tab:registered-vs-revised}은 등록 분석의 판정 대비와 판정을, 고친 점수에서 구한 같은 대비 옆에 나란히 놓아요. 이 표의 마지막 네 행은 두 기간의 차수 대비예요. 이 대비들은 두 분석 모두에서 라벨을 본 뒤에 계산했고, 판정이 없어요. 등록 열에 있는 모든 판정은 이 연구가 그 자료를 뽑기 전에 정한 시험의 결과예요. 수정 열에 있는 판정은 어느 것도 그렇지 않아요. 규칙을 등록 분석의 결과를 안 뒤에 정의한 점수에 적용했기 때문이에요. 모든 계수, 구간, 판정을 담은 등록 분석의 전체 결과는 계산한 그대로 \texttt{data/3-published-results-for-thesis/}(\texttt{results\_digest.md})에 공개했어요. 수정한 결과는 그 아래 폴더 \texttt{revised-usd/}에 있어요.

\IfFileExists{results/tables/registered-vs-revised.tex}{% Generated by experiments/scripts/export_latex.py -- do not edit by hand.
% Source: experiments/data/2-processed-tables-and-evaluations/same-window/eval_summary.json
% Source: experiments/data/2-processed-tables-and-evaluations/holdout-w0/eval_summary.json
% Source: experiments/data/2-processed-tables-and-evaluations/registered-w1/eval_summary.json
% Source: experiments/data/2-processed-tables-and-evaluations/revised-usd/same-window/eval_summary.json
% Source: experiments/data/2-processed-tables-and-evaluations/revised-usd/holdout-w0/eval_summary.json
% Source: experiments/data/2-processed-tables-and-evaluations/revised-usd/window-w1/eval_summary.json
\begin{table}[htbp]
\centering
\caption[등록 분석과 수정 분석의 결정 대비.]{등록 분석의 결정 대비와 판정이에요(금액은 원래 단위, 모든 ERC-20 토큰, $b=1$). 분석 계획과 W1 등록 문서에 따라 얻은 그대로예요. 그 옆에 수정 분석의 점수에서 구한 같은 대비를 놓았어요(금액은 달러(USD), 목록에 오른 토큰만). 수정은 등록 분석의 결과를 본 뒤에 정했어요. 수정 분석 열은 등록한 규칙과 기준값을 수정 분석의 점수에 적용한 것이고, 등록한 검사가 아니에요. 마지막 묶음에는 두 분석의 어떤 규칙도 판정에 쓰지 않는 차수 대비가 있어요. 이 대비들은 두 분석 모두에서 보고해요. $\Delta\tau$와 95\% 짝지은 부트스트랩 구간을 적었어요(다시 뽑기 400번과 2{,}000번, 시드 42). 규칙과 그 기준값은 \ref{sec:holdout-protocol}절과 \ref{sub:fresh-holdout}절에 있어요.}
\label{tab:registered-vs-revised}
\fitwidth{%
\begin{tabular}{lcc}
\toprule
항목 & 등록 분석 & 수정 분석 \\
\midrule
매칭 코호트 (계약 수) & 27{,}844 & 15{,}107 \\
W0 스펜더 코호트 & 33{,}101 & 17{,}050 \\
W1 스펜더 코호트 & 34{,}587 & 18{,}152 \\
W1 트레이더 & 70 & 70 \\
EndorseRank 대 AWP, 같은 기간 ($\tau$) & 0.217 [0.209, 0.225] & 0.174 [0.163, 0.185] \\
\midrule
A, W0: C-PR 빼기 C-PR ($\lambda=1$), 새 허락 쌍 & $+0.008$ [0.003, 0.013] & $+0.029$ [0.022, 0.036] \\
A, W0: C-PR 빼기 C-PR ($\lambda=0$), 새 송금자 & $-0.030$ [-0.034, -0.026] & $-0.034$ [-0.040, -0.029] \\
규칙 A, 주된 부분 & 충족 안 됨 & 충족 안 됨 \\
규칙 A, 보조 부분 & 충족 안 됨 & 충족 안 됨 \\
\midrule
B, F1: C-PR 빼기 C-PR ($\lambda=0$), 새 허락 쌍 & $+0.101$ [0.092, 0.111] & $+0.110$ [0.098, 0.122] \\
B, F2: C-PR 빼기 C-PR ($\lambda=0$), 새 송금자 & $-0.025$ [-0.029, -0.021] & $-0.028$ [-0.033, -0.023] \\
B, F3: C-PR 빼기 C-PR ($\lambda=0$), 청산 없이 닫은 비율 & $-0.044$ [-0.100, -0.002] & $+0.000$ [-0.014, 0.010] \\
규칙 B (F1, F2, F3) & 충족 안 됨 & 충족 안 됨 \\
규칙 B, 민감도 1 (AWP를 비교 대상으로 함) & 충족 안 됨 & 충족 안 됨 \\
규칙 B, 민감도 2 (외톨이 점) & 충족 안 됨 & 충족 안 됨 \\
F6 (규칙 A, 주된 부분, W1에서) & 불성립 & 불성립 \\
\midrule
P, W1: 받은 허락 수 빼기 송금 입차수 & $+0.063$ [-0.026, 0.148] & $+0.084$ [-0.013, 0.176] \\
Q, W1: EndorseRank (활동에 따른 재시작) 빼기 AWP & $+0.051$ [-0.085, 0.213] & $-0.041$ [-0.135, 0.050] \\
R1 (개수 수준의 앞섬) & 충족 안 됨 & 충족 안 됨 \\
R2 (페이지랭크 수준의 앞섬) & 충족 안 됨 & 충족 안 됨 \\
R3 (특이성) & 적용됨 & 적용됨 \\
\midrule
W0: EndorseRank 빼기 받은 허락 수, 새 쌍 & $-0.011$ [-0.016, -0.006] & $-0.019$ [-0.026, -0.011] \\
W0: AWP 빼기 송금 입차수, 새 송금자 & $+0.038$ [0.035, 0.042] & $+0.043$ [0.038, 0.048] \\
W1: EndorseRank 빼기 받은 허락 수, 새 쌍 & $-0.029$ [-0.032, -0.025] & $-0.034$ [-0.038, -0.029] \\
W1: AWP 빼기 송금 입차수, 새 송금자 & $+0.020$ [0.016, 0.023] & $+0.022$ [0.017, 0.027] \\
\bottomrule
\end{tabular}
}
\end{table}
}{}

% !TEX root = ../main.tex
\section{윤리 승인 증명서}
\label{ch:appendix-ethics}

\dissertationSubheading[sec:ethics-certificate]{증명서}

이 연구는 누구나 볼 수 있는 공개 블록체인 자료만 써요. 그리고 사람을 연구 참여자로 두지 않아요(\ref{sec:ethics}절). 아래에는 유니사(Unisa) 컴퓨팅 학부 / 농업·환경 과학 대학의 연구 윤리 심사 위원회가 발급한 윤리 승인(연구가 윤리 기준을 지킨다는 허가) 증명서를 옮겨 실어요.

% TODO (user action): replace the placeholder box with the certificate PDF once
% the reference number, date and scanned document are available, e.g.
%   \includegraphics[width=\textwidth,page=1]{Figures/ethics-clearance.pdf}
% Do not invent a reference number or date.
\begin{center}
\fbox{\parbox{0.9\textwidth}{\centering\vspace{2em}
윤리 승인 증명서\\[0.5em]
참조 번호: [Unisa SoC/CAES 윤리 승인 참조 번호: TBD]\\
발급일: [TBD]\\[0.5em]
(스캔한 증명서를 여기에 넣을 예정이에요.)
\vspace{2em}}}
\end{center}

